% Modernised reading — fonds Alexandre Grothendieck, shelfmark 1, pages 1-149.
%
% This is an INTERPRETATION, not a transcription. It re-reads the pages in
% today's notation and vocabulary, states the hypotheses the manuscript leaves
% implicit, and drops the critical apparatus so that the mathematics reads
% straight through. Everything the brackets and underlines carried — a doubtful
% reading, a notation he used differently, a missing passage — moves into a
% footnote.
%
% It is held to being mathematically correct as it stands, not merely faithful
% to the page. Where the manuscript is loose or mistaken, the true statement is
% what appears here, and the footnote says what the page has. In any dispute of
% substance the transcription governs: whoever cites Grothendieck must cite the
% batch-NN.fr.tex files.
%
% Scope: the folder was read WHOLE before any of this was written — all eight
% batches, pages 1-149, in one pass — and it is written as ONE file for the
% whole shelfmark, as 19.modern.tex and 29.modern.tex are for theirs.
%
% Notation fixed for the whole folder, in the "Conventions" section: s_n(u)
% for the singular values (his ρ_n(u), r_ν(A)); S_p(E,F) for the Schatten
% classes (his L^p, L^{(1)}, L^{(2)}, 𝓛^p(H)); "opérateur à trace" everywhere
% he writes "opérateur de Fredholm", which in 1953 meant nuclear; M, τ,
% 𝔞, 𝔟 for his 𝒜, T, a, b; his φ_f, Φ_f, Ψ_f, Δ_f kept.
%
% Substantive departures from the pages, each footnoted where it occurs:
%   - page 7: the two formulas for φ_f(t∓0) have their ≤ and < swapped; the
%     page 26 version is the right way round and is what is given;
%   - page 28: "Δ(A) < ∞ ⟺ B ∈ a" holds only as ⇐ (counterexample given);
%   - page 30: Prop. 9 is stated with τ(E) ≤ t and one projector; it needs
%     τ(E) = t and, for non-normal A, a pair of equivalent projectors — the
%     proof of Théorème 1 (page 31) is rewritten with that form;
%   - page 90: Th. 3 needs f increasing (hypothesis a), not b) as its proof
%     says; b) holds only in finite dimension;
%   - page 147: the last step of the lemma integrates n(r)/r down to 0,
%     which diverges; the argument is repaired by a split at δ;
%   - page 148: P_n(u) carries the sign (−1)^k the page omits, and 2) is not
%     proved by "decreasing terms" but by the lemma; 1) is restricted to
%     normal u, the only case the lemma covers;
%   - page 96: an inclusion between ℓ_{N'} and ℓ_{N''} is reversed.
%
% Announced and never established: the plan's « N° 7, autres inégalités »
% (page 2); the marginal « ② Dual de 𝓛^p ; réflexivité » of page 126 (done for
% general Schatten norms at pages 98-104, not in that theorem); the
% inventory's « lettre (1953) », which no transcribed page identifies.
%
% Chronology, argued in the spine section from the pages themselves: pages
% 123-149 predate pages 37-104 (page 129 knows only an analytic proof of
% Σ|λ_i|^p ≤ Σ s_i^p, page 88 has it from Weyl's lemma); pages 106-118 serve
% the cancelled proof of page 79.
%
% Pass: Fable 5.1 (claude-fable-5-1), 2026-09-10 — first pass, unchecked
% against the pages by a human.
\documentclass[11pt,a4paper]{article}
% Le préambule commun à toutes les transcriptions du fonds Grothendieck.
%
% Il tient en une idée : **l'incertitude se compose, elle ne se résout pas.**
% Une transcription qui lisse un mot douteux a détruit la seule chose pour
% laquelle on la faisait — un lecteur doit pouvoir savoir, à chaque endroit,
% ce qui était sur le feuillet et ce qui a été ajouté. Les six macros qui
% suivent sont donc l'appareil critique, et non de la décoration.
%
% Les mêmes commandes sont reconnues par `scripts/render.mjs`, qui produit la
% vue de lecture du site. Une macro ajoutée ici doit l'être aussi là-bas,
% faute de quoi le rendu échouera bruyamment — ce qui est voulu : un rendu
% silencieusement faux est pire qu'un rendu absent.

\ProvidesPackage{grothendieck}[2026/08/08 Transcriptions du fonds Grothendieck]

\usepackage{amsmath,amssymb,amsthm}
\usepackage{mathrsfs}
\usepackage{stmaryrd}
% Les diagrammes commutatifs sont partout dans ce fonds : c'est la langue
% même de la géométrie algébrique de Grothendieck, et une transcription qui
% les rendrait en prose ne transcrirait rien.
\usepackage{tikz-cd}
% Français par défaut — c'est la langue des feuillets — mais l'anglais est
% chargé aussi : la traduction et le résumé sont des documents anglais, et la
% typographie française (espaces avant les deux-points) y serait une faute.
% Ces documents basculent par \selectlanguage{english} en tête de corps.
\usepackage[english,french]{babel}
\usepackage{fontspec}
\usepackage{geometry}
\geometry{a4paper,margin=25mm,marginparwidth=28mm}
\usepackage{marginnote}
\usepackage{xcolor}
% ulem, not soul. `soul`'s \st and \ul are built on a character-by-character
% scan that XeLaTeX's font machinery breaks: the document fails with an
% undefined control sequence rather than degrading. ulem does the same two jobs
% and is Unicode-engine safe. `normalem` stops it hijacking \emph, which the
% transcriptions use for Grothendieck's own emphasis.
\usepackage[normalem]{ulem}
\usepackage{hyperref}
\hypersetup{colorlinks=true,linkcolor=black,urlcolor=black,pdfborder={0 0 0}}

% --- Métadonnées du lot -----------------------------------------------------
% Reprises telles quelles par le script de rendu, qui les affiche en tête de
% la vue de lecture. Elles fixent ce que le fichier prétend couvrir : sans
% elles, une transcription flotte sans point d'attache dans un fonds de seize
% mille feuillets.

\newcommand{\folder}[1]{\gdef\grothFolder{#1}}
\newcommand{\batch}[1]{\gdef\grothBatch{#1}}
\newcommand{\leaves}[2]{\gdef\grothFirst{#1}\gdef\grothLast{#2}}
\newcommand{\foldertitle}[1]{\gdef\grothFoldertitle{#1}}
\gdef\grothFolder{?}\gdef\grothBatch{?}\gdef\grothFirst{?}\gdef\grothLast{?}\gdef\grothFoldertitle{}

% --- Le résumé --------------------------------------------------------------
%
% Il ouvre la lecture modernisée, sous le titre « Résumé », et il est composé
% en petit. Ce n'est pas un ornement : c'est le seul passage du document qui
% ne soit pas des mathématiques, et un lecteur doit pouvoir le reconnaître
% avant de l'avoir lu — puis le sauter s'il connaît déjà le sujet, ou s'y
% arrêter s'il ne le connaît pas. Le filet qui le suit marque la reprise du
% corps.

\newenvironment{resume}
  {\small\setlength{\parskip}{0.45em}}
  {\par\medskip\hrule\medskip}

% --- Mots-clés ---------------------------------------------------------------
%
% En anglais, en clôture du résumé de la lecture modernisée : le vocabulaire
% moderne sous lequel chercher ce que les feuillets construisent. Le manifeste
% du site (`npm run manifest`) les extrait pour étiqueter la cote entière —
% cette ligne est la source unique des tags, il n'y a pas de fichier à côté.
\newcommand{\keywords}[1]{\par\smallskip\noindent
  {\footnotesize\textsc{Keywords} --- \textit{#1}}\par}

% --- L'appareil critique ----------------------------------------------------

% Le numéro de feuillet, en marge. C'est l'ancre qui permet de régler un
% désaccord sur une formule en regardant *un* feuillet plutôt qu'une chemise
% de six cents. Le site s'en sert aussi pour tourner les pages du fac-similé
% quand on fait défiler la transcription.
\newcommand{\leaf}[1]{%
  \marginnote{\footnotesize\color{gray!70}\textsf{#1}}%
  \label{leaf:#1}\ignorespaces}

% Illisible : on ne devine pas. Un mot inventé qui se lit comme les autres est
% le pire résultat possible.
\newcommand{\ill}{\textcolor{red!60!black}{[\,\ldots\,]}}

% Lecture douteuse : le mot est donné, et signalé comme incertain.
\newcommand{\uncertain}[1]{\uline{#1}}

% Ajout de l'éditeur — une lettre manquante, un mot sous-entendu.
\newcommand{\add}[1]{\textcolor{blue!50!black}{[#1]}}

% Biffé par Grothendieck lui-même. Ce qu'il a rayé fait partie du document :
% une rature dit souvent où la pensée a changé de direction.
\newcommand{\struck}[1]{\sout{#1}}

% Note du transcripteur, sur l'état matériel du feuillet.
\newcommand{\note}[1]{\par\smallskip{\small\color{gray!60!black}\textit{[#1]}}\par\smallskip}

% Note marginale de Grothendieck, distincte du corps du texte.
\newcommand{\marginal}[1]{\par\smallskip\noindent\hspace{1em}%
  {\small\color{gray!60!black}#1}\par\smallskip}

% --- Titre ------------------------------------------------------------------

\AtBeginDocument{%
  \begin{center}
    {\large\bfseries Fonds Alexandre Grothendieck --- cote n\textsuperscript{o}~\grothFolder}\\[2pt]
    {\itshape\grothFoldertitle}\\[4pt]
    {\small feuillets \grothFirst--\grothLast\ (lot \grothBatch)}\\[2pt]
    {\footnotesize\color{gray!60!black}Transcription établie d'après le fac-similé numérisé
     par l'Université de Montpellier.\\
     Les lectures incertaines sont soulignées, les illisibles notées [\,\ldots\,],
     les ajouts entre crochets.}
  \end{center}
  \vspace{1em}}

\endinput


\folder{1}
\pages{1}{149}
\dating{1953}
\watermark{Édition de démonstration\\interprétation personnelle de l'œuvre}
\foldertitle{Analyse fonctionnelle : notes manuscrites (s.d.), lettre (1953) — lecture modernisée du dossier entier}

\begin{document}

\section*{Résumé}

\begin{resume}
Le carton porte un mot de l'inventaire, « analyse fonctionnelle », et une
date, 1953. Dedans, cent quarante-neuf pages qui tournent toutes autour
d'une même idée, simple à dire : mesurer la taille d'un opérateur non par
un seul nombre, sa norme, mais par toute une suite de nombres, et savoir ce
que cette suite fait quand on ajoute deux opérateurs ou qu'on les compose.

Voici l'image. Un opérateur $u$ sur un espace de Hilbert --- disons une
matrice infinie --- est \emph{compact} s'il écrase presque tout : il existe
une suite décroissante $s_1(u) \geq s_2(u) \geq \cdots$, tendant vers zéro,
telle que $u$ se décompose en une somme de « rangs un » de tailles $s_1$,
$s_2$, \dots\ Ce sont ses \emph{valeurs singulières}. Elles disent combien
de dimensions $u$ retient vraiment, et à quel prix : $s_1$ est la norme,
$\sum s_i$ est la trace de $|u|$, $\sum s_i^2$ est la norme de
Hilbert--Schmidt. Toute fonction raisonnable de la suite est une manière
de mesurer $u$, et les classes d'opérateurs ainsi définies --- les classes
de Schatten, $\sum s_i^p < \infty$ --- sont pour les opérateurs ce que les
espaces $\ell^p$ sont pour les suites.

Le problème est que les valeurs singulières se comportent mal. Celles de
$u + v$ ne sont pas des fonctions simples de celles de $u$ et de $v$, ni
celles de $uv$. Ce qui reste vrai, et c'est la découverte de Weyl en 1949,
c'est une \emph{comparaison en moyenne} : les sommes partielles $s_1 +
\cdots + s_n$ de $u + v$ sont majorées par celles de $u$ plus celles de $v$,
et les produits partiels $s_1 \cdots s_n$ de $uv$ par ceux de $u$ fois
ceux de $v$. Et il y a un lemme, dû à Hardy, Littlewood et Pólya, qui dit
qu'une comparaison en moyenne suffit à comparer toute quantité convexe et
symétrique. Mettez les deux bout à bout : toute norme raisonnable sur les
valeurs singulières donne une norme sur les opérateurs, avec l'inégalité
triangulaire, l'inégalité de Hölder, la dualité. C'est la théorie de
Schatten, refaite en dix pages, aux pages 81 à 104.

Le dossier ne s'arrête pas là, et c'est ce qui le rend intéressant. Il
contient deux versions de la même théorie, et la plus ancienne, aux pages
123 à 149, ne connaît pas encore le lemme de Weyl. Elle passe par une
autre route, celle du \emph{déterminant} : pour un opérateur à trace,
$\det(1 + zu)$ est une fonction entière de $z$, dont les zéros sont les
inverses des valeurs propres et dont la croissance mesure les valeurs
singulières. Toute la théorie des fonctions entières --- factorisation de
Weierstrass et de Hadamard, genre, ordre, formule de Jensen --- est alors
disponible, et le carton la rappelle aux pages 106 à 118 pour s'en servir.
Ce détour donne plus que la théorie de Schatten : il atteint les
\emph{classes pour $p < 1$}, qui ne sont plus des espaces normés, et il y
démontre l'inégalité $\sum s_i(u+v)^p \leq \sum s_i(u)^p + \sum s_i(v)^p$
et l'inégalité de Hölder pour tous les exposants. Il démontre aussi, en
passant, que le déterminant est bien le produit $\prod(1 + \lambda_i)$ des
valeurs propres, ce qui est la formule de la trace ; et une inégalité,
$\det(1 + |u + v|) \leq \det(1 + |u|)\det(1 + |v|)$, que la littérature
date de 1975.

Il y a une troisième couche, la plus ambitieuse, aux pages 6 à 35 : refaire
tout cela dans une \emph{algèbre de von Neumann} munie d'une trace, là où
il n'y a plus de suite de valeurs singulières parce que les projecteurs
n'ont plus une dimension entière mais une trace réelle. La suite devient
une fonction décroissante $\varphi_A(t)$ sur $]0,\infty[$, le
\emph{réarrangement} de $A$, et le lemme de Weyl devient une inégalité
entre intégrales. On y trouve le déterminant de Fuglede et Kadison, une
inégalité $\Delta_{AB}(t) \leq \Delta_A(t)\Delta_B(t)$ sur des déterminants
partiels, et l'inégalité de Hölder non commutative, en trois lignes. Ce
sont les nombres singuliers généralisés de Fack et Kosaki, et les espaces
$L^p$ non commutatifs, tels qu'on les écrit depuis les années 1980.

Trois feuillets à part, pages 3 à 5, identifient les multiplicateurs d'une
$C^{*}$-algèbre dans son bidual à ceux qu'elle a dans n'importe quelle
représentation ; deux pages, 119 à 121, étudient une fonction holomorphe
sur $\ell^1$ qui converge partout mais n'est bornée sur aucune boule de
rayon un.

Le ton est celui d'un cahier de travail. « Je ne connais qu'une méthode
analytique », écrit-il d'une inégalité qu'une page plus jeune du même
carton obtient en deux lignes ; « c'est une inégalité remarquable, mais
n'est pas celle que nous cherchons », d'une étape qui s'avère être la
bonne ; « est-ce utile ? » et « pas d'importance » en marge de deux
corollaires. Et, sous un trait au crayon, une question sur les valeurs
propres et les valeurs singulières que le texte ne tranche pas, et que
cette lecture tranche en note. Les noms à chercher ensuite sont ceux de
Weyl, Ky Fan et Horn pour les inégalités, de Schatten pour les classes,
de Fuglede--Kadison et Fack--Kosaki pour les algèbres de von Neumann, et
de la théorie de Fredholm pour le déterminant.

\keywords{singular values, Schatten--von Neumann classes, symmetric gauge function, unitarily invariant norm, Ky Fan maximum principle, Ky Fan inequalities, Weyl majorant theorem, Horn inequality, Hardy--Littlewood--Pólya majorization, decreasing rearrangement, rearrangement-invariant functional, Krein--Šmulian theorem, Hölder inequality for Schatten classes, quasi-normed operator ideal, McCarthy--Rotfel'd inequality, Seiler--Simon determinant inequality, Fredholm determinant, Fredholm minor, Lidskii trace formula, exterior power, antisymmetric Fock space, Hadamard inequality, Hilbert tensor product, Hilbert--Schmidt operator, trace class, Schatten duality, Köthe dual, generalized singular numbers, Fuglede--Kadison determinant, semifinite trace, noncommutative $L^p$ space, von Neumann trace inequality, multiplier algebra, idealizer, entire function of genus zero, Hadamard factorization, exponent of convergence, Jensen formula, Gâteaux holomorphy, radius of boundedness, Kadec--Klee property, Birman--Koplienko--Solomyak inequality, Mirsky inequality}
\end{resume}

\section*{Le fil du dossier, et les conventions}

Cent quarante-neuf pages, huit lots de transcription, au moins quatre
papiers et trois encres : le carton n'est pas un texte mais un chantier, et
il faut dire d'abord ce qu'il construit, dans quel ordre il le construit, et
avec quels mots.

\pagerange{2}{2}
\subsection*{Les stations}

\begin{itemize}
  \item \textbf{Page 2} --- le plan, de sa main, en sept numéros :
    réarrangements de fonctions, fonctions de Weyl, inégalités, mesures
    spectrales d'un opérateur, compléments sur le déterminant, les
    inégalités pour $AB$, autres inégalités.
  \item \textbf{Pages 3 à 5} --- un fragment étranger au plan, sur les
    multiplicateurs d'une $C^{*}$-algèbre dans son bidual.
  \item \textbf{Pages 6 à 35} --- les numéros 1 à 6 du plan : réarrangements
    décroissants de fonctions, fonctions de Pólya et de Weyl, inégalités de
    majoration, puis tout cela transporté dans une algèbre de von Neumann
    munie d'une trace --- réarrangement spectral d'un opérateur, déterminant,
    inégalité $\Delta_{AB}(t) \leq \Delta_A(t)\,\Delta_B(t)$ et ses
    conséquences ; une seconde mouture des propositions clefs aux pages 34
    et 35.
  \item \textbf{Pages 37 à 41} --- des « Préliminaires » sur les opérateurs
    compacts d'un espace de Hilbert : décomposition polaire, valeurs
    singulières, inégalités de Ky Fan, opérateurs à trace, dual des
    opérateurs compacts.
  \item \textbf{Pages 43 à 57} --- « 2. Produits tensoriels hilbertiens » :
    l'espace conjugué, $E \mathbin{\hat\otimes}_2 F$, les opérateurs de
    Hilbert--Schmidt et les opérateurs à trace, leurs normes et leur dualité.
  \item \textbf{Pages 58 à 69} --- « 3. Produits extérieurs hilbertiens » :
    $\Lambda^n E$ comme espace de Hilbert, l'inégalité de Hadamard, les
    puissances extérieures d'un opérateur, et les inégalités de Weyl
    $|\lambda_1 \cdots \lambda_n| \leq s_1 \cdots s_n$ et de Horn
    $\prod_{i \leq n} s_i(uv) \leq \prod_{i \leq n} s_i(u)\,s_i(v)$.
  \item \textbf{Pages 70 à 79} --- « Puissances extérieures d'opérateurs de
    Fredholm » en deux états, la mise au net (70 à 73) et son brouillon (76 à
    79) : $\det(1+zu) = \prod(1+z\lambda_i)$, la formule de trace
    $\operatorname{Tr} u = \sum \lambda_i$, et les inégalités
    $\det(1+r|u+v|) \leq \det(1+r|u|)\det(1+r|v|)$. Entre les deux, aux pages
    74 et 75, la sous-additivité de $\sum f(s_i)$ et une question ; page 80,
    une feuille sur les fonctions convexes.
  \item \textbf{Pages 81 à 92} --- « 4. Inégalités de convexité » : le lemme
    de majoration de Weyl et ses trois théorèmes de convexité dans l'espace
    de Hilbert, dont l'inégalité de Hölder et l'inégalité triangulaire des
    normes de Schatten.
  \item \textbf{Pages 93 à 104} --- « 5. Classes remarquables d'opérateurs »
    et « Espaces de Schatten » : normes symétriques, espaces $\ell^N$ et
    $\ell_N$, norme polaire, dualité, puis les classes
    $\mathcal{S}^N(E,F)$ et $\mathcal{S}_N(E,F)$ et le théorème récapitulatif
    en cinq points.
  \item \textbf{Pages 106 à 118} --- des rappels sur les fonctions entières :
    Weierstrass, Hadamard--Borel, genre et ordre, exposant de convergence,
    fonction de comptage des zéros et formule de Jensen.
  \item \textbf{Pages 119 à 121} --- « Sur une certaine fonction holomorphe
    sur l'espace de Banach $\ell^1$ ».
  \item \textbf{Pages 123 à 125} --- cinq corollaires encadrés, jusqu'à
    $\det(1+|A+B|) \leq \det(1+|A|)\det(1+|B|)$.
  \item \textbf{Pages 126 à 139} --- un « Théorème » sur les classes
    $\mathcal{S}_p(H)$ pour $0 < p \leq \infty$, et sa démonstration
    analytique complète.
  \item \textbf{Pages 141 à 147} --- la convergence dans $\mathcal{S}_p$ :
    critère pour les suites d'hermitiens, continuité de $A \mapsto |A|^r$,
    un lemme sur les petites valeurs singulières, et des questions.
  \item \textbf{Pages 148 et 149} --- l'opérateur $R(u) = \det(1+u)\,(1+u)^{-1}$
    et ses bornes.
\end{itemize}

\subsection*{L'ordre du carton n'est pas celui de l'écriture}

Trois indices, pris dans les pages elles-mêmes, fixent une chronologie que
la cote ne respecte pas. Page 129, pour l'inégalité $\sum |\lambda_i|^p \leq
\sum s_i^p$, il écrit « je ne connais qu'une méthode analytique », et en donne
une par les zéros de $\det(1 - zA)$ ; page 88, la même inégalité est un
corollaire en deux lignes du lemme de Weyl. Page 93, il annonce « un exposé de
la théorie de Schatten, simplifié par les inégalités de convexité de Weyl ».
Page 132, il note qu'une preuve directe de l'inégalité de Hölder « suffirait
pour démontrer directement, comme nous le verrons plus bas, la formule (6)
dans le cas général », et les pages 81 à 92 sont ce « plus bas ». Les pages
123 à 149 sont donc l'état ancien, celui d'avant le lemme de Weyl ; les pages
37 à 104 l'état récent ; et les rappels sur les fonctions entières des pages
106 à 118 servent la démonstration de la page 79. Cette lecture suit l'ordre
des pages, mais dit à chaque fois lequel des deux états on lit.

\subsection*{Les conventions, valables pour tout le dossier}

$E$, $F$, $G$, $H$ sont des espaces de Hilbert complexes ; $L(E,F)$ les
opérateurs bornés, $L_0(E,F)$ les compacts. Pour $u$ compact, $|u| = (u^{*}u)^{1/2}$
et $s_1(u) \geq s_2(u) \geq \cdots$ sont ses \emph{valeurs singulières}, les
valeurs propres de $|u|$ rangées en décroissant, répétées selon leur
multiplicité.\footnote{Le dossier les note $\rho_n(u)$ des pages 37 à 104 et
$r_\nu(A)$ aux pages 146 et 147 ; on écrit $s_n$, qui est l'usage. Les
valeurs propres $\lambda_i(u)$ sont rangées par modules décroissants, avec
leur multiplicité algébrique.} $\mathcal{S}_p(E,F)$, $0 < p \leq \infty$, est la classe
des $u$ compacts avec $\|u\|_p = (\sum s_i(u)^p)^{1/p} < \infty$, et
$S_p(u) = \|u\|_p^p$ ; $\mathcal{S}_\infty = L_0$ avec la norme
d'opérateur.\footnote{Ses notations : $L^p(E,F)$, $L^{(1)}$, $L^{(2)}$ aux
pages 37 à 104, $\mathcal{L}^p(H)$ aux pages 126 à 147. Les
$\mathcal{S}_p$ sont les classes de Schatten--von Neumann ; le nom vient de
Schatten (1950) et de von Neumann (1937), que la page 93 nomme.}
Un point de vocabulaire domine tout : ce que ces pages appellent
\emph{opérateur de Fredholm}, c'est un opérateur \emph{à trace}, un élément
de $\mathcal{S}_1$ --- l'opérateur pour lequel le déterminant de Fredholm
$\det(1+zu)$ a un sens. Ce n'est pas l'opérateur « de Fredholm » de la
théorie de l'indice, et on écrit ici \emph{à trace} partout.\footnote{L'usage
est celui de sa thèse et de \emph{La théorie de Fredholm} (1956), où
« application de Fredholm » désigne une application nucléaire. Les renvois du
dossier à « Chap.\ 1, § 2, n° 2 » et « Introduction, 2, formule (8) » (pages
63, 64, 76, 77) visent un texte qui n'est pas dans le carton ; la structure
annoncée est celle de ce mémoire, sans qu'on puisse l'affirmer.}
$\operatorname{Tr}$ est la trace, $\det(1+u) = \prod (1+\lambda_i(u))$ le
déterminant de Fredholm, $\alpha_n(u)$ le coefficient de $z^n$ dans
$\det(1+zu)$, et $\Lambda^n u$ la puissance extérieure.

Pour les fonctions (pages 6 à 17) : $(M,m)$ est un espace mesuré,
$\varphi_f$ le \emph{réarrangement décroissant} d'une fonction mesurable
positive $f$ --- ce que l'on note aujourd'hui $f^{*}$ ---, et
$\Phi_f(t) = \int_0^t \varphi_f$, $\Psi_f(t) = \int_0^t \log\varphi_f$,
$\Delta_f = \exp\Psi_f$. Ces lettres sont celles du dossier et on les garde,
parce que les mêmes reviennent pour les opérateurs.

Pour les algèbres (pages 18 à 35) : $\mathcal{M}$ est une algèbre de von
Neumann, $\tau$ une trace normale, semi-finie et fidèle sur $\mathcal{M}^{+}$,
$\mathfrak{a}$ son idéal de définition, $\mathfrak{b}$ l'adhérence en norme de
$\mathfrak{a}$, $\mathcal{P}$ l'ensemble des projecteurs de
$\mathcal{M}$.\footnote{Il écrit $\mathcal{A}$, $T$, $\underline{a}$,
$\underline{b}$ ; « idéal de définition » est le mot de Dixmier (1953), dont
ces pages prennent la théorie des traces.} Le \emph{réarrangement spectral}
$\varphi_A$ d'un $A \in \mathfrak{b}$ est ce que l'on appelle depuis
Fack et Kosaki (1986) la fonction $t \mapsto \mu_t(A)$ des \emph{nombres
singuliers généralisés} ; $\Delta_A(t) = \exp\int_0^t \log\varphi_{|A|}$ et
$\Delta(A) = \Delta_A(\tau(1))$ est le déterminant de Fuglede--Kadison
(1952).

Deux choses annoncées ne sont pas dans le carton. Le plan de la page 2
promet un « N° 7, autres inégalités » qui n'existe nulle part. Et la page
126 annonce en marge « ② Dual de $\mathcal{L}^p$ ; réflexivité » sans que ce
point soit traité dans ce théorème-là ; il l'est, pour les normes de
Schatten générales, aux pages 98 à 104. L'inventaire mentionne enfin une
« lettre (1953) » : aucune page transcrite ne porte d'adresse ni de signature,
et cette lecture ne désigne pas laquelle serait la lettre.

\pagerange{3}{5}
\section*{Les multiplicateurs d'une $C^{*}$-algèbre (pages 3 à 5)}

Trois pages pliées, d'une encre à part, qui ne tiennent au reste que par le
sujet des algèbres d'opérateurs. Soit $C$ une $C^{*}$-algèbre et $C''$ son
bidual, muni de sa structure d'algèbre de von Neumann enveloppante ; $C$ s'y
plonge. Un endomorphisme linéaire continu $T$ de $C$ qui commute aux
translations à droite, $T(xy) = T(x)\,y$, a un bitransposé $T''$ qui commute
aux mêmes translations sur $C''$, donc est de la forme $x \mapsto ux$ pour un
$u \in C''$ ; ces $u$ forment l'algèbre $\Gamma_g$ des $u \in C''$ tels que
$uC \subset C$, et $T \mapsto u$ est un isomorphisme d'algèbres normées de
l'algèbre de ces endomorphismes sur $\Gamma_g$.\footnote{Ce sont les
\emph{multiplicateurs à gauche} de $C$. La page dit « permutables aux
translations à gauche » et écrit $T(xy) = T(x)y$, qui est la commutation aux
translations à droite ; on suit la formule. La prose de ces trois pages est
en grande partie perdue, les formules non.} Soit $\Gamma$ l'ensemble des
$u \in C''$ tels que $uC \subset C$ et $Cu \subset C$ : c'est une
sous-$C^{*}$-algèbre de $C''$, l'algèbre des \emph{multiplicateurs} de $C$,
$M(C)$ en notation d'aujourd'hui.

\emph{Théorème.} Soit $R$ une algèbre de von Neumann contenant $C$ comme
sous-algèbre faiblement dense, et $\Gamma_1$ l'ensemble des $x \in R$ tels
que $xC \subset C$ et $Cx \subset C$. Alors $\Gamma_1$ et $\Gamma$ sont
canoniquement isomorphes.

Voici les deux applications. Pour $u \in \Gamma_1$, la translation
$x \mapsto ux$ est un endomorphisme de $C$ commutant aux translations à
droite, donc définit $\varphi(u) \in \Gamma_g$ ; comme $u^{*} \in \Gamma_1$
aussi, $\varphi(u) \in \Gamma$, et $\varphi$ est un morphisme d'algèbres
involutives. On a $\varphi(u) = \lim ux = \lim xu$ au sens faible, $x$
parcourant une unité approchée de $C$, soit
$\langle \varphi(u), f\rangle = \lim \langle ux, f\rangle$ pour $f \in C'$.
Dans l'autre sens, $C$ étant faiblement dense dans $R$, l'inclusion se
prolonge en un morphisme normal $\psi : C'' \to R$, et $\psi(\Gamma) \subset
\Gamma_1$ ; le prédual $R_{*}$ s'identifie à une partie de $C'$ par
restriction. Alors $\psi\varphi = 1$ sur $\Gamma_1$ : pour $u \in \Gamma_1$,
$v = \varphi(u)$ et $f \in R_{*}$, $\langle \psi v, f\rangle = \langle v,
f|_C\rangle = \langle u, f\rangle$. Et $\varphi\psi = 1$ sur $\Gamma$ : il
suffit qu'un $x \in \Gamma$ avec $\psi(x) = 0$ soit nul, or pour $y \in C$
on a $xy \in C$, donc $xy = \psi(xy) = \psi(x)\,y = 0$, et $x = 0$ puisque
$C$ a une unité approchée. Donc $\varphi$ est un isomorphisme de $\Gamma_1$
sur $\Gamma$, d'inverse $\psi$.\footnote{C'est l'identification de $M(C)$ au
\emph{idéaliseur} de $C$ dans toute représentation fidèle non dégénérée ---
énoncé standard depuis Busby (1968) et Johnson (1964) sur les
centralisateurs, et que l'on trouve dans Pedersen, \emph{$C^{*}$-algebras
and their automorphism groups}, 3.12. Rien n'indique que ces pages aient été
publiées.}

\pagerange{6}{10}
\section*{Réarrangements décroissants et déterminant d'une fonction (pages 6 à 10)}

C'est le « N° 1 » du plan. Soit $f$ une fonction mesurable complexe sur
$(M,m)$. On dit que $f$ \emph{admet une mesure spectrale} si l'image de $m$
par $f$ est une mesure de Radon sur $\mathbf{C}^{*} = \mathbf{C} \smallsetminus
\{0\}$, autrement dit si $|f|^{-1}([a,b])$ est de mesure finie pour
$0 < a < b$ ; cette mesure est notée $\mu_f$, et
\[
  \langle \alpha, \mu_f\rangle = \int_M \alpha\circ f\, dm
\]
pour $\alpha$ continue à support compact dans $\mathbf{C}^{*}$. Elle est
positive, portée par $\mathbf{R}^{*}$ si $f$ est réelle, par
$\mathbf{R}^{*}_{+}$ si $f$ est positive. Deux fonctions, sur deux espaces
mesurés, sont \emph{spectralement équivalentes} --- on dit aujourd'hui
\emph{équimesurables} --- si elles ont même mesure spectrale.

\emph{Proposition 1.} Une fonction positive décroissante $\varphi$ sur
$\mathbf{R}^{*}_{+}$ admet une mesure spectrale si et seulement si elle tend
vers $0$ à l'infini, et alors $\int_1^\infty d\mu_\varphi < \infty$.
Réciproquement, toute mesure positive $\mu$ sur $\mathbf{R}^{*}_{+}$ avec
$\int_1^\infty d\mu < \infty$ est la mesure spectrale d'une fonction positive
décroissante et d'une seule, à un ensemble négligeable près : celle qui
vérifie, pour presque tout $t$,
\[
  t = \int_{\varphi(t)}^{\infty} d\mu(s),
\]
c'est-à-dire l'inverse généralisée de $s \mapsto \mu(]s,\infty[)$.

\emph{Corollaire.} Une fonction mesurable positive $f$ sur $M$ est
équimesurable à une fonction décroissante sur $\mathbf{R}^{*}_{+}$ si et
seulement si $f^{-1}([a,\infty[)$ est de mesure finie pour tout $a > 0$ ;
cette fonction décroissante est alors unique, c'est le \emph{réarrangement
décroissant} $\varphi_f$ de $f$, et deux telles fonctions sont équimesurables
si et seulement si elles ont même réarrangement. On a, en identifiant une
partie mesurable $E$ à sa fonction caractéristique,
\[
  \varphi_f(t{-}0) = \inf_{m(E) < t} \|(1-E)f\|_\infty, \qquad
  \varphi_f(t{+}0) = \inf_{m(E) \leq t} \|(1-E)f\|_\infty .
\]
\footnote{La page 7 écrit ces deux formules avec $\leq$ pour la limite à
gauche et $<$ pour la limite à droite, le premier signe étant surchargé
d'après la transcription. C'est l'inverse : l'infimum sur les $E$ de mesure
$< t$ porte sur moins d'ensembles, donc est le plus grand des deux, donc
est la limite à gauche. La page 26 écrit la formule pour les opérateurs avec
le signe strict et la limite à gauche, ce qui est cohérent avec ce qu'on
donne ici.}
De là les règles, pour $f$, $g$ positives admettant un réarrangement :
$\varphi_f \leq \varphi_g$ si $f \leq g$ ; $\varphi_{\lambda f} =
\lambda\varphi_f$ pour $\lambda \geq 0$ ; $\varphi_{\alpha(f)} =
\alpha(\varphi_f)$ pour $\alpha$ croissante sur $\mathbf{R}_{+}$ nulle en
$0$, en particulier $\varphi_{f^\alpha} = \varphi_f^\alpha$ ;
$\varphi_{f+g} \leq \varphi_f + \varphi_g$ ; $\varphi_{fg} \leq
\varphi_f\varphi_g$.\footnote{Les deux dernières se lisent au sens de la
\emph{sous-majorisation} : $\int_0^t \varphi_{f+g} \leq \int_0^t
(\varphi_f + \varphi_g)$ pour tout $t$, et de même en produit logarithmique
pour $fg$. Ponctuellement, $\varphi_{f+g}(t) \leq \varphi_f(t) + \varphi_g(t)$
est faux --- $f$ et $g$ indicatrices d'ensembles disjoints de même mesure $1$
donnent $\varphi_{f+g} = 1$ sur $]0,2[$ et $\varphi_f + \varphi_g = 0$ sur
$]1,2[$ --- et c'est bien sous forme intégrée que la page 7 les utilise
ensuite (formules (10) et (11)). Le passage qui les commente est en grande
partie illisible.} Une bulle marginale note ce qui justifie tout : $\int_M
f\,dm = \int_0^\infty \varphi_f(t)\,dt$, et de même pour toute fonction
symétrique de $f$, $N_p(f) = N_p(\varphi_f)$.

Supposons $\varphi_f$ sommable au voisinage de $0$, ce qui revient à
$\int^\infty s\,d\mu_f(s) < \infty$, et posons
\[
  \Phi_f(t) = \int_0^t \varphi_f(s)\,ds, \qquad
  \Psi_f(t) = \int_0^t \log\varphi_f(s)\,ds, \qquad
  \Delta_f(t) = \exp\Psi_f(t) .
\]
$\Phi_f$ est continue, croissante et concave. $\Psi_f$ a un sens dès que
$\log f$ est sommable sur $\{f \geq 1\}$ ; elle est concave, à valeurs dans
$[-\infty, +\infty[$, et vaut $-\infty$ à partir de $t_0 = m(\{f \neq 0\})$ si
celui-ci est fini. $\Delta_f$ est continue, croissante puis nulle à partir de
$t_0$. Les règles ci-dessus donnent les inégalités correspondantes, dont les
plus importantes sont
\[
  \Psi_{fg} \leq \Psi_f + \Psi_g, \qquad \Delta_{fg} \leq \Delta_f\,\Delta_g .
\]
Pour une fonction mesurable quelconque on pose enfin
\[
  \Delta(f) = \Delta_{|f|}(m(M)) = \exp\int_M \log|f|\,dm,
\]
qui existe, avec $0 \leq \Delta(f) < \infty$, dès que $\log|f|$ est sommable
sur $\{|f| > 1\}$ ; $\Delta$ a les propriétés d'un module de déterminant,
$\Delta(fg) = \Delta(f)\Delta(g)$ et $\Delta(1) = 1$.\footnote{C'est la
moyenne géométrique de $|f|$, le déterminant de Fuglede--Kadison de l'algèbre
commutative $L^\infty(M)$ : la page 27 le transportera tel quel dans une
algèbre de von Neumann.} Si $m(M) < \infty$ et $f \geq 0$, on a
$\varphi_{1+f} = 1 + \varphi_f$ sur $]0, m(M)[$, ce qui permet de définir
$\Delta_{1+f}(t) = \exp\int_0^t \log(1+\varphi_f)$ dès que $\Phi_f$ existe,
avec $\Delta_{1+f}(m(M)) = \Delta(1+f)$.

Pour une fonction \emph{réelle} $f$ sur un $M$ de mesure finie, la même
proposition 1 --- appliquée à $\mu_f$ sur $\mathbf{R}^{*}$, qui porte la masse
totale $m(M)$ --- donne une fonction décroissante $\varphi_f$ et une seule sur
$]0, m(M)[$ équimesurable à $f$, encore appelée réarrangement décroissant.
Le résultat suivant est classique :

\emph{Proposition 2} (Hardy--Littlewood). Pour $f$, $g$ réelles mesurables
sur $M$ de mesure finie, telles que $fg$ soit sommable,
\[
  \int_M fg\,dm \leq \int_0^{m(M)} \varphi_f\,\varphi_g\,ds .
\]
\footnote{La page ajoute, d'une lecture incertaine, que la sommabilité de
$\varphi_f\varphi_g$ est requise, ce qui est bien la condition pour que le
second membre soit fini ; l'inégalité vaut dans $]-\infty,+\infty]$ sans
elle.}

\pagerange{11}{15}
\section*{Fonctions de Pólya et fonctions de Weyl (pages 11 à 15)}

Le « N° 2 » construit les fonctionnelles pour lesquelles les réarrangements
sont la bonne variable. Soit $I = ]0,\omega[$ avec $0 < \omega \leq \infty$,
et $K^\infty(I)$ l'espace des fonctions mesurables bornées à support de
mesure finie dans $I$.

\emph{Définition 1.} Une \emph{fonction de Pólya} sur $I$ est une fonction
$P : K^\infty(I) \to \mathbf{R}$ telle que
\begin{itemize}
  \item[$P_1$)] $P$ est \emph{symétrique} : $P(f) = P(g)$ si $f$ et $g$ sont
    équimesurables ;
  \item[$P_2$)] $P$ est convexe ;
  \item[$P_3$)] pour tout $K \subset I$ de mesure finie et toute suite
    bornée $(f_i)$ dans $L^\infty(K)$ convergeant en mesure vers $f$, on a
    $P(f) \leq \liminf P(f_i)$.
\end{itemize}
En fait $P_3$ équivaut à la condition apparemment plus forte $P_3'$ : $P$ est
semi-continue inférieurement pour la topologie faible $\sigma(L^\infty(K),
L^1(K))$ sur les parties bornées de $L^\infty(K)$. Car pour tout $a$
l'ensemble convexe $A = \{f \in L^\infty(K) : P(f) \leq a\}$ est, par $P_3$,
fermé pour la convergence en mesure sur les boules de $L^\infty(K)$, donc
fermé pour la norme de $L^1(K)$ sur ces boules ; or un convexe de
$L^\infty(K) = (L^1(K))'$ dont l'intersection avec chaque boule est fermée
pour la norme de $L^1$ --- \emph{a fortiori} pour $\sigma(L^\infty, L^1)$ ---
est faiblement fermé.\footnote{C'est le théorème de Krein--Šmulian, que la
page 11 invoque en citant, dans un passage biffé, « Bourbaki » et
« Dieudonné » ; le nom moderne est fourni ici.} Une fonction sur
$K^\infty(]0,\infty[)$ est de Pólya si et seulement si sa restriction à chaque
$L^\infty(]0,\omega[)$, $\omega < \infty$, l'est.

\emph{Définition 2.} Une \emph{fonction de Weyl} sur les fonctions positives
de $K^\infty(I)$ est une fonction $W$ telle que $f \mapsto W(e^{f})$ soit
une fonction de Pólya. Les fonctions de Weyl et les fonctions de Pólya se
correspondent donc par $W(g) = P(\log g)$, $P(f) = W(e^f)$.\footnote{Ces
deux noms sont les siens et ne sont pas passés dans l'usage. Une fonction de
Pólya est une fonctionnelle convexe, \emph{invariante par réarrangement} et
semi-continue ; les fonctions de Weyl sont leurs transformées par le
logarithme, celles que le lemme de Weyl de la page 81 fait comparer par les
produits partiels plutôt que par les sommes partielles.}

\emph{Représentation.} Soit $K \subset I$ de mesure finie et $P$ une fonction
de Pólya sur $L^\infty(K)$. L'épigraphe $\{(\lambda, f) : \lambda \geq P(f)\}$
est convexe ($P_2$) et fermé ($P_3'$) dans $\mathbf{R} \times L^\infty(K)$
pour la topologie faible, donc, par Hahn--Banach, intersection des
demi-espaces fermés qui le contiennent : $P(f) = \sup_i (L_i(f) + a_i)$ avec
$L_i \in L^1(K)$ et $a_i \in \mathbf{R}$. Par $P_1$, $P(f) \geq L_i(g) + a_i$
pour toute $g$ équimesurable à $f$, et le supremum de $\int L_i g$ sur ces
$g$ est $\int \varphi_{L_i}\varphi_f$ par la proposition 2 --- le
sup étant atteint --- d'où, en posant $\varphi_i = \varphi_{L_i}$,
\[
  P(f) = \sup_{i \in I}\Bigl( a_i + \int_0^{\omega} \varphi_i\,\varphi_f \Bigr),
\]
les $\varphi_i$ étant des fonctions décroissantes sommables sur $]0, |K|]$.
Réciproquement toute fonction de cette forme est de Pólya, et elle est
croissante si les $\varphi_i$ sont positives. On peut prendre les $\varphi_i$
nulles sur $[|K|, \infty[$, et la formule s'écrit $P(f) = P^{*}(\varphi_f)$ où
$P^{*}(\varphi) = \sup_i (a_i + \int_0^\infty \varphi_i\varphi)$ est une
fonction de Pólya sur $\mathbf{R}^{*}_{+}$.\footnote{En langage
d'aujourd'hui : une fonctionnelle convexe s.c.i.\ invariante par
réarrangement est le sup des fonctionnelles affines invariantes par
réarrangement, et celles-ci sont données par le produit scalaire des
réarrangements. C'est l'argument qui, pour les normes, donne la
représentation d'une norme invariante par réarrangement par sa norme duale
(Lorentz, Luxemburg) ; ces pages le font sans le nommer, et deux ans avant
Lorentz, mais rien n'y renvoie à une publication.}

\emph{Exemples de fonctions de Weyl.} Soit $\alpha$ une fonction sur
$\mathbf{R}_{+}$, nulle en $0$, telle que $s \mapsto \alpha(e^s)$ soit
convexe. Alors $W_\alpha(f) = \int_0^\omega \alpha(\varphi_f)$ est une
fonction de Weyl, croissante si et seulement si $\alpha$ l'est. Pour
$\alpha_p(s) = s^p$, $0 < p < \infty$, on trouve $W_{\alpha_p}(f) = \int
|f|^p = N_p(f)^p$. Une fonction de Weyl croissante sur les fonctions
positives de $K^\infty$ s'étend à toute fonction mesurable positive $f$ sur
un espace mesuré quelconque par $W(f) = \sup\{W(\varphi_g) : 0 \leq g \leq f,\ g \in K^\infty\}$, et l'on a $W(f) = W(\varphi_f)$, $W(f) = W(\varphi_{|f|})$
pour $f$ complexe.

\pagerange{16}{17}
\section*{Les inégalités de majoration (pages 16 et 17)}

Le « N° 3 » est court, et il contient la moitié de tout ce qui suit. Soit
$0 < \omega < \infty$.

\emph{Proposition 4.} Soient $f$, $g$ décroissantes et localement sommables
sur $]0,\omega[$ telles que
\[
  \int_0^t f \leq \int_0^t g \qquad \text{pour } 0 \leq t \leq \omega .
\]
Alors $P(f) \leq P(g)$ pour toute fonction de Pólya \emph{croissante} $P$ sur
$L^\infty(]0,\omega[)$. (La réciproque est triviale : pour tout $t$,
$f \mapsto \int_0^t \varphi_f$ est une fonction de Pólya croissante.)

Par la représentation du numéro 2, il suffit de prouver $\int_0^\omega f\varphi
\leq \int_0^\omega g\varphi$ pour $\varphi$ positive, décroissante et sommable,
ce qui résulte de l'hypothèse par intégration par parties : avec
$F(t) = \int_0^t f$,
\[
  \int_0^\omega f\varphi = \varphi(\omega)F(\omega) + \int_0^\omega F\,(-d\varphi),
\]
et les deux termes croissent avec $F$ puisque $\varphi(\omega) \geq 0$ et
$-d\varphi \geq 0$.

\emph{Remarque.} Si de plus $\int_0^\omega f = \int_0^\omega g$, le premier
terme est le même pour $f$ et $g$ et l'on n'a plus besoin de $\varphi \geq 0$ :
la proposition vaut alors pour toute fonction de Pólya, croissante ou
non.\footnote{C'est le théorème de Hardy, Littlewood et Pólya sur la
majorisation, en dimension infinie : $f \prec_w g$ (sous-majorisation faible)
entraîne $\Phi(f) \leq \Phi(g)$ pour $\Phi$ convexe, symétrique et
croissante, et $f \prec g$ (majorisation) l'entraîne pour $\Phi$ convexe
symétrique. La page 81 en redonnera la version finie, attribuée à Weyl.}

\emph{Corollaire 1.} Soient $f$, $g$ positives sur $M$ telles que
$\Delta_f(t) \leq \Delta_g(t) < \infty$ pour $0 \leq t \leq \omega$. Alors
$W(f) \leq W(g)$ pour toute fonction de Weyl croissante $W$ sur les fonctions
positives de $K^\infty(]0,\omega[)$. On se ramène à $M = ]0,\omega[$ et à
$f = \varphi_f$, $g = \varphi_g$, puis on applique la proposition 4 à
$\log f$, $\log g$ et à la fonction de Pólya croissante $h \mapsto W(e^h)$.

\pagerange{18}{27}
\section*{Réarrangement spectral d'un opérateur dans une algèbre de von Neumann (pages 18 à 27)}

Le « N° 4 » transporte le numéro 1 dans une algèbre de von Neumann
$\mathcal{M}$ munie d'une trace $\tau$ normale, semi-finie et fidèle : une
fonction $\tau : \mathcal{M}^{+} \to [0,+\infty]$, additive, positivement
homogène, invariante par les unitaires, telle que $\tau(\sup A_i) = \sup
\tau(A_i)$ pour toute famille filtrante croissante majorée, telle que tout
$A > 0$ majore un $B > 0$ avec $\tau(B) < \infty$, et telle que $\tau(A) = 0$
entraîne $A = 0$. Les $A \geq 0$ avec $\tau(A) < \infty$ forment le cône
positif d'un idéal bilatère involutif $\mathfrak{a}$, sur lequel $\tau$ se
prolonge en une forme linéaire vérifiant $\tau(AB) = \tau(BA)$ pour
$A \in \mathfrak{a}$, $B \in \mathcal{M}$. On note $\mathfrak{b}$
l'adhérence en norme de $\mathfrak{a}$ ; c'est un idéal bilatère involutif
fermé, égal à $\mathcal{M}$ si et seulement si $\tau$ est finie.\footnote{La
page 19, biffée, est reprise page 20. Pour $\mathcal{M} = L(H)$ et $\tau$ la
trace usuelle, $\mathfrak{a}$ est $\mathcal{S}_1$ et $\mathfrak{b}$ les
opérateurs compacts : tout ce numéro est la théorie des valeurs singulières,
écrite sans hypothèse de type.}

\emph{Proposition 1.} Un projecteur de $\mathfrak{b}$ est dans $\mathfrak{a}$.
Si $\mathcal{C} = C_0(M)$ est une sous-algèbre commutative involutive de
$\mathfrak{b}$, alors $K(M) \subset \mathfrak{a}$, où $K(M)$ désigne les
fonctions continues à support compact.

Pour la première assertion, si $E = \lim A_i$ en norme avec $A_i \in
\mathfrak{a}$, on a $E = \lim A_i^{*}A_i = \lim E A_i^{*} A_i E$, donc $E$ est
limite de $D_i \in \mathfrak{a}$ avec $0 \leq D_i \leq E$ ; dès que
$\|E - D_i\| \leq \frac12$ on a $E \leq 2D_i$, donc $E \in \mathfrak{a}$.
Pour la seconde, une $f \in K^{+}(M)$ est majorée par la fonction
caractéristique d'un compact $K$, elle-même majorée par une $g \in K^{+}(M)$ ;
l'opérateur $U_{1_K}$ de multiplication par $1_K$ est un projecteur, majoré
par $U_g \in \mathfrak{b}$, donc dans $\mathfrak{b}$, donc dans
$\mathfrak{a}$, et $U_f \leq \|f\|_\infty U_{1_K}$ est dans $\mathfrak{a}$.
Page 21, la fin de la démonstration établit que si $\mathcal{C}$ est
commutative \emph{maximale} dans $\mathfrak{b}$, le commutant de $\mathcal{C}$
dans $E\mathcal{M}E$, pour $E = U_{1_K}$, se réduit à $\mathcal{C}(K)$.

\emph{Corollaire 1.} Si $\mathcal{C} = C_0(M)$ est une sous-algèbre
commutative maximale de $\mathfrak{b}$, la trace induit une mesure positive
$\mu$ sur $M$ : $f \in C_0(M)$ est dans $\mathfrak{a}$ si et seulement si $f$
est $\mu$-intégrable, et alors $\tau(f) = \int f\,d\mu$.

\emph{Corollaire 2.} Si $A$ est un opérateur normal de $\mathfrak{b}$ et $f$
une fonction continue bornée sur $\mathbf{C}$, nulle au voisinage de $0$,
alors $f(A) \in \mathfrak{a}$.

Soit maintenant $A$ normal dans $\mathfrak{b}$. Il est contenu dans une
sous-algèbre commutative involutive maximale $C_0(M)$ de $\mathfrak{b}$,
comme fonction $f$ sur $M$ ; la mesure spectrale $\mu_f$ de $f$ pour la
mesure $\mu$ du corollaire 1 existe et ne dépend pas de $\mathcal{C}$, car
\[
  \langle \varphi, \mu_f\rangle = \int_M \varphi\circ f\,d\mu = \tau(\varphi(A))
  \qquad (\varphi \in K(\mathbf{C}^{*})) .
\]
On l'appelle \emph{mesure spectrale de $A$} et on la note $\mu_A$ ; sa masse
est finie si et seulement si $A \in \mathfrak{a}$, et alors
$\int z^n\,d\mu_A(z) = \tau(A^n)$ pour $n \geq 1$. Si $A \geq 0$, la fonction
décroissante $\varphi_A$ sur $\mathbf{R}^{*}_{+}$ de mesure spectrale $\mu_A$
est le \emph{réarrangement spectral} de $A$ ; on a $\|\varphi_A\|_\infty =
\|A\|$, et pour $A \in \mathfrak{b}$ quelconque on pose $\varphi_{|A|}$, puis
\[
  \Phi_A(t) = \int_0^t \varphi_A, \qquad
  \Psi_A(t) = \int_0^t \log\varphi_A, \qquad
  \Delta_A(t) = \exp\int_0^t \log\varphi_{|A|}(s)\,ds,
\]
et, si $\omega = \tau(1) < \infty$, $\Delta(A) = \Delta_A(\omega)$, avec
$0 \leq \Delta(A) < \infty$.\footnote{$\varphi_A(t) = \mu_t(A)$ de Fack et
Kosaki : $\mu_t(A) = \inf\{\|A(1-E)\| : E \in \mathcal{P},\ \tau(E) \leq t\}$,
et c'est la proposition 6 ci-dessous, trente ans plus tôt et pour les
seuls $A \in \mathfrak{b}$. Pour $\mathcal{M} = L(H)$, $\varphi_A$ est la
fonction en escalier $t \mapsto s_{\lceil t\rceil}(A)$.}

Les pages 23 à 25, en grande partie biffées et reprises, remarquent qu'un
point $t \in M$ de masse $\mu(\{t\}) > 0$ définit un projecteur minimal
$E = U_{1_{\{t\}}}$ de $\mathcal{M}$, que $\mathcal{M}$ se décompose alors en
une somme directe dont le facteur porté par $E$ est de type I avec $\tau$
proportionnelle à la trace usuelle, et que, inversement, une algèbre à trace
finie fidèle se plonge dans une algèbre continue en posant $\tau(f) =
\int_0^1 \tau(f(s))\,ds$ sur $L^\infty([0,1]) \mathbin{\bar\otimes}
\mathcal{M}$ ; le réarrangement d'un $A \in \mathcal{M}$ est le même dans
l'une et dans l'autre. C'est l'« astuce » que la page 34 nommera, et qui
permet de supposer la trace continue.\footnote{La prose de ces trois pages
est presque entièrement perdue ; on n'en retient que ce que les formules et
la page 34 confirment.}

\emph{Proposition 6.} Pour $A \in \mathfrak{b}^{+}$ et $t > 0$,
\[
  \varphi_A(t{-}0) = \inf_{\substack{E \in \mathcal{P}\\ \tau(E) < t}}
  \ \sup_{x \in (1-E)H,\ \|x\| \leq 1} (Ax, x)
  = \inf_{\substack{E \in \mathcal{P}\\ \tau(E) < t}} \|(1-E)A(1-E)\| .
\]
Prenant $E$ dans une sous-algèbre commutative maximale $C_0(M)$ contenant
$A$, les projecteurs sont des fonctions caractéristiques et la formule (4)
du numéro 1 donne l'inégalité $\geq$ ; et pour tout projecteur $E$ de
$\mathcal{M}$ avec $\tau(E) < t$, si $F$ est le projecteur spectral de $A$
sur $[\varphi_A(t{-}0), \infty[$, on a $\tau(F) \geq t > \tau(E)$, donc
$(1-E)H \cap FH \neq 0$ et le sup est $\geq \varphi_A(t{-}0)$.\footnote{Ce
second argument est celui de la page 34 (seconde mouture) : deux projecteurs
$E$, $F$ avec $\tau(E) < \tau(F)$ ne peuvent avoir $FH \cap (1-E)H = 0$, car
$F \ominus (F \wedge (1-E)) \sim E \ominus (E \wedge (1-F))$ donnerait
$F \prec E$. La page 26 ne donne que la première moitié.}

\emph{Corollaire 1.} Pour $A \in \mathfrak{b}$ quelconque,
$\varphi_{|A|}(t{-}0) = \inf_{\tau(E) < t} \|A(1-E)\|$.

\emph{Corollaire 2.} Pour $A \in \mathfrak{b}$, $B \in \mathcal{M}$ :
$\varphi_{|BA|} \leq \|B\|\,\varphi_{|A|}$ et $\varphi_{|AB|} \leq
\|B\|\,\varphi_{|A|}$.

\emph{Corollaire 3.} $\varphi_A$ est fonction croissante de
$A \in \mathfrak{b}^{+}$.

\emph{Corollaire 4.} Les fonctions $A \mapsto \varphi_{|A|}(t{-}0)$ sont
semi-continues supérieurement sur $\mathfrak{b}$ pour la topologie de la
norme.\footnote{La fin de l'énoncé est perdue au bord de la page ; la
topologie est celle que l'infimum de normes rend évidente.}

\emph{Proposition 7.} Pour $A \in \mathfrak{b}$, $\mu_{|A|} = \mu_{|A^{*}|}$,
c'est-à-dire $\varphi_{|A|} = \varphi_{|A^{*}|}$, et par suite $\Phi_{|A|} =
\Phi_{|A^{*}|}$, $\Psi_{|A|} = \Psi_{|A^{*}|}$, $\Delta_A = \Delta_{A^{*}}$.
Il suffit de voir $\mu_{A^{*}A} = \mu_{AA^{*}}$, et par continuité de le voir
pour $A \in \mathfrak{a}$ ; les deux mesures sont alors, hors de $0$, de
masse finie et à support dans $[0,\|A\|^2]$, donc déterminées par leurs
moments $\int x^p\,d\mu$, $p \geq 1$ entier, et $\tau((A^{*}A)^p) =
\tau((AA^{*})^p)$ résulte de $\tau(RS) = \tau(SR)$.

\pagerange{27}{32}
\section*{Le déterminant (pages 27 à 32)}

C'est le « N° 6 » de la page 27, le « N° 5 » du plan. Pour $A = \lambda 1 + B$
avec $B \in \mathfrak{b}$, on définit $\varphi_{|A|}$ par la formule du
corollaire 1 de la proposition 6, et
\[
  \Delta(A) = \exp\int_0^{\tau(1)} \log\varphi_{|A|}(s)\,ds \in [0, +\infty] .
\]
\footnote{La page définit d'abord $\Delta$ sur $\mathfrak{b}$, puis sur les
$\lambda 1 + B$, notés $\tilde{\mathfrak{b}}$ ; pour $\lambda \neq 0$ et
$\tau(1) = \infty$ la mesure spectrale de $|A|$ n'est plus de Radon sur
$\mathbf{R}^{*}_{+}$, et c'est la caractérisation par les projecteurs qui
donne un sens à $\varphi_{|A|}$. Ce choix est le nôtre.}
a) Si $\tau$ est finie, $\Delta(A)$ est défini pour tout $A \in \mathcal{M}$,
et $\Delta(\lambda A) = |\lambda|^{\tau(1)}\Delta(A)$ : c'est le déterminant
de Fuglede et Kadison.
b) Si $\tau$ est semi-finie et non finie, pour $A = \lambda 1 + B$ avec
$B \in \mathfrak{b}$ hermitien et $A \geq 0$ : $\Delta(A) = 0$ si
$|\lambda| < 1$, $\Delta(A) = +\infty$ si $|\lambda| > 1$, et si
$|\lambda| = 1$ on a $0 \leq \Delta(A) < \infty$ dès que $B \in
\mathfrak{a}$.\footnote{La page 28 écrit dans ce dernier cas « $\Delta(A) <
+\infty \Longleftrightarrow B \in \mathfrak{a}$ ». L'implication réciproque
est fausse : dans $L(H)$ avec $H$ de dimension infinie, $A = 1 - \frac12 E$
pour un projecteur $E$ de rang et de corang infinis a $\varphi_{|A|} \equiv
1$, donc $\Delta(A) = 1$, et $B = -\frac12 E$ n'est pas à trace. Elle est
vraie si $B \geq 0$, puisqu'alors $\varphi_{|1+B|} = 1 + \varphi_B$ et
$\log(1+\varphi_B)$ est sommable si et seulement si $\varphi_B$ l'est.}
Seules importent donc les valeurs de $\Delta$ sur $1 + \mathfrak{a}$.

\emph{Proposition 8.} Pour $A$, $B \in 1 + \mathfrak{a}$ :
\[
  \Delta(A) = \Delta(A^{*}), \qquad \Delta(AB) = \Delta(A)\,\Delta(B), \qquad
  \Delta(1) = 1 .
\]
Seule la seconde demande une démonstration. Par continuité on peut
supposer qu'il existe un projecteur $E \in \mathfrak{a}$ avec $A - 1$,
$B - 1 \in E\mathcal{M}E$, et l'on est ramené à l'algèbre
$\mathcal{M}_E = E\mathcal{M}E$, d'unité $E$, munie de la trace finie
$\tau_E$ induite ; pour $A \in \mathcal{M}_E$, on note $\Delta^{E}(A)$ son
déterminant relatif à $\mathcal{M}_E$, et l'on a $\Delta(1 + A) =
\Delta^{E}(1_E + A)$. Le cas d'une trace finie est celui de Fuglede et
Kadison.\footnote{Le manuscrit ne cite pas Fuglede et Kadison (1952), dont
la multiplicativité en trace finie est le théorème principal ; la page 29
commence une démonstration directe par $\Delta(A)^2 = \Delta(AA^{*})$, la
biffe, et la page 30 la remplace par les déterminants relatifs qui suivent.
Le passage de $1 + \mathfrak{a}$ au cas fini « par raison de continuité »
demande la semi-continuité de $A \mapsto \Delta(A)$, qui n'est pas écrite ;
elle résulte du corollaire 4 de la proposition 6.}

\emph{Déterminants relatifs.} Supposons la trace finie. Soient $X$, $Y$ deux
projecteurs équivalents de $\mathcal{M}$, $U \in \mathcal{M}$ une isométrie
partielle avec $U^{*}U = X$, $UU^{*} = Y$, et $A \in \mathcal{M}$ tel que
$AX(H) \subset Y(H)$, soit $\sigma(AX) \leq Y$ avec $\sigma$ le support à
gauche. On pose
\[
  \Delta^{X,Y}(A) = \Delta^{X}(U^{*}AX) = \Delta^{Y}(YAU^{*}),
\]
qui ne dépend pas du choix de $U$, deux tels $U$ différant par un unitaire
de $\mathcal{M}_X$. Si $Z$ est un projecteur équivalent à $Y$ et $B \in
\mathcal{M}$ avec $\sigma(BY) \leq Z$, on a
\[
  \Delta^{X,Z}(BA) = \Delta^{X,Y}(A)\,\Delta^{Y,Z}(B) .
\]

\emph{Proposition 9.} Soit $A \in \mathfrak{b}$ et $t = \tau(E_0)$ pour un
projecteur $E_0$. Alors
\[
  \Delta_A(t) = \sup_{\substack{E \simeq F\\ \tau(E) = t}} \Delta^{E,F}(FAE),
\]
et le sup est atteint ; si $A$ est normal, on peut se borner aux
$\Delta^{E}(EAE)$.\footnote{La page 30 écrit $\Delta_A(t) = \sup_{\tau(E)
\leq t} \Delta^{E}(EAE)$ pour $A \in \mathfrak{b}$ quelconque, et la
démontre « par diagonalisation ». Deux corrections. D'abord $\tau(E) = t$ et
non $\leq t$ : pour $A = \frac12 1$ en trace finie, $\Delta_A(t) = 2^{-t}$
décroît, tandis que le sup sur $\tau(E) \leq t$ vaut $1$ ; la page 31 elle-même
prend $t = \tau(E)$. Ensuite, la formule avec un seul projecteur exige $A$
normal : pour $A = e_{12}$ dans $M_2(\mathbf{C})$, $\Delta_A(1) = 1$ alors
que $|\langle Av, v\rangle| \leq \frac12$ pour tout vecteur unitaire $v$,
donc $\Delta^{E}(EAE) \leq \frac12$ pour tout projecteur de rang $1$. Avec
deux projecteurs équivalents on a l'égalité en général : si $A = U|A|$ et $E$
est le projecteur spectral de $|A|$ sur ses $t$ premières valeurs, $F =
UEU^{*}$ donne $\Delta^{E,F}(FAE) = \Delta^{E}(E|A|E) = \Delta_A(t)$. C'est
cette forme qu'exige la démonstration du théorème 1.}
L'inégalité $\leq$ est l'égalité qu'on vient de dire. L'inégalité $\geq$ est
le corollaire suivant, car $\varphi_{|EU^{*}AE|} \leq \varphi_{|A|}$ par le
corollaire 2 de la proposition 6, appliqué deux fois.

\emph{Corollaire.} Si $E \simeq F$, $\tau(E) = t$ et $A \in \mathcal{M}$,
alors $\Delta^{E,F}(FAE) \leq \Delta_A(t)$.

\emph{Théorème 1.} Pour $A$, $B \in \mathfrak{b}$ et $t \geq 0$,
\[
  \Delta_{AB}(t) \leq \Delta_A(t)\,\Delta_B(t) .
\]
Les deux membres étant fonctions semi-continues de $A$ et $B$, on peut
supposer la trace finie et $t = \tau(E)$. Par la proposition 9,
$\Delta_{AB}(t) = \sup \Delta^{E,G}(GABE)$ sur les couples $E \simeq G$ de
trace $t$ ; soit $F$ un projecteur équivalent à $E$ majorant $\sigma(BE)$,
$A' = GAF$, $B' = FBE$, de sorte que $GABE = A'B'$ et, par la
multiplicativité relative,
\[
  \Delta^{E,G}(GABE) = \Delta^{E,F}(FBE)\,\Delta^{F,G}(GAF)
  \leq \Delta_B(t)\,\Delta_A(t)
\]
par le corollaire.\footnote{C'est l'inégalité de Horn $\prod_{i \leq n}
s_i(AB) \leq \prod_{i \leq n} s_i(A)\,s_i(B)$ (1950), sous forme intégrée et
dans une algèbre de von Neumann semi-finie quelconque, où elle est due à
Fack (1982) et Fack--Kosaki (1986). Telle qu'elle est écrite page 31, la
démonstration passe par la proposition 9 à un seul projecteur appliquée au
produit $AB$, non normal ; on l'a réécrite avec les couples de projecteurs
équivalents, ce qui est exactement ce que la fin de la page 31 et la page 32
manipulent. Pour $t$ hors de $\tau(\mathcal{P})$ --- en type I, entre deux
entiers --- l'énoncé reste vrai par la caractérisation de Fack--Kosaki, mais
l'argument des pages ne le couvre pas.}

\pagerange{32}{33}
\section*{Les inégalités fondamentales (pages 32 et 33)}

Le théorème 1 s'écrit aussi $\Delta_{\varphi_{|AB|}} \leq
\Delta_{\varphi_{|A|}\varphi_{|B|}}$, puisque $\Delta_{fg} = \Delta_f\Delta_g$
pour deux fonctions décroissantes sur $\mathbf{R}^{*}_{+}$. Par le
corollaire 1 du numéro 3 :

\emph{Théorème 2.} Si $A$, $B \in \mathfrak{b}$ et si $W$ est une fonction de
Weyl croissante sur $K^\infty(]0,\tau(1)])$, alors
\[
  W(AB) \leq W(\varphi_{|A|}\,\varphi_{|B|}),
\]
où $W(AB)$ désigne $W(\varphi_{|AB|})$.

\emph{Corollaire 1.} Pour $t \geq 0$, $\int_0^t \varphi_{|AB|} \leq \int_0^t
\varphi_{|A|}\varphi_{|B|}$ ; en particulier
\[
  \int_0^\infty \varphi_{|AB|} \leq \int_0^\infty \varphi_{|A|}\,\varphi_{|B|},
  \qquad\text{d'où}\qquad
  |\tau(AB)| \leq \tau(|AB|) = \int \varphi_{|AB|} \leq \int \varphi_{|A|}\varphi_{|B|} .
\]
La plupart des inégalités importantes sont des conséquences immédiates de
la première.\footnote{C'est l'inégalité de trace de von Neumann (1937),
$|\operatorname{Tr} AB| \leq \sum s_i(A)s_i(B)$, dans une algèbre semi-finie ;
Fack--Kosaki, théorème 4.2.}

\emph{Corollaire 2} (Hölder pour les opérateurs). Pour $p$, $q$, $r > 0$ avec
$\frac1p + \frac1q = \frac1r$, et $\|A\|_r = \tau(|A|^r)^{1/r}$ :
$\|AB\|_r \leq \|A\|_p\,\|B\|_q$. Car $W(f) = \int |f|^r$ est une fonction de
Weyl croissante, d'où $\tau(|AB|^r) \leq \int (\varphi_{|A|}\varphi_{|B|})^r$,
et l'inégalité de Hölder classique majore le second membre.\footnote{Pour
$r < 1$ et $p$, $q$ quelconques, ce qui est remarquable : c'est le théorème
de la page 126, obtenu ici en trois lignes, et dans une algèbre de von Neumann
semi-finie.}

\emph{Corollaire 3.} Pour $A \in \mathfrak{b}$,
\[
  \Phi_{|A|}(t) = \sup\bigl\{ |\tau(BA)| : B \in \mathcal{M},\ \|B\| \leq 1,\   \tau(\sigma(B)) \leq t \bigr\} .
\]
L'inégalité $\geq$ : si $E = \sigma(B)$, $\varphi_{|B|} = \varphi_{|EB|} \leq
\varphi_E$, et par le corollaire 1, $|\tau(BA)| \leq \int \varphi_E
\varphi_{|A|} = \int_0^t \varphi_{|A|}$. L'inégalité $\leq$ : si $A = U|A|$,
il existe un projecteur $E$ commutant à $|A|$, de trace $t$, avec
$\tau(E|A|) = \Phi_{|A|}(t)$, et $B = EU^{*}$ convient.\footnote{C'est le
principe du maximum de Ky Fan (1951), $\sum_{i \leq n} s_i(A) = \sup
|\operatorname{Tr} BA|$ sur les $B$ de norme $\leq 1$ et de rang $\leq n$,
que la page 39 énonce dans $L(H)$. Le projecteur $E$ de trace exactement
$t$ existe si $t \in \tau(\mathcal{P})$ ; sinon on prend le sup.}

\emph{Corollaire 4.} Les fonctions $A \mapsto \Phi_{|A|}(t)$ sont
sous-additives sur $\mathfrak{b}$ : $\Phi_{|A+B|} \leq \Phi_{|A|} +
\Phi_{|B|}$. Elles sont des sup de semi-normes par le corollaire 3.

\pagerange{34}{35}
\section*{Seconde mouture (pages 34 et 35)}

Deux feuillets plus petits reprennent les propositions 6 et 7 et la
proposition 9, sous les numéros 1, 2 et une proposition sans numéro, avec
la trace supposée \emph{continue} --- c'est l'« astuce » de la page 34 : on
peut toujours s'y ramener par $L^\infty([0,1]) \mathbin{\bar\otimes}
\mathcal{M}$. On y trouve, sous cette hypothèse, $\varphi_A(t{+}0)$ à la place
de $\varphi_A(t{-}0)$ dans la proposition 6,\footnote{Avec la trace continue,
$\inf_{\tau(E) < t}$ et $\inf_{\tau(E) \leq t}$ diffèrent encore --- l'exemple
$A = \frac12 1 + \frac12 E$ le montre --- de sorte que le « $t + 0$ » de la
page 34, en face du signe strict, est la même inversion qu'à la page 7. On
garde la limite à gauche.} l'argument par les projecteurs équivalents
rapporté plus haut, les corollaires 1 à 3 et 5, et en marge un
\emph{corollaire 4} : $\varphi_{\alpha(A)} = \alpha(\varphi_A)$ pour $\alpha$
croissante, $\varphi_{|1+A|} \leq 1 + \varphi_{|A|}$, $\varphi_{1+|A|} = 1 +
\varphi_{|A|}$. Puis, avec la notation $\Delta_t(A)$ pour $\Delta_A(t)$ :
\[
  \Delta_t(A) = \sup_{\tau(E) = t} \Delta^{E}(EAE) \quad (A \geq 0), \qquad
  \Delta_t(1 + |A|) = \sup_{\substack{\tau(\sigma(B)) \leq t\\ \|B\| \leq 1}}
  \Delta^{E}\bigl(1_E + BAE\bigr) \quad (E = \sigma(B)),
\]
la seconde formule étant celle de la page 35 lue avec les corrections de la
proposition 9.\footnote{La page écrit $\operatorname{Tr}|B| < t$ sous le
sup, « $BAB$ » avec le dernier facteur noirci, et la première formule pour
$A$ quelconque avec $\operatorname{Tr} E < t$. On lit le dernier facteur
comme $E$ : l'inégalité $\Delta^{E}(1_E + BAE) \leq \Delta_t(1+|A|)$ résulte
de $\varphi_{|E + BAE|} \leq 1 + \varphi_{|BAE|} \leq 1 + \varphi_{|A|}$
sur $]0, \tau(E)[$, l'intégrande étant positive ; et si $A = U|A|$ et $E$ est
un projecteur spectral de $|A|$ de trace $t$, $B = EU^{*}$ donne
$BAE = E|A|E$ et l'égalité. Ici le sup sur $\tau(E) \leq t$ est légitime,
$\Delta_t(1+|A|)$ croissant avec $t$.} Elle est ce que la page 32 appelle « corollaire prop.\ 9 ».

\pagerange{37}{41}
\section*{Préliminaires sur les opérateurs compacts (pages 37 à 41)}

Quatre feuillets jaunis, encadrés au crayon rouge, qui rappellent ce que
tout le reste suppose ; le second porte « II ». Les pages 43 à 104 en sont
la rédaction. Si $E$ est un espace de Hilbert, son conjugué $\bar{E}$ l'est
et $E' = \bar{E}$. Les opérateurs compacts sont les limites en norme
d'opérateurs de rang fini. L'adjoint $u^{*} : F \to E$ de $u : E \to F$
vérifie $\|u^{*}u\| = \|u\|^2$. Un hermitien compact se diagonalise, $u =
\sum \lambda_i\, \bar{e}_i \otimes e_i$ avec $\lambda_i$ réels tendant vers
$0$ et $(e_i)$ orthonormal, et $u \geq 0$ si et seulement si tous les
$\lambda_i \geq 0$.

\emph{Décomposition polaire.} Pour $u : E \to F$ on pose $|u| =
(u^{*}u)^{1/2} \geq 0$ dans $E$ ; comme $\|\,|u|x\| = \|ux\|$, il existe une
isométrie partielle $U : E \to F$ et une isométrie partielle $V : F \to E$
telles que
\[
  u = U|u|, \qquad |u| = Vu, \qquad \|U\| \leq 1,\ \|V\| \leq 1,
\]
et en particulier $\|u\| = \|\,|u|\,\|$. Si $u$ est compact, $|u| = \sum
s_i\,\bar{e}_i \otimes e_i$ avec $s_i > 0$, et $u = \sum s_i\,\bar{e}_i
\otimes f_i$ où $f_i = Ue_i$ est orthonormal ; réciproquement une telle
écriture donne $|u|$, $u^{*} = \sum s_i\,\bar{f}_i \otimes e_i$ et
$|u^{*}| = \sum s_i\,\bar{f}_i \otimes f_i$, de sorte que $|u|$ et $|u^{*}|$
ont les mêmes valeurs propres.\footnote{C'est la décomposition de Schmidt ;
la fin de la page 37 est perdue dans une déchirure.}

\emph{Valeurs singulières.} Pour $u$ compact, $s_n(u) = \lambda_n(|u|)$, suite
décroissante positive ; $|\lambda_1(u)| \leq s_1(u) = \|u\|$ et $s_n(u) =
s_n(u^{*})$. Pour $h \geq 0$ compact, le principe du minimax donne
\[
  s_n(h) = \inf_{\dim E_{n-1} = n-1}\ \sup_{x \perp E_{n-1},\ \|x\| \leq 1} (hx, x),
  \qquad
  s_n(u) = \inf_{\dim E_{n-1} = n-1}\ \sup_{x \perp E_{n-1},\ \|x\| \leq 1} \|ux\|,
\]
d'où $s_n(uv) \leq \|u\|\,s_n(v)$, $s_n(uv) \leq \|v\|\,s_n(u)$, et $s_n(h)$
croît avec l'hermitien positif $h$. Puis les \emph{inégalités de Ky Fan},
que la page 143 démontre :
\[
  s_{m+n-1}(AB) \leq s_m(A)\,s_n(B), \qquad
  s_{m+n-1}(A+B) \leq s_m(A) + s_n(B) .
\]
\footnote{Ky Fan, 1951 ; la page 38 porte le nom dans une lecture incertaine,
et un point d'interrogation rouge en marge. Voir la page 143 pour la
démonstration et la correction d'un signe.}

\emph{Opérateurs à trace.} $u \in \mathcal{S}_1(E,F)$ si et seulement si
$\sum s_i(u) < \infty$, et alors $\|u\|_1 = \sum s_i(u)$ ; si $u$ est à trace,
$|u|$ l'est et $\|u\|_1 = \|\,|u|\,\|_1 = \operatorname{Tr}|u|$ ;
$|\operatorname{Tr} u| \leq \sum s_i(u)$. Pour $u$ compact,
\[
  \sum_{i=1}^n s_i(u) = \sup\bigl\{ |\operatorname{Tr} v_n u| : v_n \in L(F,E),\   \|v_n\| \leq 1,\ \operatorname{rang} v_n \leq n \bigr\},
\]
le sup étant atteint pour $v_n = \sum_{j \leq n} \bar{f}_j \otimes e_j$, et
pour un projecteur hermitien de rang $n$ si $u$ est hermitien positif. D'où,
pour $u$, $v$ compacts, $\sum_{i \leq n} s_i(u+v) \leq \sum_{i \leq n}
(s_i(u) + s_i(v))$, $\|u+v\|_1 \leq \|u\|_1 + \|v\|_1$, et $\sum s_i$ est une
norme sur $\mathcal{S}_1$. Enfin la définition des classes $\mathcal{S}_p(E,F)$
par $\|u\|_p = (\sum s_i(u)^p)^{1/p}$ et $S_p(u) = \sum s_i(u)^p$, avec les
cas $p = 1$ et $p = \infty$.

\emph{Dual des opérateurs compacts.} Le dual de $L_0(E,F)$ s'identifie
isométriquement à $\mathcal{S}_1(F,E)$ par $\langle u, v\rangle =
\operatorname{Tr} vu$. Que l'application $\mathcal{S}_1(F,E) \to
(L_0(E,F))'$ soit isométrique résulte de la formule précédente ; qu'elle
soit surjective : une forme continue $\Phi$ sur $L_0(E,F)$ définit par
$(x', y) \mapsto \Phi(x' \otimes y)$ une forme bilinéaire continue sur
$E' \times F$, donc un $v \in L(F,E)$ avec $\Phi(u) = \operatorname{Tr} vu$
pour $u$ de rang fini, et $|\operatorname{Tr} vu| \leq M\|u\|$ ; par
décomposition polaire on se ramène à $v = w \geq 0$ dans $L(E)$, puis $w$
est compact --- sinon $w$ serait inversible sur un sous-espace stable de
dimension infinie, ce qui contredit la borne --- et $w = \sum s_i\,
\bar{e}_i \otimes e_i$ avec $u = \sum_{i \leq n} \bar{e}_i \otimes e_i$ donne
$\sum_{i \leq n} s_i \leq M$.\footnote{Schatten (1946, 1950). La page 41
écrit $v = Ww$ avec $w = (vv^{*})^{1/2}$, ce qui est la décomposition polaire
à gauche ; et l'isométrie est vérifiée avec « $u = \sum_1^n \rho_i \bar{e}_i
\otimes f_i$ », un $\rho_i$ de trop que la transcription signale. Une
« méthode plus élémentaire » pour $\mathcal{S}_p$ et $\mathcal{S}_q$ est
esquissée en trois mots et ne se lit pas.}

\pagerange{43}{57}
\section*{Produits tensoriels hilbertiens (pages 43 à 57)}

Le numéro 2 d'un chapitre dont le numéro 1, « Rappels et notations.
Opérateurs de Fredholm », n'est qu'une ligne de plan. Il est, dit
l'interligne de la page 43, « vraiment classique », et on le résume.

\emph{L'espace conjugué.} Pour $E$ un espace de Banach complexe, $\bar{E}$
est le même groupe additif avec la loi $(\lambda, x) \mapsto \bar\lambda x$ ;
l'identité $E \to \bar{E}$ est antilinéaire, et $\bar{\bar{E}} = E$. Une
application sesquilinéaire $E \times F \to G$ est une application bilinéaire
$E \times \bar{F} \to G$, donc une application linéaire $E \otimes \bar{F}
\to G$, et $L(E, \bar{F}) = \overline{L(E,F)}$. On a canoniquement
$\overline{E \otimes F} = \bar{E} \otimes \bar{F}$. Si $E$ est préhilbertien,
$(\bar{x}, \bar{y}) = (y, x)$ fait de $\bar{E}$ un espace préhilbertien de
même norme, et si $E$ est de Hilbert, $y \mapsto (\,\cdot\,, y)$ est une
bijection antilinéaire isométrique de $E$ sur $E'$, soit
\[
  E' = \bar{E} \qquad (E \text{ espace de Hilbert}).
\]

\emph{Le produit tensoriel hilbertien.} Sur $E \otimes F$, la forme
sesquilinéaire $(x_1 \otimes y_1, x_2 \otimes y_2) = (x_1, x_2)(y_1, y_2)$ est
définie positive : sur $E_0 \otimes F_0$ pour $E_0$, $F_0$ de dimension finie,
les $e_i \otimes f_j$ formés de bases orthonormales sont orthonormaux, et
$E \otimes F$ est réunion de ces sous-espaces. Le complété $E
\mathbin{\hat\otimes}_2 F$ est le \emph{produit tensoriel hilbertien}.

\emph{Proposition 1.} Si $(e_i)$ et $(f_j)$ sont des bases hilbertiennes de
$E$ et $F$, $(e_i \otimes f_j)$ est une base hilbertienne de $E
\mathbin{\hat\otimes}_2 F$. De plus $\|x \otimes y\|_2 = \|x\|\,\|y\|$, et
$\overline{E \mathbin{\hat\otimes}_2 F} = \bar{E} \mathbin{\hat\otimes}_2 \bar{F}$.

\emph{Théorème 1.} Pour $\mu$ une mesure positive sur un espace localement
compact et $E$ un espace de Hilbert, $L^2(\mu) \mathbin{\hat\otimes}_2 E =
L^2_E(\mu)$, l'espace des fonctions de carré intégrable à valeurs dans $E$,
avec $(f, g) = \int (f(t), g(t))\,d\mu(t)$ : l'application $\varphi \otimes a
\mapsto \varphi\,a$ conserve les produits scalaires et a une image dense.
\emph{Corollaire 1} : $\ell^2(I) \mathbin{\hat\otimes}_2 E = \ell^2_E(I)$.
\emph{Corollaire 2} : $L^2(\mu) \mathbin{\hat\otimes}_2 L^2(\nu) = L^2(\mu
\otimes \nu)$, puisque $L^2(\mu \otimes \nu) = L^2_{L^2(\nu)}(\mu)$.

\emph{Opérateurs de Hilbert--Schmidt.} L'application $(\bar{a} \otimes b)\,x =
(x, a)\,b$ identifie $\bar{E} \otimes F$ aux opérateurs de rang fini de $E$
dans $F$, et pour $u \in \bar{E} \otimes F$ on a $(ux, y) = (u, \bar{x}
\otimes y)$, donc $|(ux, y)| \leq \|u\|_2\,\|x\|\,\|y\|$ : l'application
$\bar{E} \mathbin{\hat\otimes}_2 F \to L(E,F)$ est continue de norme $\leq 1$,
et injective, car dans des bases hilbertiennes elle envoie une matrice de
carré sommable sur l'opérateur qu'elle définit.

\emph{Définition 2.} Un \emph{opérateur de Hilbert--Schmidt} de $E$ dans $F$
est un opérateur défini par un élément de $\bar{E} \mathbin{\hat\otimes}_2 F$ ;
leur espace, muni de cette structure hilbertienne, est noté
$\mathcal{S}_2(E,F)$, et $\|u\|_2 = (\sum_{i,j} |u_{ij}|^2)^{1/2}$ dans des
bases hilbertiennes. On a $\|u\| \leq \|u\|_2$, et ces opérateurs sont
compacts.

\emph{Théorème 2.} Si $u \in \mathcal{S}_2(E,F)$ et $A \in L(F,G)$, alors
$Au \in \mathcal{S}_2(E,G)$ et $\|Au\|_2 \leq \|A\|\,\|u\|_2$, et de même
pour $uA$ ; $u \in \mathcal{S}_2(E,F)$ si et seulement si $u^{*} \in
\mathcal{S}_2(F,E)$, et $\|u\|_2 = \|u^{*}\|_2$.

\emph{Théorème 3.} $u$ est de Hilbert--Schmidt si et seulement si $\sum
s_i(u)^2 < \infty$, et alors $\|u\|_2 = \|u^{*}\|_2 = (\sum s_i(u)^2)^{1/2}$ ;
en particulier $u \in \mathcal{S}_2$ si et seulement si $u^{*}u \in
\mathcal{S}_1$, et $\|u\|_2 = \sqrt{\operatorname{Tr} u^{*}u} =
\sqrt{\|u^{*}u\|_1}$. La démonstration passe par $u = Uh$, $h = Vu$ avec
$h = |u|$ et $U$, $V$ des isométries partielles, puis diagonalise $h$.

\emph{Théorème 3 bis.} Pour $u \in L(E,F)$ : $u$ est à trace si et seulement
si $|u| = \sqrt{u^{*}u}$ l'est, si et seulement si $\sum s_i(u) < \infty$, et
alors
\[
  \|u\|_1 = \|\,|u|\,\|_1 = \sum s_i(u) = \operatorname{Tr}\sqrt{u^{*}u}
  = \operatorname{Tr}\sqrt{uu^{*}} .
\]
\footnote{La page 54 numérote ce théorème « 3 » comme le précédent ; on le
distingue. Sa démonstration suppose connue une norme $\|u\|_1$ sur les
opérateurs à trace pour laquelle $\|uA\|_1 \leq \|u\|_1\|A\|$ et
$|\operatorname{Tr} u| \leq \|u\|_1$ : c'est la norme nucléaire du numéro 1
absent, et $\sum_{i \leq n} s_i = \operatorname{Tr} p_n h \leq \|p_n h\|_1
\leq \|h\|_1$ en découle. Que $\sqrt{u^{*}u}$ et $\sqrt{uu^{*}}$ aient les
mêmes valeurs propres est vérifié sur $u = \sum s_i\,\bar{a}_i \otimes b_i$.}
\emph{Corollaire 1} : les opérateurs à trace de $E$ dans $F$ sont exactement
les $u = \sum s_i\,\bar{a}_i \otimes b_i$ avec $(a_i)$, $(b_i)$ orthonormaux
et $(s_i)$ positive sommable, et alors $\|u\|_1 = \sum s_i$.
\emph{Corollaire 2} : $\|u\|_1 = \sup_{\|A\| \leq 1} |\operatorname{Tr} Au|$
sur $A \in L(F,E)$.

\emph{Théorème 4.} Si $u \in \mathcal{S}_2(E,F)$ et $v \in \mathcal{S}_2(F,G)$,
alors $vu \in \mathcal{S}_1(E,G)$ et $\|vu\|_1 \leq \|v\|_2\,\|u\|_2$ ; et pour
$u$, $v \in \mathcal{S}_2(E,F)$,
\[
  (u, v) = \operatorname{Tr} uv^{*} = \operatorname{Tr} v^{*}u .
\]
Par continuité on se ramène aux rangs finis ; alors $|\operatorname{Tr} Avu|
\leq \|Av\|_2\|u\|_2 \leq \|A\|\,\|v\|_2\|u\|_2$ pour $\|A\| \leq 1$ donne la
première formule par le corollaire 2, et la seconde se vérifie sur
$u = \bar{a} \otimes b$, $v = \bar{a}' \otimes b'$, où $v^{*}u = (b, b')\,\bar{a}
\otimes a'$ et $\operatorname{Tr} v^{*}u = (a', a)(b, b') = (u, v)$.
\emph{Corollaire 1} : $u \in L(E,F)$ est de Hilbert--Schmidt si et seulement
si $vu$ est à trace pour tout $v \in \mathcal{S}_2(F,E)$, et alors $\|u\|_2 =
\sup_{\|v\|_2 \leq 1} |\operatorname{Tr} vu|$ ; la réciproque passe par le
théorème du graphe fermé, qui rend $v \mapsto vu$ continue de
$\mathcal{S}_2(F,E)$ dans $\mathcal{S}_1(E)$, donc $w \mapsto
\operatorname{Tr} w^{*}u$ continue sur $\bar{E} \mathbin{\hat\otimes}_2 F$,
donc représentée par un $u_0$, qui coïncide avec $u$.
\emph{Corollaire 2} : $u \in L(E)$ est à trace si et seulement si $u = vw$
avec $v$, $w \in \mathcal{S}_2(E)$, et l'on peut prendre $\|v\|_2 = \|w\|_2 =
\sqrt{\|u\|_1}$, en écrivant $u = U W^2$ avec $W = \sqrt{|u|}$ et $v = UW$.

\pagerange{58}{69}
\section*{Produits extérieurs hilbertiens (pages 58 à 69)}

Le numéro 3, « Applications : des inégalités diverses ». Le produit tensoriel
hilbertien de $n$ espaces $E_1, \ldots, E_n$ est défini par récurrence, avec
$(a_1 \otimes \cdots \otimes a_n, b_1 \otimes \cdots \otimes b_n) = \prod (a_i,
b_i)$ et pour base hilbertienne les $e^1_{\alpha_1} \otimes \cdots \otimes
e^n_{\alpha_n}$ ; il est associatif, et pour $u_i \in L(E_i, F_i)$ on a
$\|u_1 \otimes \cdots \otimes u_n\| \leq \prod \|u_i\|$, qu'on ramène par
associativité au cas de deux facteurs et d'un seul $u_i \neq 1$, où c'est le
théorème 2 du numéro 2 sous l'identification $E_1 \mathbin{\hat\otimes}_2 E_2
= \mathcal{S}_2(\bar{E}_1, E_2)$.

Sur $\bigotimes^n_2 E$, l'antisymétriseur
\[
  a_n = \frac{1}{n!} \sum_{\sigma \in \mathfrak{S}_n} \varepsilon_\sigma\,\sigma
\]
est un projecteur hermitien, l'adjoint de $\sigma$ étant $\sigma^{-1}$.

\emph{Définition.} La \emph{puissance extérieure hilbertienne} $\Lambda^n_2 E$
est l'image de $a_n$, munie du produit scalaire de $\bigotimes^n_2 E$
multiplié par $n!$. La puissance extérieure algébrique $\Lambda^n E$ s'y
plonge comme sous-espace dense, et sur les multivecteurs décomposables
\[
  (a_1 \wedge \cdots \wedge a_n,\ b_1 \wedge \cdots \wedge b_n)
  = n!\,(a_n A, a_n B) = \sum_\sigma \varepsilon_\sigma \prod_i (a_i, b_{\sigma i})
  = \det\bigl((a_i, b_j)\bigr) .
\]
\footnote{Avec $A = a_1 \otimes \cdots \otimes a_n$, $B = b_1 \otimes \cdots
\otimes b_n$. Le facteur $n!$ est le choix qui fait de $\{e_J\}$ une base
orthonormale ; sans lui, $\|e_{i_1} \wedge \cdots \wedge e_{i_n}\| =
(n!)^{-1/2}$. C'est la convention de l'espace de Fock antisymétrique.}

\emph{Proposition.} Il existe sur $\Lambda E \times \Lambda E$ une unique
forme hermitienne définie positive satisfaisant à cette formule, et son
complété est $\Lambda_2 E$. Si $(e_i)_{i \in I}$ est une base hilbertienne de
$E$, $I$ totalement ordonné, les $e_J = e_{i_1} \wedge \cdots \wedge e_{i_n}$,
$J = \{i_1 < \cdots < i_n\}$, forment une base hilbertienne de $\Lambda_2 E$.
En particulier
\[
  \|a_1 \wedge \cdots \wedge a_n\| = \det\bigl((a_i, a_j)\bigr)^{1/2}
  \leq \|a_1\| \cdots \|a_n\|,
\]
l'inégalité venant, si $a_i = \sum_j a_{ij} e_j$ et $A = (a_{ij})$, de
$\det((a_i, a_j)) = \det AA^{*} = |\det A|^2$ et du choix d'une base où $A$
est triangulaire ; ce qui donne l'\emph{inégalité de Hadamard}
$|\det A| \leq \prod_i \|a_i\|$, et le \emph{corollaire 1} :
$|\det((a_i, b_j))| \leq \prod_i \|a_i\|\,\|b_i\|$, par Cauchy--Schwarz.
Ce corollaire montre que la constante $K_n$ de la théorie générale vaut au
plus $1$ dans l'espace de Hilbert.\footnote{Renvoi à « Chap.\ 1, § 2, n° 2,
prop.\ 1 » : dans \emph{La théorie de Fredholm}, la constante qui borne
$|\det((\langle x_i, x'_j\rangle))|$ par $\prod \|x_i\|\|x'_j\|$ dans un
espace de Banach quelconque. Ces pages ne la définissent pas.}

\emph{Puissances extérieures d'opérateurs.} Pour $u_1, \ldots, u_n \in L(E,F)$,
le produit extérieur $u_1 \wedge \cdots \wedge u_n : \Lambda^n E \to \Lambda^n
F$ s'écrit $n!\,\varphi_n\,(u_1 \otimes \cdots \otimes u_n)\,a_n$, avec
$\varphi_n$ la projection de $\bigotimes^n F$ sur $\Lambda^n F$ ; comme
$a_n/\sqrt{n!}$ est une isométrie de $\Lambda_2 E$ dans $\bigotimes^n_2 E$ et
$\sqrt{n!}\,\varphi_n$ de norme $\leq 1$, on a
\[
  \|u_1 \wedge \cdots \wedge u_n\| \leq \|u_1\| \cdots \|u_n\|,
\]
et $u_1 \wedge \cdots \wedge u_n$ se prolonge à $\Lambda^n_2 E \to \Lambda^n_2
F$. Il est symétrique en les $u_i$, et l'on note $\Lambda^n u = u \wedge
\cdots \wedge u$, $\Lambda u$ son action sur $\Lambda_2 E$ tout entier.
Alors
\begin{gather*}
  (\Lambda v)(\Lambda u) = \Lambda(vu), \qquad
  (u_1 \wedge \cdots \wedge u_n)^{*} = u_1^{*} \wedge \cdots \wedge u_n^{*}, \\
  (\Lambda u)^{*} = \Lambda(u^{*}), \qquad
  (\Lambda h)^\alpha = \Lambda(h^\alpha)\quad (h \geq 0,\ \alpha > 0),
\end{gather*}
la seconde par $((\Lambda u)\,a_1 \wedge \cdots \wedge a_n, b_1 \wedge \cdots
\wedge b_n) = \det((ua_i, b_j)) = \det((a_i, u^{*}b_j))$, la troisième pour
$\alpha = 1/n$ par unicité de la racine positive, puis pour $\alpha$
rationnel, puis par continuité de $\alpha \mapsto h^\alpha$.

\emph{Proposition 1.} Résumé : $u_1 \wedge \cdots \wedge u_n$ est continue de
$\Lambda^n_2 E$ dans $\Lambda^n_2 F$, de norme $\leq \prod \|u_i\|$ ; $u
\mapsto \Lambda u$ de $L(E)$ dans $L(\Lambda_2 E)$ est multiplicative,
commute à l'involution, envoie les positifs sur des positifs et commute aux
puissances $h \mapsto h^\alpha$.

\emph{Proposition 2.} Soit $u$ compact dans $E$ et $(\lambda_i)$ la suite de
ses valeurs propres non nulles, avec multiplicités algébriques. Les valeurs
propres non nulles de $\Lambda^n u$ sur $\Lambda^n_2 E$, avec leurs
multiplicités, sont les produits $\lambda_{i_1} \cdots \lambda_{i_n}$,
$i_1 < \cdots < i_n$. Par continuité de $u \mapsto \Lambda u$ on se ramène
aux $u$ de rang fini ; alors $E = E_0 \oplus E_1$ avec $E_0$ de dimension
finie stable par $u$ et $u(E_1) = 0$, $\Lambda^n_2 E = \bigoplus_{p+q=n}
\Lambda^p_2 E_0 \otimes \Lambda^q_2 E_1$, $\Lambda^n u$ s'annule sur les
termes $q > 0$ et agit sur $\Lambda^n E_0$, où c'est l'algèbre linéaire
d'une matrice triangulaire.\footnote{Le passage à la limite demande que les
valeurs propres non nulles dépendent continûment de $u$ avec leurs
multiplicités, ce qui est vrai pour les opérateurs compacts (continuité des
projecteurs de Riesz) ; la page 66 le dit par « base hilbertienne
convenable » dans un bloc biffé.}

\emph{Corollaire 1.} Pour $u$ compact :
\begin{gather*}
  \lambda_1(u) \cdots \lambda_n(u) = \lambda_1(\Lambda^n u), \qquad
  s_1(u) \cdots s_n(u) = \|\Lambda^n u\| = s_1(\Lambda^n u), \\
  |\lambda_1(u) \cdots \lambda_n(u)| \leq s_1(u) \cdots s_n(u),
\end{gather*}
la seconde parce que $|\Lambda^n u| = ((\Lambda^n u)^{*}\Lambda^n u)^{1/2} =
\Lambda^n |u|$ a pour plus grande valeur propre $s_1 \cdots s_n$, et la
troisième par $|\lambda_1(w)| \leq s_1(w)$ appliqué à $w = \Lambda^n u$.

\emph{Corollaire 3.} Pour $u \in L_0(E,F)$, $v \in L_0(F,G)$ et tout $n$,
\[
  s_1(vu) \cdots s_n(vu) \leq \bigl(s_1(u)\,s_1(v)\bigr) \cdots \bigl(s_n(u)\,s_n(v)\bigr),
\]
car le premier membre est $\|\Lambda^n(vu)\| = \|(\Lambda^n v)(\Lambda^n u)\|
\leq \|\Lambda^n v\|\,\|\Lambda^n u\|$. L'idée est due à H.\ Weyl, dit la
page 69.\footnote{Les inégalités du corollaire 1 sont celles de Weyl (1949) ;
celle du corollaire 3 est de Horn (1950), qui l'obtient par la même méthode
des puissances extérieures. Il n'y a pas de corollaire 2 visible, et la
page 69 écrit $\rho_i(uv)$ là où $w = vu$ ; on a suivi $w$.}

\pagerange{70}{73}
\section*{Puissances extérieures d'opérateurs à trace, mise au net (pages 70 à 73)}

Une rédaction sur feuilles blanches, numérotée à partir de (1) ; les pages
76 à 79 en sont le brouillon, et l'on y renvoie pour les démonstrations
qu'elle omet.

\emph{Proposition 1.} Pour $u_1, \ldots, u_n \in \mathcal{S}_1(E,F)$,
$u_1 \wedge \cdots \wedge u_n$ est à trace et
\[
  \|u_1 \wedge \cdots \wedge u_n\|_1 \leq \frac{1}{n!}\,\|u_1\|_1 \cdots \|u_n\|_1 .
\]
Par linéarité il suffit de prendre $u_i = \bar{b}_i \otimes a_i$ avec
$\|a_i\|, \|b_i\| \leq 1$ ; alors $u_1 \wedge \cdots \wedge u_n = \frac{1}{n!}
\,\overline{(b_1 \wedge \cdots \wedge b_n)} \otimes (a_1 \wedge \cdots \wedge
a_n)$, de norme trace $\leq \frac{1}{n!}\|b_1 \wedge \cdots \wedge b_n\|\,
\|a_1 \wedge \cdots \wedge a_n\| \leq \frac{1}{n!}$ par Hadamard.

\emph{Corollaire 1.} $\alpha_n(u_1, \ldots, u_n) = \operatorname{Tr}(u_1 \wedge
\cdots \wedge u_n)$, et $\alpha_n(u) = \operatorname{Tr}\Lambda^n u$ : les
coefficients du déterminant de Fredholm sont les traces des puissances
extérieures. \emph{Corollaire 2} : $|\alpha_n(u_1, \ldots, u_n)| \leq
\frac{1}{n!}\prod \|u_i\|_1$ et $|\alpha_n(u)| \leq \frac{1}{n!}\|u\|_1^n$.
\emph{Corollaire 3} :
\[
  |\alpha_n(u)| \leq \alpha_n(|u|) = \|\Lambda^n u\|_1 \leq \frac{1}{n!}\,\|u\|_1^n,
\]
car $\alpha_n(|u|) = \operatorname{Tr}\Lambda^n|u| = \operatorname{Tr}|\Lambda^n u|$.
\emph{Corollaire 4} : $|\det(1+u)| \leq \det(1+|u|)$, et plus généralement
$|\det(1+zu)| \leq \det(1+|z|\,|u|)$.

\emph{Théorème 1.} Pour $u \in \mathcal{S}_1(E)$, la suite $(\lambda_i(u))$ est
sommable et $\det(1+zu)$ est une fonction entière de genre zéro :
\[
  \det(1+zu) = \prod_i \bigl(1 + z\lambda_i(u)\bigr) ;
\]
de façon équivalente, $\alpha_n(u) = \sum_{i_1 < \cdots < i_n} \lambda_{i_1}
\cdots \lambda_{i_n}$, et en particulier
\[
  \operatorname{Tr} u = \sum_i \lambda_i(u),
\]
d'où l'on retrouve la formule pour $\alpha_n$ en l'appliquant à
$\Lambda^n u$. De même $\alpha_n(|u|) = \|\Lambda^n u\|_1 = \sum_{i_1 <
\cdots < i_n} s_{i_1} \cdots s_{i_n}$.\footnote{C'est le théorème de la
trace, publié par Lidskii en 1959 ; Grothendieck l'établit dans
\emph{La théorie de Fredholm} (1956) pour les espaces de Hilbert, et la
démonstration de ces pages, biffée page 79 et complétée par la proposition
de la page 110, est la sienne : les coefficients $\alpha_n(u)$ sont majorés
par ceux de $\det(1+z|u|)$, fonction de genre zéro, les zéros $-1/\lambda_i$
vérifient $\sum |\lambda_i| < \infty$ par Weyl, donc $\det(1+zu)$ est de
genre zéro et vaut $\prod(1+z\lambda_i)$. L'ordre des pages cache que la
démonstration existe ; la remarque des pages 106 à 111 la rend visible.}

\emph{Proposition 2.} Pour $u$, $v \in \mathcal{S}_1$ et $p$, $q \geq 0$, en
notant $(\Lambda^p u) \wedge (\Lambda^q v)$ le produit extérieur de $p$
copies de $u$ et $q$ de $v$,
\[
  \|(\Lambda^p u) \wedge (\Lambda^q v)\|_1 \leq \frac{p!\,q!}{(p+q)!}\,
  \|\Lambda^p u\|_1\,\|\Lambda^q v\|_1
  = \frac{p!\,q!}{(p+q)!}\,\alpha_p(|u|)\,\alpha_q(|v|) .
\]
Avec $u = \sum \rho_i\,\bar{b}_i \otimes a_i$ et $v = \sum \sigma_j\,\bar{d}_j
\otimes c_j$ (décompositions de Schmidt), la multilinéarité et
l'antisymétrie donnent
\[
  (\Lambda^p u) \wedge (\Lambda^q v) = \frac{p!\,q!}{(p+q)!}
  \sum_{\substack{i_1 < \cdots < i_p\\ j_1 < \cdots < j_q}}
  \rho_{i_1} \cdots \rho_{i_p}\,\sigma_{j_1} \cdots \sigma_{j_q}\;
  \overline{(b_{i_1} \wedge \cdots \wedge d_{j_q})} \otimes
  (a_{i_1} \wedge \cdots \wedge c_{j_q}),
\]
et chaque terme est de norme trace $\leq 1$ par Hadamard.\footnote{La page
71 écrit le facteur $\frac{1}{(p+q)!}$ devant la somme et fait apparaître
$p!\,q!$ à la ligne suivante ; le compte est le même.}

\emph{Corollaire 2.} Pour tout $n$,
\[
  \alpha_n(|u+v|) \leq \sum_{p=0}^{n} \alpha_p(|u|)\,\alpha_{n-p}(|v|) ,
\]
car $\alpha_n(|u+v|) = \|\Lambda^n(u+v)\|_1$ et $\Lambda^n(u+v) =
\sum_p \binom{n}{p}\,(\Lambda^p u) \wedge (\Lambda^{n-p} v)$.

\emph{Corollaire 1.} Pour $r \geq 0$,
\begin{gather*}
  \det(1 + r|u+v|) \leq \det(1 + r|u|)\,\det(1 + r|v|), \qquad\text{i.e.} \\
  \prod_i (1 + r s_i(u+v)) \leq \prod_i (1 + r s_i(u)) \prod_i (1 + r s_i(v)) .
\end{gather*}
\footnote{La page 71 énonce ce corollaire avant le corollaire 2 qui le
démontre coefficient par coefficient. L'inégalité est celle que la
littérature attribue à Seiler et Simon (1975), qui la déduisent du théorème
de Rotfel'd (1967) ; elle est ici, avec sa démonstration par les
puissances extérieures, dans un carton daté 1953. Rien n'indique que ces
pages aient circulé, et cette lecture n'affirme rien sur ce qui a été
publié.}

\emph{Théorème 2.} Pour $u$, $v$ compacts, $r \geq 0$ et tout $n$,
\[
  \prod_{i=1}^{n} \bigl(1 + r s_i(u+v)\bigr) \leq
  \prod_{i=1}^{n} \bigl(1 + r s_i(u)\bigr)\prod_{i=1}^{n} \bigl(1 + r s_i(v)\bigr),
\]
et plus précisément, avec $\sigma_n^m(w) = \sum_{1 \leq i_1 < \cdots < i_m
\leq n} s_{i_1}(w) \cdots s_{i_m}(w)$, pour $0 \leq m \leq n$ :
\[
  \sigma_n^m(u+v) \leq \sum_{p=0}^{m} \sigma_n^p(u)\,\sigma_n^{m-p}(v) .
\]

\emph{Proposition 3.} Pour $u$ compact et $m \leq n$,
\[
  \sigma_n^m(u) = \sup\bigl\{ |\alpha_m(v_n u)| : v_n \in L(F,E),\ \|v_n\| \leq 1,\   \operatorname{rang} v_n \leq n \bigr\} = \sup_{v_n} \alpha_m(|v_n u|) .
\]
Pour $\leq$ : si $u = \sum s_i\,\bar{e}_i \otimes f_i$, $v_n = \sum_{i \leq n}
\bar{f}_i \otimes e_i$ donne $v_n u = \sum_{i \leq n} s_i\,\bar{e}_i \otimes
e_i$ et $\alpha_m(v_n u) = \sigma_n^m(u)$. Pour $\geq$ : $s_i(v_n u) \leq
s_i(u)$ et $v_n u$ est de rang $\leq n$, donc $|\alpha_m(v_n u)| \leq
\alpha_m(|v_n u|) = \sigma_n^m(v_n u) \leq \sigma_n^m(u)$.
\emph{Proposition 4} : $\prod_{i \leq n}(1 + s_i(u)) = \sup_{v_n}
\det(1 + v_n u)$ sur les mêmes $v_n$.

La démonstration du théorème 2 : pour $w \in L(F,E)$ de norme $\leq 1$ et de
rang $\leq n$, le corollaire 2 donne $|\alpha_m(w(u+v))| \leq
\sum_p \alpha_p(|wu|)\,\alpha_{m-p}(|wv|) \leq \sum_p \sigma_n^p(u)\,
\sigma_n^{m-p}(v)$, et l'on prend le sup en $w$ par la proposition 3. La page
73 juge ce système d'inégalités « les plus fortes possibles » majorant les
$s_i(u+v)$ en fonction des $s_i(u)$ et $s_i(v)$.\footnote{Elles sont, pour
$f(s) = \log(1 + rs)$, l'inégalité $\sum_{i \leq n} f(s_i(u+v)) \leq
\sum_{i \leq n} f(s_i(u)) + f(s_i(v))$ que Rotfel'd (1967) et Thompson (1976)
établissent pour toute $f$ concave croissante nulle en $0$ ; les pages 74
et 75 vont exactement dans cette direction.}

\pagerange{74}{75}
\section*{Sous-additivité des fonctions des valeurs singulières (pages 74 et 75)}

Deux feuilles brunes d'une autre plume tirent du théorème 2 ce qu'il donne
pour d'autres fonctions que $\log(1+rs)$. Écrivons-le
\[
  \sum_{i \leq n} \log(1 + r s_i(u+v)) \leq \sum_{i \leq n} \log(1 + r s_i(u))
  + \sum_{i \leq n} \log(1 + r s_i(v)) \qquad (r \geq 0) .
\]
Si $\varphi$ est une fonction positive décroissante sur $\mathbf{R}^{*}_{+}$,
et si $A \leq B$ sont deux fonctions croissantes nulles en $0$, alors
$\int \varphi(r)\,dA(r) \leq \int \varphi(r)\,dB(r)$, par intégration par
parties : $\int \varphi\,dA = [\varphi A] + \int A\,(-d\varphi)$ et les deux
termes croissent avec $A$. Appliquée à $A(r) = \sum_{i \leq n} \log(1 +
r s_i(u+v))$ et $B$ le second membre, cela donne, en posant
\[
  f_\varphi(\rho) = \int_0^{\infty} \frac{\varphi(r)}{\frac{1}{\rho} + r}\,dr
  = \int_0^{\infty} \varphi(s/\rho)\,\frac{ds}{1+s}
  \quad\text{et}\quad
  S^n_\varphi(w) = \sum_{i \leq n} f_\varphi(s_i(w)),
\]
l'inégalité $S^n_\varphi(u+v) \leq S^n_\varphi(u) + S^n_\varphi(v)$, puisque
$dA(r) = \sum_i \frac{s_i}{1 + r s_i}\,dr$.\footnote{La page 74 écrit
$A'(r) = \sum r/(1 + r\rho_i)$, coquille pour $\sum \rho_i/(1 + r\rho_i)$,
que la transcription note.} Le cas le plus important est $\varphi(r) =
r^{-p}$, $0 < p < 1$, où $f_\varphi(\rho) = \rho^p \int_0^\infty
\frac{s^{-p}}{1+s}\,ds = \frac{\pi}{\sin \pi p}\,\rho^p$ ; d'où, avec
$S^{(n)}_p(u) = \sum_{i \leq n} s_i(u)^p$,
\[
  S^{(n)}_p(u+v) \leq S^{(n)}_p(u) + S^{(n)}_p(v), \qquad
  S_p(u+v) \leq S_p(u) + S_p(v) \qquad (0 < p < 1) .
\]
\emph{Remarque} : pour $p \geq 1$ ce n'est pas la bonne inégalité, qui est
celle de Weyl (la norme). \emph{Question} : trouver toutes les fonctions $f$
d'une variable positive telles que $f^{(n)}(u) = \sum_{i \leq n} f(s_i(u))$
soit sous-additive pour tout $n$ ; il suffit que $f^{(\infty)}$ le
soit.\footnote{Toute fonction concave croissante nulle en $0$ convient :
c'est le théorème de Rotfel'd (1967) et Thompson (1976), dont les
$f_\varphi$ de la page 74 --- mélanges des noyaux concaves $\rho \mapsto
\rho/(1 + r\rho)$ --- sont un cas particulier. La sous-additivité de $f$
elle-même est nécessaire ($n = 1$, $u$ et $v$ multiples d'un même projecteur
de rang un). La réponse de la page 75, « arbitraire », est illisible au-delà
de ce mot.} La sous-additivité de $S_p$ pour $p \leq
1$ reparaîtra, avec une autre preuve, à la page 129.

\pagerange{76}{79}
\section*{Le premier état de la mise au net (pages 76 à 79)}

Feuilles brunes portant, au crayon rouge, « Mettre 1 : n° 2 et cor.\ n° --- »,
et la numérotation (13) à (23) d'un texte plus long. La proposition 3 y est
la proposition 1 des pages 70 ; les corollaires 1 et 2 y sont ceux de la
mise au net, avec la précision que, pour $u$ à trace ou de Hilbert--Schmidt,
\[
  \|\Lambda^n u\|_1 = \alpha_n(|u|) = \alpha_n\bigl(\sqrt{u^{*}u}\bigr), \qquad
  \|\Lambda^n u\|_2 = \sqrt{\alpha_n(u^{*}u)},
\]
la première parce que $((\Lambda^n u)^{*}\Lambda^n u)^{1/2} = \Lambda^n |u|$,
la seconde parce que $(\Lambda^n u)^{*}\Lambda^n u = \Lambda^n(u^{*}u)$ est à
trace si $u^{*}u$ l'est. Le corollaire 4 est $|\alpha_n(u)| \leq
\alpha_n(|u|)$ --- les coefficients de $\det(1+zu)$ sont majorés en module
par ceux de $\det(1+z|u|)$ --- et le corollaire 5 est $|\det(1+zu)| \leq
\det(1 + |z|\,|u|)$.

\emph{Théorème 2.} Pour $u$ à trace dans $H$, $\det(1+u) = \prod_i (1 +
\lambda_i)$, et $\det(1+zu) = \prod(1 + z\lambda_i)$ est de genre zéro ;
en marge, $\alpha_n(u) = \sum_{i_1 < \cdots < i_n} \lambda_{i_1} \cdots
\lambda_{i_n}$ et $\operatorname{Tr} u = \sum \lambda_i(u)$. La démonstration
des pages 78 et 79 est barrée de deux traits obliques et se lit par bribes :
$\sum |\lambda_i| \leq \|u\|_1$ par les inégalités de Weyl ; les zéros de
$\det(1 - zu)$ ; ses coefficients majorés par ceux d'une fonction de genre
zéro ; et, pour la formule des coefficients dans le cas diagonalisable,
$u = \sum \lambda_i\,\bar{e}_i \otimes e_i$ et $\alpha_n(u) = \sum
\lambda_{i_1} \cdots \lambda_{i_n}\,\alpha_n(\bar{e}_{i_1} \otimes e_{i_1},
\ldots)$ où seuls comptent les indices distincts.\footnote{La page 79
écrit « $|z_1 \cdots z_n| = |\operatorname{Tr} pu| \leq \|u\|_1$ » pour une
somme ; la transcription le note. Ce que la démonstration barrée cherche,
c'est la proposition de la page 110, qui la complète : voir plus bas.}

\pagerange{80}{80}
\section*{Une feuille sur les fonctions convexes (page 80)}

Soit $f$ convexe sur un ouvert convexe $U$ et $G_f = \{(x,y) : x \in U,\ y
\geq f(x)\}$ son épigraphe. Par un point $(a, f(a))$ passe un hyperplan
$H$ ne rencontrant pas l'intérieur de $G_f$ ; il n'est pas vertical, donc
d'équation $y = L_a(x)$ avec $L_a$ affine, et $f(x) \geq L_a(x)$ sur $U$
avec égalité en $a$. Ainsi
\[
  f(x) = \sup_a L_a(x) .
\]
Si de plus $f$ est croissante en chacune de ses variables, les coefficients
$c_i$ de $L_a$ sont $\geq 0$ : sinon, $L_a$ décroissant en $x_1$ au voisinage
de $a_1$ tandis que $f$ croît, on aurait $L_a > f$ d'un côté. C'est
l'ingrédient de la page 82.

\pagerange{81}{92}
\section*{Le lemme de Weyl et les trois théorèmes de convexité (pages 81 à 92)}

Le numéro 4, « Inégalités de convexité et classes remarquables d'opérateurs,
compléments ». Tout y descend d'un lemme que la page 81 attribue à H.\ Weyl
et qui est, sous sa forme additive, le théorème de majoration de Hardy,
Littlewood et Pólya.

\emph{Théorème 1, 1°} (lemme fondamental). Soient $(\alpha_i)_{i \leq n}$ et
$(\beta_i)_{i \leq n}$ deux suites réelles, la première décroissante, telles
que
\[
  \alpha_1 + \cdots + \alpha_m \leq \beta_1 + \cdots + \beta_m \qquad
  (m = 1, \ldots, n),
\]
et $\varphi$ une fonction convexe symétrique sur $\mathbf{R}^n$. Si a)
$\varphi$ est croissante en chaque variable, ou si b) $\sum \alpha_i = \sum
\beta_i$, alors
\[
  \varphi(\alpha_1, \ldots, \alpha_n) \leq \varphi(\beta_1, \ldots, \beta_n) .
\]
\footnote{La page 81 dit la \emph{deuxième} suite décroissante, en deux mots
d'une lecture incertaine ; l'hypothèse porte sur la première, et la
démonstration l'utilise ainsi. L'hypothèse sur $(\beta_i)$ est superflue :
ses sommes partielles réordonnées en décroissant sont plus grandes. Sans
la décroissance de $(\alpha_i)$ l'énoncé est faux : $\alpha = (0, 2)$,
$\beta = (1, 1)$, $\varphi = \max$.}
Démonstration. Une fonction convexe sur $\mathbf{R}^n$ est continue et son
épigraphe est un convexe fermé, donc $\varphi = \sup_k L_k$ pour des
fonctions affines $L_k(t) = \sum a_i t_i + b$ ; si $\varphi$ est croissante,
on peut prendre tous les $a_i \geq 0$, car un $a_i < 0$ donnerait $L_k(t) \to
+\infty$ quand $t_i \to -\infty$ alors que $\varphi(t)$ reste bornée. Si
$\varphi$ est symétrique, $\sigma L_k \leq \varphi$ pour toute permutation
$\sigma$, donc $\varphi = \sup_k \sup_\sigma \sigma L_k$, et il suffit de
prouver, pour une $L$ fixée,
\[
  \sup_\sigma (\sigma L)(\alpha) \leq \sup_\sigma (\sigma L)(\beta) .
\]
On peut supposer $b = 0$. Le sup de gauche est atteint quand les $a_i$ sont
rangés en décroissant en face des $\alpha_i$ décroissants (inégalité de
réarrangement), et l'on est ramené à $\sum a_i\alpha_i \leq \sum a_i\beta_i$
pour $a_1 \geq \cdots \geq a_n$. Or, par sommation d'Abel,
\[
  \sum a_i\alpha_i = (a_1 - a_2)\,\alpha_1 + (a_2 - a_3)(\alpha_1 + \alpha_2)
  + \cdots + (a_{n-1} - a_n)(\alpha_1 + \cdots + \alpha_{n-1})
  + a_n(\alpha_1 + \cdots + \alpha_n),
\]
et de même pour $\beta$ ; les $n-1$ premiers coefficients sont $\geq 0$ et
multiplient des sommes partielles majorées, et le dernier terme est majoré
soit parce que $a_n \geq 0$ (cas a), soit parce que les sommes totales sont
égales (cas b).

\emph{Théorème 1, 2°} (forme multiplicative). Soit $f(\rho_1, \ldots,
\rho_n)$ une fonction des $\rho_i > 0$ telle que $f(e^{t_1}, \ldots, e^{t_n})$
soit convexe et symétrique. Soient $\rho_1 \geq \cdots \geq \rho_n \geq 0$ et
$\sigma_1, \ldots, \sigma_n \geq 0$ avec
\[
  \rho_1 \cdots \rho_m \leq \sigma_1 \cdots \sigma_m \qquad (m = 1, \ldots, n) .
\]
Si a) $f$ est croissante en chaque variable, définie et continue pour
$\rho_i \geq 0$, ou si b) $\rho_1 \cdots \rho_n = \sigma_1 \cdots \sigma_n
\neq 0$, alors $f(\rho) \leq f(\sigma)$. Pour des $\rho_i$, $\sigma_i > 0$
c'est le 1° avec $\alpha_i = \log\rho_i$, $\beta_i = \log\sigma_i$ ; dans le
cas a) on atteint les suites ayant des zéros en approchant $(\rho_i),
(\sigma_i)$ par des suites admissibles strictement positives --- les pages 81
et 82 y reviennent deux fois, encadrent et barrent leur premier essai.
\footnote{La symétrie de $f$ n'est pas dans l'énoncé de la page 81 et
figure en interligne page 90 ; elle est nécessaire, et supposée ici.}

\emph{Corollaire 1a.} Sous les conditions du 2°, b), on a aussi
$f(1/\rho_1, \ldots, 1/\rho_n) \leq f(1/\sigma_1, \ldots, 1/\sigma_n)$ : de
$\rho_1 \cdots \rho_n = \sigma_1 \cdots \sigma_n$ on tire $\rho_M \cdots
\rho_n \geq \sigma_M \cdots \sigma_n$ pour tout $M$, donc les suites
$(1/\rho_n, \ldots, 1/\rho_1)$ et $(1/\sigma_n, \ldots, 1/\sigma_1)$
satisfont encore aux hypothèses.

\emph{Remarque.} Si deux suites satisfont aux hypothèses du 1° (resp.\ du
2°), leurs $k$ premiers termes y satisfont aussi pour $1 \leq k \leq n$, et
l'inégalité vaut pour les fonctions de $k$ variables.

\emph{Exemples.} Si $f$ est convexe croissante des $\rho_i$, alors $f\circ
\exp$ est convexe, comme sup de fonctions $\sum a_i e^{t_i} + b$ avec
$a_i \geq 0$. Si $f(\rho)$ est une fonction d'une variable avec $f\circ\exp$
convexe, $\sum f(\rho_i)$ satisfait aux conditions du 2°, b), et à celles de
a) si de plus $f$ est croissante et continue en $0$ ; d'où
\[
  \sum_{i=1}^{n} f(\rho_i) \leq \sum_{i=1}^{n} f(\sigma_i), \qquad\text{et sous a)}\qquad
  \sum_{i=1}^{k} f(\rho_i) \leq \sum_{i=1}^{k} f(\sigma_i)\quad (k \leq n) .
\]
Pour $f(\rho) = \rho^p$ : $\sum \rho_i^p \leq \sum \sigma_i^p$ pour $p$ réel
sous b), et pour tout $k$ si $p \geq 0$. Enfin toute \emph{norme}
symétrique $N$ sur $\mathbf{R}^n$ avec $N(x) = N(|x|)$, croissante sur
l'orthant positif --- une norme de Schatten, dira le numéro 5 --- satisfait
aux conditions du 2°, a), d'où $N(\rho) \leq N(\sigma)$ et $N(\rho_1,
\ldots, \rho_k) \leq N(\sigma_1, \ldots, \sigma_k)$.

\emph{Théorème 2} (premier théorème de convexité dans les espaces de
Hilbert). Soit $u$ un opérateur compact dans un espace de Hilbert et $f$
une fonction de $n$ variables $\geq 0$, symétrique, croissante en chacune,
continue, avec $f\circ\exp$ convexe. Alors
\[
  f(|\lambda_1|, \ldots, |\lambda_n|) \leq f(s_1, \ldots, s_n),
\]
par le théorème 1, 2°, a), et l'inégalité $|\lambda_1 \cdots \lambda_m| \leq
s_1 \cdots s_m$ du numéro 3.

\emph{Corollaire 1.} Si $f$ est croissante sur $\mathbf{R}_{+}$, continue,
avec $f\circ\exp$ convexe : $\sum_{i \leq k} f(|\lambda_i|) \leq \sum_{i
\leq k} f(s_i)$ pour tout $k$, et $\sum_i f(|\lambda_i|) \leq \sum_i f(s_i)$,
la première somme convergeant si la seconde converge.
\emph{Corollaire 2.} Pour $p > 0$, $\sum_{i \leq k} |\lambda_i|^p \leq
\sum_{i \leq k} s_i^p$ et $\sum_i |\lambda_i|^p \leq \sum_i s_i^p = \|u\|_p^p$.
\emph{Corollaire 3.} Si $u$ est à trace (resp.\ de Hilbert--Schmidt), la
suite des valeurs propres est sommable (resp.\ de carré sommable), avec
$\sum |\lambda_i| \leq \|u\|_1$ (resp.\ $(\sum |\lambda_i|^2)^{1/2} \leq
\|u\|_2$).\footnote{Weyl, 1949 ; le corollaire 3 est ce que la page 129,
antérieure, ne savait obtenir que « par voie analytique ».}
\emph{Corollaire 4.} En dimension finie $n$, si $f\circ\exp$ est convexe et
symétrique sans être croissante, l'inégalité $f(|\lambda|) \leq f(s)$ vaut
encore pourvu que $u$ soit inversible ou que $f$ soit définie et continue
sur l'orthant fermé : c'est le cas b), grâce à
\[
  |\lambda_1 \cdots \lambda_n| = |\det u| = s_1 \cdots s_n,
\]
puisque $(s_1 \cdots s_n)^2 = \det(u^{*}u) = |\det u|^2$.
\emph{Corollaire 5.} Pour $N$ une norme de Schatten sur $\mathbf{R}^k$,
$k \leq \dim E$ : $N(|\lambda_1|, \ldots, |\lambda_k|) \leq N(s_1, \ldots,
s_k)$ ; et si $N$ est une norme de Schatten sur $\mathbf{R}^{(\mathbf{N})}$
--- l'espace des suites à support fini ---, prolongée par $N(\xi) = \lim_n
N(\xi_1, \ldots, \xi_n, 0, \ldots)$, alors $N((|\lambda_i|)) \leq N((s_i))$.

\emph{Théorème 3} (deuxième théorème de convexité). Soient $u \in L_0(E,F)$,
$v \in L_0(F,G)$ et $f$ comme au théorème 2. Alors
\[
  f\bigl(s_1(vu), \ldots, s_n(vu)\bigr) \leq
  f\bigl(s_1(v)s_1(u), \ldots, s_n(v)s_n(u)\bigr),
\]
par le théorème 1, 2°, a), et l'inégalité de Horn du numéro 3.\footnote{La
page 90 énonce $f$ sans monotonie, ajoute en interligne « symétrique et
croissante », et sa démonstration renvoie à l'hypothèse b) du théorème 1,
qui exigerait $\prod_{i \leq n} s_i(vu) = \prod_{i \leq n} s_i(v)s_i(u)$ ---
vrai en dimension finie $n$ (corollaire 4 ci-dessous), faux en général. Pour
$u$, $v$ compacts quelconques c'est l'hypothèse a) qu'il faut, et la
monotonie de $f$ ; c'est ce qu'on a écrit.}
\emph{Corollaire 1.} Pour $f$ continue croissante sur $\mathbf{R}_{+}$ avec
$f\circ\exp$ convexe, $\sum_{i \leq n} f(s_i(vu)) \leq \sum_{i \leq n}
f(s_i(v)s_i(u))$ pour tout $n$, et de même pour les sommes infinies. Un
feuillet collé page 90 en donne le sens : les $s_i(vu)$ sont majorés, au
sens des fonctions de Weyl, par ce qu'ils seraient si $u$ et $v$ étaient
les multiplications par $(s_i(u))$ et $(s_i(v))$ dans $\ell^2$.
\emph{Corollaire 2} (« inégalité de Hölder pour opérateurs »). Pour
$\frac1p + \frac1q = \frac1s$, $p, q, s > 0$, et tout $n$,
\[
  \Bigl(\sum_{i \leq n} s_i(vu)^s\Bigr)^{1/s} \leq
  \Bigl(\sum_{i \leq n} s_i(u)^p\Bigr)^{1/p}\Bigl(\sum_{i \leq n} s_i(v)^q\Bigr)^{1/q},
  \qquad\text{d'où}\qquad
  \|vu\|_s \leq \|u\|_p\,\|v\|_q,
\]
en prenant $f(\rho) = \rho^s$ dans le corollaire 1 puis l'inégalité de
Hölder ordinaire. Le cas le plus important est $s = 1$ : $\|vu\|_1 \leq
\|u\|_p\|v\|_{p'}$ pour $\frac1p + \frac1{p'} = 1$.
\emph{Corollaire 3} (en marge). $|\operatorname{Tr} vu| \leq \sum_i
s_i(u)\,s_i(v)$, par $f(s) = s$ et $|\operatorname{Tr} w| \leq \sum
s_i(w)$.\footnote{L'inégalité de trace de von Neumann (1937), obtenue ici
comme cas $f(s) = s$ du corollaire 1.}
\emph{Corollaire 4} (feuillet collé). En dimension finie $n$, si $f\circ\exp$
est convexe et symétrique, et si $u$ et $v$ sont inversibles ou $f$ continue
sur l'orthant fermé, $f(s(vu)) \leq f(s(u)s(v))$, grâce à $\prod s_i(vu) =
\prod s_i(u)s_i(v)$, les deux membres valant $|\det(vu)|$ et $|\det u|\,|\det
v|$. Un crayon en marge demande : « Est-ce utile ? »

\emph{Théorème 4} (troisième théorème de convexité). Pour $u$, $v \in
L_0(E,F)$ et $\varphi$ convexe, symétrique et croissante de $n$ variables
$\geq 0$,
\[
  \varphi\bigl(s_1(u+v), \ldots, s_n(u+v)\bigr) \leq
  \varphi\bigl(s_1(u) + s_1(v), \ldots, s_n(u) + s_n(v)\bigr),
\]
cas particulier du théorème 1, 1°, a), et de $\sum_{i \leq m} s_i(u+v) \leq
\sum_{i \leq m} (s_i(u) + s_i(v))$ du numéro 1.\footnote{Ky Fan (1951). Un
corollaire $\sum \varphi(s_i(u+v)) \leq \sum \varphi(s_i(u) + s_i(v))$ est
écrit, encadré et barré ; il est vrai pour $\varphi$ convexe croissante.}

\emph{Corollaire 1.} Soit $N$ une norme de Schatten sur $\mathbf{R}^n$ et,
pour $u$ compact, $\|u\|_N = N(s_1(u), \ldots, s_n(u))$. Alors
\[
  \|u + v\|_N \leq \|u\|_N + \|v\|_N,
\]
car $N(s(u+v)) \leq N(s(u) + s(v)) \leq N(s(u)) + N(s(v))$.\footnote{Avec
$\|\lambda u\|_N = |\lambda|\,\|u\|_N$, biffé sur la page : $u \mapsto
\|u\|_N$ est une norme sur les compacts de rang $\leq n$. C'est le théorème
de von Neumann (1937) sur les normes unitairement invariantes.}
\emph{Corollaire 2.} Pour $1 \leq p \leq \infty$ et tout $n$,
$(\sum_{i \leq n} s_i(u+v)^p)^{1/p} \leq (\sum_{i \leq n} s_i(u)^p)^{1/p} +
(\sum_{i \leq n} s_i(v)^p)^{1/p}$, d'où
\[
  \|u + v\|_p \leq \|u\|_p + \|v\|_p .
\]

\pagerange{93}{100}
\section*{Classes remarquables d'opérateurs (pages 93 à 100)}

Le numéro 5 est « un exposé de la théorie de Schatten, simplifié par les
inégalités de convexité de Weyl, en y apportant quelques compléments ».

\emph{Définition 1.} Une norme $N$ sur $\mathbf{R}^n$ est une \emph{norme de
Schatten} si elle est symétrique, si $N(x_1, \ldots, x_n) = N(|x_1|, \ldots,
|x_n|)$, si elle est croissante sur l'orthant positif, et si $N(1, 0, \ldots,
0) = 1$.\footnote{Ce sont les \emph{fonctions de jauge symétriques} de von
Neumann (1937) et Schatten, les \emph{fonctions normantes symétriques} de
Gohberg et Krein. La monotonie sur l'orthant est en fait conséquence des
deux premières conditions pour une norme ; la page la met dans la
définition.} Une norme sur $\mathbf{R}^{(\mathbf{N})}$ est de Schatten si
elle induit une norme de Schatten sur chaque $\mathbf{R}^n$. Exemples : les
$N_p(x) = (\sum |x_i|^p)^{1/p}$, $1 \leq p < \infty$, et $N_\infty(x) =
\sup |x_i|$. Toute norme de Schatten $N_0$ sur $\mathbf{R}^n$ est induite
par une norme de Schatten sur $\mathbf{R}^{(\mathbf{N})}$, par exemple
$N(x) = N_0(x^{*}_1, \ldots, x^{*}_n)$ où $x^{*}$ est la suite des $|x_i|$
rangée en décroissant ; la seule chose à vérifier est l'inégalité
triangulaire, et elle résulte du corollaire 1 du théorème 4, $x$ étant vu
comme l'opérateur de multiplication par $x$ dans $\ell^2$. Il y a donc une
infinité continue de normes de Schatten sur $\mathbf{R}^{(\mathbf{N})}$
équivalentes à la norme de $c_0$.

\emph{Les espaces $\ell^N$ et $\ell_N$.} Pour $N$ une norme de Schatten sur
$\mathbf{R}^{(\mathbf{N})}$ et $x$ une suite complexe, posons $|x|_n =
(|x_1|, \ldots, |x_n|, 0, \ldots)$ et
\[
  N(x) = \lim_n N(|x|_n) \in [0, +\infty] .
\]
On a $N(\lambda x) = |\lambda| N(x)$ et $N(x+y) \leq N(x) + N(y)$, car
$|x+y|_n \leq |x|_n + |y|_n$ et $N$ croît sur l'orthant. Soit $\ell^N$
l'espace des $x$ avec $N(x) < \infty$ : $N$ est une norme sur $\ell^N$. Si
$N' \leq N''$, alors $\ell^{N''} \subset \ell^{N'}$ avec injection de norme
$\leq 1$ ; comme toute norme de Schatten est comprise entre $N_\infty$ et
$N_1$,
\[
  \ell^1 \subset \ell^N \subset \ell^\infty,
\]
les injections étant de norme $\leq 1$. On a de plus, pour toute suite
$(x^k)$ dans $\ell^N$ avec $\sum_k N(x^k) < \infty$, la sous-additivité
dénombrable $N(\sum_k x^k) \leq \sum_k N(x^k)$, où la série converge dans
$\ell^\infty$ ; et il en résulte que $\ell^N$ est \emph{complet}, une série
normalement convergente convergeant vers sa somme dans $\ell^\infty$ puis,
par cette inégalité appliquée aux restes, dans $\ell^N$.

Soit $\ell_N$ l'adhérence de $\mathbf{R}^{(\mathbf{N})}$ dans $\ell^N$,
c'est-à-dire le complété de $\mathbf{C}^{(\mathbf{N})}$ pour $N$. Si $N' \leq
N''$, $\ell_{N''} \subset \ell_{N'}$,\footnote{La page 96 écrit
l'inclusion dans l'autre sens, ce que la transcription note ; le sens est
celui des $\ell^N$, l'adhérence d'un sous-espace dans le plus petit espace
étant contenue dans son adhérence dans le plus grand.} et
\[
  \ell^1 \subset \ell_N \subset c_0,
\]
les espaces $\ell_N$ pour $N_1$ et $N_\infty$ étant $\ell^1$ et $c_0$. Une
formule $\ell_N = \ell^N \cap c_0$ est écrite, numérotée (9), puis barrée ;
elle est fausse en général.\footnote{Pour $N = N_\infty$ elle est vraie ;
pour une norme de Lorentz ou d'Orlicz elle ne l'est pas : $\ell_N$ est le
sous-espace des $x \in \ell^N$ dont les queues tendent vers zéro pour $N$,
ce que la page 98 dira exactement. La proposition 1 annoncée page 96 n'a
pas d'énoncé lisible ; d'après le plan de la page 101, c'est : $\ell_N
\subset c_0$ strictement, sauf si $N$ est équivalente à $N_\infty$.}

\emph{La norme polaire.} $\mathbf{R}^{(\mathbf{N})}$ est en dualité avec
lui-même par $\langle x, y\rangle = \sum x_i y_i$, et toute norme $N$ y a
une \emph{norme polaire}
\[
  N^{\circ}(y) = \sup_{N(x) \leq 1} |\langle x, y\rangle| ,
\]
avec $(N^{\circ})^{\circ} = N$ ; si $N$ est de Schatten, $N^{\circ}$ l'est.
Pour $x \in \ell^N$, $y \in \ell^{N^{\circ}}$ on a $\sum |x_i y_i| \leq N(x)\,
N^{\circ}(y)$, d'où un accouplement naturel entre $\ell^N$ et
$\ell^{N^{\circ}}$, et :

\emph{Le dual de $\ell_N$ est $\ell^{N^{\circ}}$}, avec sa norme, pour cet
accouplement. Une forme linéaire continue $X'$ sur $\ell_N$ se restreint
à $\ell^1$, donc est donnée par un $x' \in \ell^\infty$ ; alors
$N^{\circ}(|x'|_n) \leq \|X'\|$ pour tout $n$, donc $N^{\circ}(x') \leq
\|X'\|$. En général le dual de $\ell^N$ n'est pas $\ell^{N^{\circ}}$, car
$\ell_N$ n'est pas dense dans $\ell^N$ ; exemple $N = N_\infty$.\footnote{C'est
le dual de Köthe d'un espace de suites symétrique ; Schatten (1950), et
pour la « dualité $\ell_N$--$\ell^{N^{\circ}}$ » la théorie des espaces
de suites symétriques de Garling. Le passage de $\ell^\infty$ à
$N^{\circ}(x') \leq \|X'\|$ utilise que $N^{\circ}(|x'|_n)$ se calcule sur
$\mathbf{R}^n \subset \ell_N$.}

\emph{Proposition.} $\ell_N$ est réflexif si et seulement si $\ell_N = \ell^N$
et $\ell_{N^{\circ}} = \ell^{N^{\circ}}$. Si $\ell_N$ est réflexif, son bidual
$(\ell^{N^{\circ}})'$, qui contient $\ell^N$ isométriquement par
l'accouplement, est $\ell_N$, d'où $\ell_N = \ell^N$ ; et $\ell^{N^{\circ}}$,
dual d'un réflexif, est réflexif, d'où de même $\ell_{N^{\circ}} =
\ell^{N^{\circ}}$. Réciproquement, sous ces deux égalités, le dual de $\ell_N$
est $\ell_{N^{\circ}}$ et le dual de celui-ci est $\ell_N$, par l'accouplement
canonique.\footnote{La page 97 se reproche une « démonstration circulaire
de l'équivalence des 4 conditions », biffe quatre lignes, et en donne une
autre ; celle-ci est la sienne, mise au propre. Elle utilise $(N^{\circ})^{\circ}
= N$, écrit en marge de la page 96.}

\emph{Les classes $\mathcal{S}^N(E,F)$.} Soit $N$ une norme de Schatten sur
$\mathbf{R}^{(\mathbf{N})}$, non équivalente à $N_\infty$, et $E$, $F$ deux
espaces de Hilbert. On note $\mathcal{S}^N(E,F)$ le sous-espace de $L_0(E,F)$
des $u$ tels que $(s_i(u)) \in \ell^N$, et
\[
  N(u) = N\bigl((s_i(u))\bigr) < \infty .
\]
C'est une norme : $N(u+v) = \lim_n N(|s(u+v)|_n) \leq \lim_n \bigl(N(|s(u)|_n)
+ N(|s(v)|_n)\bigr)$ par le corollaire 1 du théorème 4. Soit
$\mathcal{S}_N(E,F)$ le sous-espace des $u$ avec $(s_i(u)) \in \ell_N$ ;
c'est l'adhérence de $\bar{E} \otimes F$ dans $\mathcal{S}^N(E,F)$, puisque
$u = \sum s_i\,\bar{e}_i \otimes f_i$ est limite de ses sommes partielles
exactement quand $N(s_{n+1}, s_{n+2}, \ldots) \to 0$, ce qui caractérise
$\ell_N$ dans $\ell^N$.

\emph{Multiplicativité.} Si $N$, $N'$, $N''$ sont des normes de Schatten
avec $N(xy) \leq N'(x)\,N''(y)$ sur $\mathbf{R}^{(\mathbf{N})}$, la même
inégalité vaut sur $\ell^{N'} \times \ell^{N''}$, et pour $u \in
\mathcal{S}^{N'}(E,F)$, $v \in \mathcal{S}^{N''}(F,G)$ on a $vu \in
\mathcal{S}^{N}(E,G)$ avec
\[
  N(vu) \leq N'(u)\,N''(v),
\]
par le corollaire 1 du théorème 3 appliqué à la norme $N$. En particulier,
pour les normes polaires, $\|vu\|_1 \leq N(u)\,N^{\circ}(v)$ ; et $N(u^{*}) =
N(u)$, $N(vu) \leq \|v\|\,N(u)$, $N(uv) \leq \|v\|\,N(u)$, par $s_i(vu) \leq
\|v\| s_i(u)$.

\emph{Le dual de $\mathcal{S}_N(E,F)$ est $\mathcal{S}^{N^{\circ}}(F,E)$}, par
$\langle u, v\rangle = \operatorname{Tr} vu$. Le dual de $\mathcal{S}_N(E,F)$
s'identifie à un sous-espace du dual de $\mathcal{S}_{N_1}(E,F) =
\mathcal{S}_1(E,F)$, donc de $L(F,E)$ ; ce sous-espace contient
$\mathcal{S}^{N^{\circ}}(F,E)$, isométriquement par l'inégalité précédente.
Réciproquement, soit $v \in L(F,E)$ tel que $|\operatorname{Tr} vu| \leq
M\,N(u)$ pour $u$ de rang fini. Par $v = U|v|$, $|v| = Vv$, on se ramène à
$F = E$ et $v \geq 0$ ; pour tout projecteur $P$ de rang fini de la
décomposition spectrale de $v$, $v = \sum \rho_i\,\bar{e}_i \otimes e_i$, les
$u$ diagonaux dans la base $(e_i)$ montrent que $(\rho_i)$ définit une forme
continue sur $\ell_N$, donc $(\rho_i) \in (\ell_N)' = \ell^{N^{\circ}}$, ce
qui est la conclusion.\footnote{Le bloc de la page 98 qui énonce cette
dualité est encadré et barré, puis repris page 99 avec un point
d'interrogation rouge sur la dernière ligne ; la page 102 la réécrit sans
le « $\circ$ » ni l'ordre $(F,E)$ ; le théorème 1 de la page 104 la donne
sous la forme ci-dessus, qui est la bonne. L'argument de la page 99 est
elliptique après « $v \geq 0$ » ; on l'a complété. C'est le théorème de
dualité de Schatten (1950).}

\emph{Corollaire 1.} $u \in \mathcal{S}^N(E,F)$ si et seulement si $|u| \in
\mathcal{S}^N(E)$, et $N(u) = N(|u|)$ ; de même pour $\mathcal{S}_N$.
\emph{Corollaire 2.} Si $u = h + ik$ avec $h$, $k$ hermitiens, $u \in
\mathcal{S}^N(E)$ si et seulement si $h$, $k \in \mathcal{S}^N(E)$, et de
même pour $\mathcal{S}_N$.

\emph{Remarque} (caractérisation axiomatique). Soit $p$ une norme sur
$\bar{E} \otimes F$, $E$ et $F$ de dimension infinie, telle que a) $p(VuU) =
p(u)$ pour $U$, $V$ unitaires, et b) $p(\bar{x} \otimes y) = \|x\|\,\|y\|$.
Pour $(e_i)$, $(f_i)$ orthonormaux infinis, $(\lambda_i) \mapsto p(\sum
\lambda_i\,\bar{e}_i \otimes f_i)$ est sur $\mathbf{R}^{(\mathbf{N})}$ une
norme symétrique (par a), normalisée (par b), fonction de $(|\lambda_i|)$
(par a), donc une norme de Schatten $N$, et $p(u) = N(u)$ sur $\bar{E}
\otimes F$ par a). Ce sont donc les normes $N(u)$ sur les opérateurs de rang
fini.\footnote{Le théorème de von Neumann sur les normes unitairement
invariantes, en dimension infinie. L'axiome b) y est plus fort qu'il ne
faut : $p(\bar{x} \otimes y) = 1$ pour $x$, $y$ unitaires suffit.}

\pagerange{101}{104}
\section*{Espaces de Schatten : le théorème récapitulatif (pages 101 à 104)}

Deux pages de plan, puis un « Théorème 1 » en cinq points dont les formules
(18) à (27) continuent la numérotation du numéro 5. Le plan ajoute au
numéro 5 : $N(x)$ est fonction croissante de $|x|$ ; $\ell_N \subsetneq c_0$
sauf si $N \sim N_\infty$ ; $(N_p)^{\circ} = N_{p'}$ et l'on retrouve la
réflexivité des $\ell^p$, $1 < p < \infty$ ; $\mathcal{S}^N(E,F) \subset
L_0(E,F)$ si $N \not\sim N_\infty$, car si $|u|$ n'était pas compact il aurait
un projecteur spectral $P$ de rang infini sur lequel $|u| \geq \varepsilon
> 0$, et $v = |u|^{-1}P$, hermitien borné, donnerait $v|u| = P \in
\mathcal{S}^N$, donc $(1, 1, \ldots) \in \ell^N$ et
$N \sim N_\infty$ ; $\mathcal{S}^{N_p} = \mathcal{S}_p$ et, pour $p < \infty$,
$\mathcal{S}_{N_p} = \mathcal{S}_p$ ; les corollaires 1 et 2 sur $|u|$ et
$h + ik$ ; et la remarque que $\ell^N$ se plonge isométriquement dans
$\mathcal{S}^N(E,F)$ par $(\lambda_i) \mapsto \sum \lambda_i\,\bar{e}_i
\otimes f_i$ quand $E$ et $F$ sont de dimension infinie.

\emph{Théorème 1.} Soit $N$ une norme de Schatten sur $\mathbf{R}^{(\mathbf{N})}$
et $E$, $F$ deux espaces de Hilbert. Pour $u \in L_0(E,F)$ posons $N(u) =
N((s_i(u)))$, et pour $u \in L(E,F)$
\[
  N(u) = \sup\bigl\{ |\operatorname{Tr} vu| : v \in \bar{F} \otimes E,\ N^{\circ}(v)
  \leq 1 \bigr\} ;
\]
les deux expressions coïncident sur $L_0(E,F)$.\footnote{Sur les compacts,
c'est la dualité $\ell_{N^{\circ}}$--$\ell^{N}$ jointe à l'inégalité de trace
de von Neumann ; la seconde définition étend $N$ à tout $L(E,F)$, où elle
vaut $+\infty$ hors de $\mathcal{S}^N$ dès que $N \not\sim N_\infty$.}

1. Soit $\mathcal{S}^N(E,F)$ le sous-espace des $u \in L(E,F)$ avec $N(u) <
\infty$ : $N$ en fait un espace de Banach. Un $u$ compact y est si et
seulement si $(s_i(u)) \in \ell^N$. L'adhérence $\mathcal{S}_N(E,F)$ de
$\bar{E} \otimes F$ dans $\mathcal{S}^N(E,F)$ est l'ensemble des $u$ compacts
avec $(s_i(u)) \in \ell_N$.

2. $N(u) = N(u^{*})$, donc $u \in \mathcal{S}^N(E,F)$ si et seulement si
$u^{*} \in \mathcal{S}^N(F,E)$, et de même pour $\mathcal{S}_N$ ; et pour
$u \in L(E,F)$, $v \in L(F,G)$, $N(vu) \leq \|v\|\,N(u)$ et $N(vu) \leq
\|u\|\,N(v)$.\footnote{La page 103 écrit les deux inégalités avec les rôles
de $u$ et $v$ échangés d'une à l'autre, et un point d'interrogation rouge en
marge ; ce sont les deux inégalités du plan de la page 101.} Sur les
hermitiens positifs compacts, $N(u)$ est fonction croissante de $u$.

3. Si $N(xy) \leq N'(x)\,N''(y)$ sur $\mathbf{R}^{(\mathbf{N})}$, alors pour
$u \in L(E,F)$, $v \in L(F,G)$ : $N(vu) \leq N'(u)\,N''(v)$ ; donc $u \in
\mathcal{S}^{N'}(E,F)$ et $v \in \mathcal{S}^{N''}(F,G)$ entraînent $vu \in
\mathcal{S}^N(E,G)$, et $vu \in \mathcal{S}_N(E,G)$ si de plus $u \in
\mathcal{S}_{N'}$ ou $v \in \mathcal{S}_{N''}$. En particulier $\|vu\|_1 \leq
N(u)\,N^{\circ}(v)$ pour $u \in L(E,F)$, $v \in L(F,E)$, et $|\operatorname{Tr}
vu| \leq N(u)\,N^{\circ}(v)$ pour $u \in \mathcal{S}^N(E,F)$, $v \in
\mathcal{S}^{N^{\circ}}(F,E)$. La démonstration pour $u$, $v$ non compacts
passe par $N(u) = \lim N(pup)$ sur les projecteurs $p$ de rang fini croissant
vers $1$, puis par le deuxième théorème de convexité en dimension finie.

4. Par l'accouplement $\langle u, v\rangle = \operatorname{Tr} vu$, le dual de
$\mathcal{S}_N(E,F)$ s'identifie à $\mathcal{S}^{N^{\circ}}(F,E)$. Si $E$ et
$F$ sont de dimension infinie, les conditions suivantes sont équivalentes :
a) $\mathcal{S}_N(E,F)$ est réflexif ; b) $\mathcal{S}^N(E,F)$ est réflexif ;
c) $\ell_N$ est réflexif ; d) $\ell^N$ est réflexif ; e) $\ell_{N^{\circ}}$
est réflexif ; f) $\ell^{N^{\circ}}$ est réflexif ; g) $\mathcal{S}_N(E,F) =
\mathcal{S}^N(E,F)$ et $\mathcal{S}_{N^{\circ}}(F,E) = \mathcal{S}^{N^{\circ}}(F,E)$.
Car $\ell^N$ (resp.\ $\ell_N$) est un sous-espace fermé de
$\mathcal{S}^N(E,F)$ (resp.\ $\mathcal{S}_N(E,F)$), et un sous-espace fermé
d'un réflexif est réflexif ; $\ell_N$ réflexif donne $\ell_N = \ell^N$ et
$\ell_{N^{\circ}} = \ell^{N^{\circ}}$ par la proposition de la page 97, donc
$\mathcal{S}_N = \mathcal{S}^N$ et de même pour $N^{\circ}$, donc le dual de
$\mathcal{S}_N(E,F)$ est $\mathcal{S}_{N^{\circ}}(F,E)$ et le dual de celui-ci
est $\mathcal{S}_N(E,F)$, qui est réflexif.\footnote{La liste de la page 104
est remaniée, avec deux « d) » ; l'ordre et la lettre g) sont de cette
lecture. Sans l'hypothèse de dimension infinie les $\ell$ ne se plongent pas
et l'implication de $\mathcal{S}$ vers $\ell$ tombe ; l'implication de
$N$ réflexive vers $\mathcal{S}_N(E,F)$ réflexif reste vraie (corollaire 2 de
la page 102).} \emph{Corollaire} : pour $1 < p < \infty$, $\mathcal{S}_p(E,F)$
est réflexif, de dual $\mathcal{S}_{p'}(F,E)$.

5. Pour $u$ compact, $N((\lambda_i(u))) \leq N(u)$, et $u \in
\mathcal{S}^N(E,F)$ entraîne $(\lambda_i(u)) \in \ell^N$ : c'est le
corollaire 5 du premier théorème de convexité.

\pagerange{106}{111}
\section*{Fonctions entières : factorisation, genre, ordre (pages 106 à 111)}

D'autres feuillets, d'une écriture posée, sur du papier dont le verso
imprimé transparaît. Ce sont des rappels, et ils servent à une chose : la
démonstration du théorème 1 de la page 70, que la proposition de la page
110 achève.

\emph{Théorème 1} (Weierstrass). Soit $f$ entière avec $f(0) \neq 0$ et
$(a_n)$ la suite de ses zéros rangés par modules croissants. Il existe des
entiers $p_n \geq 0$ et une fonction entière $g$ tels que
\[
  f(z) = e^{g(z)} \prod_n \Bigl(1 - \frac{z}{a_n}\Bigr)
  \exp\Bigl(\frac{z}{a_n} + \frac{z^2}{2a_n^2} + \cdots + \frac{z^{p_n}}{p_n a_n^{p_n}}\Bigr),
\]
le produit convergeant absolument et uniformément sur tout borné ; toute
suite $(p_n)$ avec $\sum_n |r/a_n|^{p_n + 1} < \infty$ pour tout $r$ convient,
et si $\sum 1/|a_n|^{k+1} < \infty$ on peut prendre $p_n = k$ pour tout $n$.

\emph{Genre.} Si $\sum 1/|a_n|^{k+1} < \infty$ pour un $k$, le plus petit
étant $k$, et si dans cette représentation $g$ est un polynôme de degré $r$,
le \emph{genre} de $f$ est $\max(r, k)$. Genre zéro signifie $f(z) = C
\prod_n (1 - z/a_n)$ avec $\sum 1/|a_n| < \infty$.

\emph{Ordre.} Avec $M(r) = \sup_{|z| = r} |f(z)|$, l'\emph{ordre} de $f$ est
$\omega = \limsup \frac{\log\log M(r)}{\log r} \in [0, +\infty]$, borne
inférieure des $\sigma$ tels que $M(r) \leq e^{r^\sigma}$ pour $r$ grand.

\emph{Théorème 2} (Hadamard--Borel). Si $\omega < \infty$, alors
$\sum' 1/|a_n|^{\omega + \varepsilon} < \infty$ pour tout $\varepsilon > 0$,
la somme portant sur les zéros non nuls, et $f(z) = e^{P(z)}\prod_n (1 -
z/a_n)\exp(\cdots + z^p/(p a_n^p))$ avec $p \leq \omega$ et $P$ polynôme de
degré $\leq \omega$ : le genre est $\leq$ l'ordre. \emph{Corollaire} : si
$\omega < 1$, $f$ est de genre zéro. \emph{Remarque} : inversement l'ordre
est $\leq$ genre $+ 1$, d'où $\omega - 1 \leq g \leq \omega$, et $g$ est
déterminé par $\omega$ sauf si $\omega$ est entier. Exemples : $e^{z^\omega}$
est d'ordre et de genre $\omega$ ; $\prod (1 - z/n^\alpha)$, $\alpha > 1$,
est d'ordre $1/\alpha$ et de genre zéro. L'ordre d'un produit ou d'une somme
est au plus le maximum des ordres, de sorte que les fonctions d'ordre $<
\rho$ (resp.\ $\leq \rho$) forment une algèbre.

\emph{Théorème 3} (Hadamard). Pour $f(z) = \sum c_n z^n$, l'ordre se lit sur
les coefficients ; la page 109 s'interrompt au bord du feuillet après
« si $n^\alpha\sqrt[n]{|c_n|} \to 0$ ».\footnote{L'énoncé complet : $\omega =
\limsup \frac{n\log n}{-\log|c_n|}$. Ce que la suite utilise est le cas
limite : une fonction de genre zéro vérifie $|f(z)| \leq C_\varepsilon
e^{\varepsilon|z|}$ pour tout $\varepsilon > 0$ (elle est d'ordre $\leq 1$ et
de type minimal), ce que la page 110 invoque comme connu.}

\emph{Proposition.} Soit $f$ entière dont les coefficients sont majorés en
module par ceux d'une fonction $h$ de genre zéro à coefficients positifs,
et dont les zéros non nuls $z_i$ vérifient $\sum 1/|z_i| < \infty$. Alors $f$
est de genre zéro.

Démonstration. On peut supposer $h(z) = K\prod (1 + \alpha_i z)$ avec
$\alpha_i \geq 0$ sommables ; alors $|f(z)| \leq h(|z|) \leq C_\varepsilon
e^{\varepsilon |z|}$ pour tout $\varepsilon$, $f$ est d'ordre $\leq 1$ et,
par Hadamard--Borel et l'hypothèse sur les zéros, $f(z) = e^{\alpha + \beta
z} z^m \prod_i (1 - z/z_i)$. Il faut voir $\beta = 0$. Sinon, quitte à
remplacer $f(z)$ par $f(\theta z)$ avec $\theta = \bar\beta/|\beta|$, on a
$\beta > 0$. Posons $g(z) = z^m\prod_i (1 - z/z_i)$, de genre zéro, donc
$|g(z)| \leq K' e^{(a/3)|z|}$ pour tout $a > 0$ et $K'$ dépendant de $a$. Sur
le secteur $|\arg z| \leq \theta$ avec $\pi/3 < \theta < \pi/2$, on a
$\operatorname{Re} z \geq c|z|$ avec $c > 0$, donc $|f(z)| \geq K e^{\beta c
|z|}\,|g(z)|$, et avec $\varepsilon < \beta c$ il vient $|g(z)| \leq K
e^{-a|z|}$ sur ce secteur, $a = \beta c - \varepsilon > 0$. Soit $\omega =
e^{2i\pi/3}$ et $G(z) = g(z)\,g(\omega z)\,g(\omega^2 z)$ : pour tout $z$,
l'un des trois points $z$, $\omega z$, $\omega^2 z$ est dans le secteur,
puisque $\theta > \pi/3$, et
\[
  |G(z)| \leq \bigl(K e^{-a|z|}\bigr)\bigl(K' e^{\frac{a}{3}|z|}\bigr)^2
  = K'' e^{-\frac{a}{3}|z|} .
\]
$G$ est donc bornée, donc constante ; mais $G(z) = z^{3m}\prod_i (1 -
z^3/z_i^3)$, ce qui force $m = 0$ et l'absence de zéros : $g$ est constante,
$f(z) = Ce^{\beta z}$, incompatible avec $|f(z)| \leq C_\varepsilon
e^{\varepsilon|z|}$ pour $\varepsilon < \beta$.\footnote{Le second angle de la
page 111 est serré en bout de ligne et lu $\pi/2$ ; c'est ce que l'argument
demande, et le croquis en marge --- un disque coupé en trois secteurs par
trois rayons --- est celui de $G$. La fin de la démonstration est écrite de
bas en haut dans la marge.}

\pagerange{112}{118}
\section*{Exposant de convergence, fonction de comptage (pages 112 à 118)}

Soit $(\alpha_n)$ une suite décroissante de nombres positifs tendant vers
zéro. Soit $\rho$ la borne inférieure des $s > 0$ tels que $\alpha_n =
O(n^{-1/s})$, et $\sigma$ celle des $s > 0$ tels que $\sum \alpha_n^s <
\infty$ --- l'\emph{exposant de convergence} ; on a aussi $1/\rho = \limsup
\frac{\log(1/\alpha_n)}{\log n}$. Le critère de Riemann donne $\sigma \leq
\rho$, et l'on a $\rho \leq \sigma$, par le

\emph{Lemme.} Si $\sum \alpha_n^s < \infty$ et $(\alpha_n)$ décroît, alors
$\alpha_n = o(n^{-1/s})$. On se ramène à $s = 1$ ; si l'on avait
$\alpha_{n_i} \geq 1/n_i$ le long d'une suite strictement croissante
$(n_i)$, alors $\alpha_{n_{i-1}+1} + \cdots + \alpha_{n_i} \geq (n_i -
n_{i-1})/n_i =: \varepsilon_i$ avec $\sum \varepsilon_i < \infty$, et $n_i =
n_{i-1}/(1 - \varepsilon_i) = \cdots = n_0/\prod_{k \leq i}(1 - \varepsilon_k)$
serait bornée.

Si $\sum \alpha_i < \infty$ et $f(z) = \prod (1 + \alpha_i z)$, alors $\rho$
est aussi l'ordre de $f$. On sait $\rho \leq$ ordre ; pour l'inégalité
inverse il suffit du cas $\alpha_n = n^{-\beta}$, $\beta > 1$, et d'y
prouver ordre $\leq 1/\beta$ : on coupe $f(r) = \prod_{n \leq \nu}(1 +
r/n^\beta)\prod_{n > \nu}(1 + r/n^\beta)$, le second facteur est $\leq
\exp(r\sum_{n > \nu} n^{-\beta}) \leq \exp\bigl(r\,\nu^{1-\beta}/(\beta -
1)\bigr)$, d'ordre $\leq 1/\beta$ en $r$ si $\nu = [r^{1/\beta}] + 1$, et le
premier est $\leq (1+r)^\nu = e^{\nu\log(1+r)}$, d'ordre $\leq 1/\beta$
aussi.\footnote{C'est le théorème de Borel : l'ordre d'un produit canonique
de genre zéro est l'exposant de convergence de ses zéros.} Si $f(z) =
\prod(1 + \lambda_i z)$ est de genre zéro et $F(z) = \prod(1 + |\lambda_i|z)$,
$f$ et $F$ ont même ordre, puisque ordre de $f \leq$ ordre de $F$ = ordre de
la suite $(|\lambda_i|) \leq$ ordre de $f$ par Hadamard.\footnote{La page
114 écrit « de genre 1 » et souligne le mot ; avec $\sum|\lambda_i| < \infty$
c'est le genre zéro qui est en jeu, et c'est ce que la conclusion utilise.}
On peut donc parler de l'\emph{ordre} d'un ensemble dénombrable de nombres
complexes, borne inférieure des $s > 0$ avec $\sum |z_i|^s < \infty$ ; une
suite à décroissance rapide est une suite d'ordre zéro. Suivent trois
lignes sur l'ordre d'une suite dans un espace de Banach et d'une partie
bornée, puis dix lignes de crayon effacé.\footnote{Ces trois lignes ---
« deux définitions suivant qu'on prend la sommabilité ou la sommabilité
absolue », et l'ordre d'une partie bornée comme borne inférieure des $s$
tels qu'elle soit contenue dans l'enveloppe convexe fermée d'une suite
d'ordre $s$ --- sont la première apparition, dans ce carton, des idées de
\emph{Produits tensoriels topologiques} sur les suites d'ordre $s$ ; on
ne les développe pas.}

\emph{La fonction de comptage.} Soit $(r_\nu)$ une suite croissante de
nombres $> 0$ tendant vers l'infini et $n(r)$ le nombre des $\nu$ avec
$r_\nu \leq r$. Pour $0 \leq a \leq b < \infty$ et $\alpha > 0$, l'intégrale
de Stieltjes $\int_a^b r^{-\alpha}\,dn(r)$ donne
\[
  \sum_{a < r_\nu < b} \frac{1}{r_\nu^\alpha}
  = \frac{n(b-0)}{b^\alpha} - \frac{n(a+0)}{a^\alpha}
  + \alpha\int_a^b \frac{n(r)\,dr}{r^{\alpha+1}} .
\]
\emph{Proposition.} $\sum_\nu r_\nu^{-\alpha} < \infty$ si et seulement si
$\int_0^\infty n(r)\,r^{-\alpha-1}\,dr < \infty$ ; alors $n(r)/r^\alpha \to 0$
et
\[
  \sum_{r_\nu > a} \frac{1}{r_\nu^\alpha} = -\frac{n(a+0)}{a^\alpha}
  + \alpha\int_a^\infty \frac{n(r)\,dr}{r^{\alpha+1}} .
\]
Le sens direct vient de la formule avec $a = 0$ ; réciproquement, si
l'intégrale converge, $n(b_i)/b_i^\alpha$ est borné le long d'une suite
$b_i \to \infty$, d'où la convergence de la série ; et $n(r)/r^\alpha \to 0$,
sinon $n(b_i) - n(b_{i-1}) \geq \varepsilon b_i^\alpha - M b_{i-1}^\alpha$ le
long d'une suite qu'on peut choisir avec $M(b_{i-1}/b_i)^\alpha \leq
\varepsilon/2$, et les tranches $\sum_{b_{i-1} < r_\nu \leq b_i}
r_\nu^{-\alpha}$ resteraient $\geq \varepsilon/2$. Plus généralement, pour
$\lambda$ décroissante, continûment dérivable et tendant vers zéro :
\begin{gather*}
  \sum_\nu \lambda(r_\nu) < \infty \iff \int_0^\infty n(r)\,(-\lambda'(r))\,dr < \infty, \\
  \sum_{r_\nu > a} \lambda(r_\nu) = -\lambda(a)\,n(a+0) + \int_a^\infty n(r)\,(-\lambda'(r))\,dr,
\end{gather*}
et alors $\lambda(r)\,n(r) \to 0$.

Si les $r_\nu$ sont les modules des zéros d'une fonction entière $f$, le
principe de l'argument donne $2\pi\,n(r) = \int_0^{2\pi} \frac{f'}{f}(re^{i\varphi})\,
re^{i\varphi}\,d\varphi$, et, en intégrant contre $r\mu(r)$ puis en prenant
les parties réelles,
\[
  \int_a^b n(r)\,\mu(r)\,dr = \bigl[r\mu(r)\,V(r)\bigr]_a^b - \int_a^b V(r)\,d(r\mu(r)),
  \qquad V(r) = \frac{1}{2\pi}\int_0^{2\pi} \log|f(re^{i\varphi})|\,d\varphi,
\]
pour $\mu$ continûment dérivable ; avec $\mu = -\lambda'$,
\[
  \sum_{a < r_\nu < b} \lambda(r_\nu) = \lambda(b)\,n(b-0) - \lambda(a)\,n(a+0)
  + \bigl[r\mu(r)\,V(r)\bigr]_a^b + \int_a^b V(r)\,d(r\lambda'(r)),
\]
et si $r^2\mu(r) \to 0$ en $0$ et $|f(0)| = 1$, la formule vaut avec $a = 0$
et sans le terme en $a$. Comme $V(r) \leq \log M(r)$, on obtient une
majoration de $\int_a^b n(r)\mu(r)\,dr$ par $M(r)$.\footnote{C'est la
formule de Jensen, $\int_0^r n(t)\,dt/t = V(r) - \log|f(0)|$, intégrée
contre $\mu$. La page 118 écrit $|V(r)| \leq M(r)$ pour $V(r) \leq \log
M(r)$, ce que la transcription note. Ces formules sont exactement celles
que la page 130 invoque comme « majoration classique » pour $\sum
|\lambda_i|^p$, et que la page 146 réutilise.}

\pagerange{119}{121}
\section*{Une fonction holomorphe sur $\ell^1$ (pages 119 à 121)}

Pour $\lambda = (\lambda_i) \in \ell^1$ et $n \geq 0$, soit
\[
  \alpha_n(\lambda) = \sum_{i_1 < \cdots < i_n} \lambda_{i_1} \cdots \lambda_{i_n},
  \qquad |\alpha_n(\lambda)| \leq \frac{\|\lambda\|_1^n}{n!} ;
\]
ce sont les coefficients de la fonction entière de genre zéro $f(z) =
\prod(1 + \lambda_i z)$, et l'on obtient ainsi les coefficients de toutes
les fonctions entières de genre zéro valant $1$ en $0$.

\emph{Proposition.} Pour tout $\lambda \in \ell^1$ la série $S(\lambda) =
\sum_{n \geq 0} n!\,\alpha_n(\lambda)$ converge absolument. La fonction $S$
est holomorphe au voisinage de $0$, de rayon de convergence égal à $1$, et
sa restriction à toute droite est entière.

\emph{Corollaire.} Pour tout $\lambda \in \ell^1$, $\sqrt[n]{|\alpha_n(\lambda)|}
= o(1/n)$ ; autrement dit les coefficients d'une fonction entière de genre
zéro s'écrivent $\alpha_n = \varepsilon_n^n/n!$ avec $\varepsilon_n \to 0$. Car
$\sqrt[n]{n!\,|\alpha_n(\lambda)|} \to 0$ exprime que $\sum n!\,\alpha_n(\lambda)
z^n = \sum n!\,\alpha_n(\lambda z)$ converge pour tout $z$, et
$\sqrt[n]{n!} \sim n/e$ (Stirling).\footnote{En termes classiques : une
fonction de genre zéro est de type minimal pour l'ordre $1$, et le type
d'une fonction d'ordre $1$ est $\limsup (n/e)\,|c_n|^{1/n}$.}

Démonstration. On peut supposer les $\lambda_i > 0$ et décroissants. Soit
$f(z) = \prod(1 + \lambda_i z) = \sum \alpha_n z^n$. Par l'inégalité de Cauchy,
$\alpha_n \leq M(R)/R^n$ avec $M(R) = f(R) = \prod(1 + \lambda_i R)$, et
$\prod_{i \geq n+2}(1 + \lambda_i R) \leq e^{\rho_{n+1}R}$ où $\rho_{n+1} =
\sum_{k \geq n+2}\lambda_k$. Avec $R = (n+1)/\rho_{n+1}$ :
\[
  \alpha_n \leq \prod_{k=1}^{n+1}\Bigl(\frac{\rho_{n+1}}{n+1} + \lambda_k\Bigr) e^{n+1}
  \leq \prod_{k=1}^{n+1}\Bigl(\frac{\rho_k}{k} + \lambda_k\Bigr) e^{n+1} .
\]
Or $\lambda_k = o(1/k)$ (lemme de la page 112) et $\rho_k/k = o(1/k)$, donc
$\rho_k/k + \lambda_k \leq \varepsilon_k/k$ avec $\varepsilon_k \to 0$, et
\[
  n!\,\alpha_n \leq \frac{\varepsilon_1 \cdots \varepsilon_{n+1}}{n+1}\,e^{n+1},
  \qquad
  \sqrt[n]{n!\,\alpha_n} \to 0 ;
\]
la série converge absolument.\footnote{Les indices et exposants des pages
120 et 121 oscillent entre $n$ et $n+1$ ; on a suivi $R = (n+1)/\rho_{n+1}$,
qui donne $n+1$ facteurs.} Ensuite $S$ est holomorphe près de $0$ puisque
$|\alpha_n(\lambda)| \leq \|\lambda\|_1^n$ ; mieux, la forme $n$-linéaire
symétrique associée à $\alpha_n$ est
\[
  \alpha_n(\lambda^{(1)}, \ldots, \lambda^{(n)}) = \frac{1}{n!}
  \sum_{i_1, \ldots, i_n \text{ distincts}} \lambda^{(1)}_{i_1} \cdots \lambda^{(n)}_{i_n},
\]
de norme $\leq 1/n!$, donc $\|n!\,\alpha_n\| \leq 1$ et le rayon $R_0$ de la
série de Taylor de $S$ en $0$ est $\geq 1$. Enfin, pour $a^{(m)} \in \ell^1$
valant $1/m$ sur les $m$ premiers indices et $0$ ailleurs, $\|a^{(m)}\|_1 = 1$
et
\[
  n!\,\alpha_n(a^{(m)}) = \frac{m(m-1)\cdots(m-n+1)}{m^n}
  = \Bigl(1 - \frac{1}{m}\Bigr)\cdots\Bigl(1 - \frac{n-1}{m}\Bigr) \to 1
  \quad (m \to \infty),
\]
donc $S(a^{(m)}) \to \infty$ : $S$ n'est pas bornée sur la boule unité, le
rayon est $\leq 1$, et $R_0 = R_1 = 1$.\footnote{$R_1$ est, dans les termes
d'aujourd'hui, le \emph{rayon de bornitude} de $S$ en $0$ (Dineen) ; la
proposition dit que $S$ est G-holomorphe sur $\ell^1$ tout entier, holomorphe
au sens de Fréchet sur la boule unité ouverte, et non bornée sur cette
boule --- l'exemple type de la distinction entre les deux rayons. On peut
ajouter que $S(\lambda) = \int_0^\infty e^{-t}\prod_i(1 + \lambda_i t)\,dt$,
transformée de Laplace du produit canonique, ce qui rend la convergence
évidente à partir du type minimal ; la page ne le dit pas.}

\pagerange{123}{125}
\section*{Le déterminant d'une somme (pages 123 à 125)}

Trois feuillets d'une main plus lente, antérieurs à la mise au net de la
page 70, dont ils sont la source. La page 123 s'ouvre au milieu d'un
énoncé : pour $u_1, \ldots, u_n$ à trace, en notant $|u|' = |u^{*}|$,
\[
  \|u_1 \wedge \cdots \wedge u_n\|_1^2 \leq \alpha_n(|u_1|, \ldots, |u_n|)\;
  \alpha_n(|u_1|', \ldots, |u_n|'), \qquad
  |\alpha_n(u_1, \ldots, u_n)| \leq \alpha_n(|u_1|, \ldots, |u_n|) .
\]
Avec $u_i = \sum_\alpha \lambda^{(i)}_\alpha\,\bar{a}^{(i)}_\alpha \otimes
b^{(i)}_\alpha$ (Schmidt), on a $|u_i| = \sum \lambda^{(i)}_\alpha\,
\bar{a}^{(i)}_\alpha \otimes a^{(i)}_\alpha$, $|u_i|' = \sum \lambda^{(i)}_\alpha\,
\bar{b}^{(i)}_\alpha \otimes b^{(i)}_\alpha$, et
\[
  u_1 \wedge \cdots \wedge u_n = \frac{1}{n!}\sum_{\alpha_1, \ldots, \alpha_n}
  \lambda^{(1)}_{\alpha_1} \cdots \lambda^{(n)}_{\alpha_n}\,
  \overline{(a^{(1)}_{\alpha_1} \wedge \cdots \wedge a^{(n)}_{\alpha_n})} \otimes
  (b^{(1)}_{\alpha_1} \wedge \cdots \wedge b^{(n)}_{\alpha_n}),
\]
d'où $\|u_1 \wedge \cdots \wedge u_n\|_1 \leq \frac{1}{n!}\sum \prod
\lambda\,\|a_{\alpha_1} \wedge \cdots\|\,\|b_{\alpha_1} \wedge \cdots\|$,
tandis que $\alpha_n(|u_1|, \ldots, |u_n|) = \frac{1}{n!}\sum \prod\lambda\,
\|a_{\alpha_1} \wedge \cdots\|^2$ et de même avec les $b$ ; c'est
Cauchy--Schwarz.

\emph{Corollaire 1.} $|\alpha_n(u_1, \ldots, u_n)| \leq \|u_1 \wedge \cdots
\wedge u_n\|_1 \leq \frac12\bigl(\alpha_n(|u_1|, \ldots, |u_n|) +
\alpha_n(|u_1|', \ldots, |u_n|')\bigr)$.
\emph{Corollaire 2.} $|\alpha_n(A+B)|^2 \leq \alpha_n(|A| + |B|)\;
\alpha_n(|A|' + |B|')$ : on développe $\alpha_n(A+B) = \sum_k \binom{n}{k}
\alpha_n(A, \ldots, A, B, \ldots, B)$, on majore chaque terme par la racine
du produit des deux $\alpha_n$ correspondants, et Cauchy--Schwarz sur la
somme en $k$ refait les deux $\alpha_n(|A| + |B|)$ et $\alpha_n(|A|' +
|B|')$.
\emph{Corollaire 3.} $|\det(1 + A + B)| \leq \frac12\bigl(\det(1 + |A| + |B|)
+ \det(1 + |A|' + |B|')\bigr)$.
\emph{Corollaire 4.} $|\det(1 + A + B)| \leq \det(1 + |A|)\,\det(1 + |B|)$.
Il suffit, par le corollaire 3 et $\det(1 + |A|) = \det(1 + |A|')$, de le
prouver pour $A = H$ et $B = K$ hermitiens positifs ; alors $1 + H + K =
(1 + H)\bigl(1 + (1+H)^{-1}K\bigr)$, donc $\det(1 + H + K) = \det(1 + H)\,
\det(1 + H'K)$ avec $H' = (1+H)^{-1}$ hermitien positif de norme $\leq 1$, et
il reste $\det(1 + H'K) \leq \det(1 + K)$, soit $\alpha_n(H'K) \leq
\alpha_n(K)$ pour tout $n$, ce qui résulte de $\Lambda^n(H'K) =
(\Lambda^n H')(\Lambda^n K)$ et $\|\Lambda^n(H'K)\|_1 \leq \|\Lambda^n H'\|\,
\|\Lambda^n K\|_1 \leq \|K\|_1$-type, avec $\|\Lambda^n H'\| \leq \|H'\|^n
\leq 1$.\footnote{La page 124 écrit $\det(1 + H + K) = \det(1 + L)(1 +
K/(1+L))$ avec une lettre $R$ ou $H$ que la suite note $L$ et $H = 1/(1+L)$ ;
on a renommé pour ne pas donner deux sens à $H$. Le corollaire 4 vaut sans
$A$, $B$ hermitiens parce que $|\det(1+A+B)| \leq \det(1 + |A+B|)$ n'est
pas utilisé : c'est le corollaire 3 qui fait le passage.}
\emph{Corollaire 5.} $\det(1 + |A + B|) \leq \det(1 + |A|)\,\det(1 + |B|)$.
Il contient le précédent, puisque $|\det(1 + A + B)| \leq \det(1 + |A + B|)$,
et il en résulte : $|A + B| = U(A + B)$ pour une isométrie partielle $U$,
et par le corollaire 4, $\det(1 + UA + UB) \leq \det(1 + |UA|)\det(1 + |UB|)$
avec $|UA| = |A|$ sur le support, donc $\det(1 + |UA|) \leq \det(1 +
|A|)$.\footnote{Le feuillet collé de la page 125 écrit « $U$ unitaire » et
$|UA| = U|A|U^{-1}$, ce qui suppose $U$ unitaire ; pour une isométrie
partielle $|UA| = (A^{*}U^{*}UA)^{1/2} \leq |A|$, ce qui suffit. Le
corollaire 5 est l'inégalité de Seiler--Simon (1975) déjà rencontrée à la
page 71 ; les deux démonstrations sont différentes.}

\pagerange{126}{131}
\section*{Le théorème sur les classes $\mathcal{S}_p(H)$, $0 < p \leq \infty$ (pages 126 à 131)}

Quatorze pages d'un seul tenant, l'état ancien de la théorie, avant le
lemme de Weyl : un énoncé en trois points et sa démonstration, numérotée
(1) à (14).

\emph{Théorème.} 1) Pour $0 < p \leq \infty$, $\mathcal{S}_p(H)$ est un idéal
bilatère de $L(H)$, stable par $A \mapsto A^{*}$ :
\[
  \|A\|_p = \|A^{*}\|_p, \qquad \|BA\|_p \leq \|B\|\,\|A\|_p, \qquad
  \|AB\|_p \leq \|B\|\,\|A\|_p \qquad (A \in \mathcal{S}_p(H),\ B \in L(H)) .
\]
Si $1 \leq p \leq \infty$, $\|\cdot\|_p$ est une norme sur $\mathcal{S}_p(H)$ :
$\|A + B\|_p \leq \|A\|_p + \|B\|_p$, et $\mathcal{S}_p(H)$ est complet. Si
$0 < p \leq 1$, $S_p(A) = \|A\|_p^p$ vérifie
\[
  S_p(A + B) \leq S_p(A) + S_p(B),
\]
et pour la distance invariante $S_p(A - B)$, $\mathcal{S}_p(H)$ est un espace
vectoriel topologique métrisable complet, non localement convexe si $p < 1$
et $\dim H = \infty$.
2) Si $A \in \mathcal{S}_p(H)$, $0 < p < \infty$, la suite $(\lambda_i)$ de
ses valeurs propres est de puissance $p$-ième sommable et
\[
  \sum |\lambda_i|^p \leq S_p(A) .
\]
$S_p$ et $\|\cdot\|_p$ sont croissantes sur les hermitiens positifs : $0 \leq
H \leq K$ entraîne $S_p(H) \leq S_p(K)$.
3) Si $0 < p, q, r \leq \infty$ et $\frac1p + \frac1q \geq \frac1r$, alors
$A \in \mathcal{S}_p(H)$, $B \in \mathcal{S}_q(H)$ entraînent $AB \in
\mathcal{S}_r(H)$ et
\[
  \|AB\|_r \leq \|A\|_p\,\|B\|_q .
\]
\footnote{Les trois points sont aujourd'hui classiques : pour $p \geq 1$,
Schatten et von Neumann ; pour $p < 1$, la sous-additivité de $S_p$ et
l'inégalité de Hölder sont attribuées à McCarthy (1967) et Rotfel'd (1967),
l'inégalité de Weyl du point 2 à Weyl (1949). Tout est ici démontré, par
une méthode qui n'est celle d'aucun de ces auteurs. Les quatre lignes de
crayon en tête de la page 126 --- « le déterminant de Fredholm »,
« notation $\mathcal{L}_0$ », « sont connues » --- et la marge « ② Dual de
$\mathcal{L}^p$ ; réflexivité » annoncent un second théorème qui n'est pas
ici.}

\emph{Démonstration de 1).} La première formule est évidente, $|A|$ et
$|A^{*}|$ ayant les mêmes valeurs propres. La seconde est un cas particulier
de 3), mais se prouve directement : si $(k_i)$ et $(h_i)$ sont les valeurs
propres décroissantes de $|BA|$ et $|A|$, $k_i \leq \|B\|\,h_i$ pour tout $i$.
L'inégalité triangulaire pour $p \geq 1$ se déduit de 3), démontré
indépendamment : si $\frac1p + \frac1{p'} = 1$ et $\|B\|_{p'} \leq 1$, alors
$\|AB\|_1 \leq \|A\|_p$, donc $|\operatorname{Tr} AB| \leq \|A\|_p$ ; et
l'inégalité $\sup_{\|B\|_{p'} \leq 1} |\operatorname{Tr} AB| \leq \|A\|_p$
est une égalité, car avec $U$ isométrie partielle telle que $UA = |A|$ et,
par Hölder classique sur la décomposition spectrale de $|A|$, un hermitien
positif $H \in \mathcal{S}_{p'}$ de norme $\leq 1$ avec $\operatorname{Tr}(H|A|)
= \|A\|_p$, l'opérateur $B = HU$ a $\|B\|_{p'} \leq 1$ et $\operatorname{Tr}
AB = \operatorname{Tr}(HUA) = \|A\|_p$. Ainsi
\[
  \|A\|_p = \sup_{\|B\|_{p'} \leq 1} |\operatorname{Tr} AB| \qquad (1 \leq p \leq \infty),
\]
sup de semi-normes, donc semi-norme.\footnote{La page 128 écrit
$\operatorname{Tr}(HUA)$ puis « $B = UA$ » ou « $HU$ » ; c'est $HU$ que
l'argument demande, et $\operatorname{Tr}(A\,HU) = \operatorname{Tr}(HUA)$
par cyclicité.} Une telle démonstration n'est plus praticable pour $p < 1$,
et l'on prouve la sous-additivité de $S_p$ « par voie analytique », grâce
au

\emph{Lemme 1.} Soit $(\alpha_i)$ une suite de nombres $\geq 0$, $0 < p <
1$, et $f(z) = \prod_i (1 + \alpha_i z)$. Alors
\[
  \sum_i \alpha_i^p = \frac{p\sin\pi p}{\pi}\int_0^\infty \frac{\log f(v)}{v^{1+p}}\,dv,
\]
les deux membres étant simultanément finis ou infinis.\footnote{Non démontré
sur la page ; c'est le calcul $\int_0^\infty \log(1 + \alpha v)\,v^{-1-p}\,dv
= \alpha^p\int_0^\infty \log(1+s)\,s^{-1-p}\,ds = \frac{\alpha^p}{p}
\int_0^\infty \frac{ds}{(1+s)s^p} = \frac{\alpha^p}{p}\,\frac{\pi}{\sin\pi p}$,
sommé sur $i$ par convergence monotone.}

\emph{Corollaire.} Pour $A \in \mathcal{S}_1(H)$, soit $M_A(v) = \sup_{|z| =
v} |\det(1 + zA)|$. Pour $0 < p < 1$,
\[
  S_p(A) = \frac{p\sin\pi p}{\pi}\int_0^\infty \frac{\log M_{|A|}(v)}{v^{1+p}}\,dv,
\]
en prenant $f(z) = \det(1 + z|A|)$, dont $M_{|A|}(v) = f(v)$.

Pour la sous-additivité : $S_p(A + B) = S_p(|A + B|)$ est donné par cette
intégrale, et $M_{|A+B|}(v) = \det(1 + v|A + B|) \leq \det(1 + v|A|)\det(1 +
v|B|) = M_{|A|}(v)\,M_{|B|}(v)$ par le corollaire 5 de la page 124 ; on
intègre et l'on réapplique le corollaire à $A$ et $B$.\footnote{C'est,
mot pour mot, la déduction de l'inégalité de McCarthy--Rotfel'd à partir de
l'inégalité de Seiler--Simon --- dans l'ordre inverse de la littérature, qui
déduit la seconde de la première. Pour $A$, $B \notin \mathcal{S}_1$ le
corollaire s'applique par troncature.} Les assertions topologiques sont des
conséquences faciles des formules.

\emph{Démonstration de 2).} « Je ne connais qu'une méthode analytique. »
Pour $A$ compact et $0 < p < \infty$, posons $s_p(A) = \sum |\lambda_i|^p$ ;
on a $s_p(A) = s_{p/n}(A^n)$ pour tout entier $n > 0$. Si $p < 1$, la
majoration classique des inverses des modules des zéros de la fonction
entière $\det(1 - zA)$ --- ce sont les $1/\lambda_i$ --- par la formule de
Jensen donne
\[
  s_p(A) \leq p^2\int_0^\infty \frac{\log M_A(r)}{r^{p+1}}\,dr
  \leq p^2\int_0^\infty \frac{\log M_{|A|}(v)}{v^{p+1}}\,dv
  = \frac{p\pi}{\sin p\pi}\,S_p(A),
\]
puisque $M_A \leq M_{|A|}$ par $|\det(1 + zA)| \leq \det(1 + |z||A|)$, et par
le corollaire.\footnote{La « majoration classique » est celle des pages 115
à 118 : $\sum r_\nu^{-p} = p\int n(r)\,r^{-p-1}\,dr$, et $\int_0^r n(t)\,dt/t
\leq \log M(r)$ par Jensen, d'où le facteur $p^2$ après une seconde
intégration par parties.} Le facteur $\frac{p\pi}{\sin p\pi}$ est $> 1$ mais
tend vers $1$ quand $p \to 0$. Pour $p$ quelconque, soit $n$ avec $p/n < 1$ :
\[
  s_p(A) = s_{p/n}(A^n) \leq \frac{\frac{p}{n}\pi}{\sin\frac{p}{n}\pi}\,
  S_{p/n}(A^n) \leq \frac{\frac{p}{n}\pi}{\sin\frac{p}{n}\pi}\,S_p(A),
\]
car $\|A^n\|_{p/n} \leq \|A\|_p^n$ par 3) généralisé à $n$ facteurs, soit
$S_{p/n}(A^n) \leq S_p(A)$ ; et $n \to \infty$ donne l'inégalité. La
monotonie sur les positifs vient de ce que $0 \leq H \leq K$ entraîne
$\lambda_i(H) \leq \lambda_i(K)$ pour tout $i$, par la caractérisation
extrémale de Fischer.

\pagerange{132}{135}
\section*{Hölder pour les exposants dyadiques, et le lemme (pages 132 à 135)}

\emph{Démonstration de 3).} C'est la plus difficile. On commence par une
preuve directe et très simple du cas $p, q \geq 1$, $r = 1$ --- qui
s'applique sans changement à une trace au sens de von Neumann sur une
algèbre auto-adjointe, les opérateurs « de puissance $p$-ième sommable »
étant les $A$ avec $\tau(|A|^p) < \infty$.\footnote{Ce que la page 33
réalise en trois lignes, dans le second état.} On prouve d'abord, pour
$2^n$ opérateurs compacts $A_1, \ldots, A_{2^n}$, $n \geq 1$,
\[
  \|A_1 A_2 \cdots A_{2^n}\|_1 \leq \|A_1\|_{2^n}\,\|A_2\|_{2^n} \cdots \|A_{2^n}\|_{2^n},
\]
par récurrence sur $n$ : pour $n = 1$ c'est $\|A_1A_2\|_1 \leq
\|A_1\|_2\|A_2\|_2$ ; ensuite, groupant les facteurs deux par deux et
appliquant l'hypothèse aux $2^{n-1}$ produits, on est ramené à $\|AB\|_{2^{n-1}}
\leq \|A\|_{2^n}\|B\|_{2^n}$. Or
\[
  \|AB\|_{2^{n-1}}^{2^{n-1}} = \operatorname{Tr}\bigl[(ABB^{*}A^{*})^{2^{n-2}}\bigr]
  = \operatorname{Tr} C = \operatorname{Tr} C', \qquad
  C' = (A^{*}A)(BB^{*})(A^{*}A)\cdots(BB^{*}),
\]
$C'$ ayant $2^{n-1}$ facteurs, obtenu par cyclicité ; et par l'hypothèse
de récurrence, $|\operatorname{Tr} C'| \leq \|C'\|_1 \leq \bigl(\|A^{*}A\|_{2^{n-1}}
\|BB^{*}\|_{2^{n-1}}\bigr)^{2^{n-2}} = (\|A\|_{2^n}\|B\|_{2^n})^{2^{n-1}}$,
puisque $\|A^{*}A\|_{2^{n-1}} = \|A\|_{2^n}^2$. Une note de bas de page
observe que ce résultat suffirait à démontrer directement la formule 3) dans
le cas général.

Soient $\frac1p + \frac1q = 1$, $A \in \mathcal{S}_p$, $B \in \mathcal{S}_q$,
qu'on peut supposer de rang fini. Écrivons $A = UH$, $B = KV$ avec $U$, $V$
partiellement isométriques, et pour $\alpha + \beta = 2^n$, $h = H^{1/\alpha}$,
$k = K^{1/\beta}$ : $AB = Uh^\alpha k^\beta V$ et
\[
  \|AB\|_1 \leq \|h^\alpha k^\beta\|_1 \leq \|h\|_{2^n}^\alpha\,\|k\|_{2^n}^\beta
  = (\operatorname{Tr} H^{\lambda})^{1/\lambda}(\operatorname{Tr} K^{\mu})^{1/\mu},
  \qquad \lambda = \frac{2^n}{\alpha},\ \mu = \frac{2^n}{\beta},
\]
$\frac1\lambda + \frac1\mu = 1$. L'inégalité vaut pour tous les couples
d'exposants conjugués dont les inverses sont des multiples de $2^{-n}$, et
par densité et continuité de $p \mapsto \|H\|_p$ pour tous : $\|AB\|_1 \leq
\|A\|_p\|B\|_q$. Une note interlinéaire ajoute la généralisation $\|A_1
\cdots A_n\|_1 \leq \prod \|A_i\|_{p_i}$ si $\sum 1/p_i = 1$, par le même
argument avec $n$ exposants dyadiques.\footnote{La page 134 invoque « Lemme
1 » pour la formule à $2^n$ facteurs, qui n'a pas de numéro ; et elle écrit
$(\operatorname{Tr} H^{p'})^{1/p'}$ pour $K$. On a corrigé les lettres. La
réduction au rang fini n'est pas justifiée ; elle se fait par troncature
spectrale, $\|AB\|_1$ étant limite croissante des $\|P_nAB Q_n\|_1$.} On en
conclut, pour $p \geq 1$, la formule de dualité $\|A\|_p = \sup_{\|B\|_{p'}
\leq 1} |\operatorname{Tr} AB|$ et l'inégalité triangulaire.

\emph{Lemme 2.} Soient $A$, $B \in \mathcal{S}_{2m}(H)$, $H = A^{*}A$, $K =
B^{*}B$, $L = (AB)^{*}(AB)$, et $m$, $n$ des entiers $\geq 1$. Alors $L^m$,
$H^m$, $K^m$ sont à trace et
\[
  \alpha_n(L^m) \leq \alpha_n(H^m)\,\alpha_n(K^m) .
\]
En remplaçant $A$, $B$ par $\Lambda^n A$, $\Lambda^n B$ --- ce qui remplace
$H$, $K$, $L$ par $\Lambda^n H$, $\Lambda^n K$, $\Lambda^n L$ et $\alpha_n$
par $\operatorname{Tr}$ --- on est ramené à $n = 1$ :
\[
  \operatorname{Tr}\bigl[(ABB^{*}A^{*})^m\bigr] \leq
  \bigl(\|A^{*}A\|_m\bigr)^m\bigl(\|B^{*}B\|_m\bigr)^m
  = \operatorname{Tr}(H^m)\,\operatorname{Tr}(K^m) .
\]
Or le premier membre est la trace du produit de $2m$ facteurs $(A^{*}A)(BB^{*})
\cdots(BB^{*})$, majorée par sa norme trace, donc par $\bigl(\|A^{*}A\|_{2m}
\|BB^{*}\|_{2m}\bigr)^m$ par Hölder à $2m$ facteurs, et $\|A^{*}A\|_{2m} \leq
\|A^{*}A\|_m$, $\|BB^{*}\|_{2m} = \|B^{*}B\|_{2m} \leq \|B^{*}B\|_m$.

\pagerange{136}{139}
\section*{Hölder en toute généralité (pages 136 à 139)}

Soient $A \in \mathcal{S}_p$, $B \in \mathcal{S}_q$ avec $\frac1p + \frac1q
\geq \frac1r$, $p$ et $q$ finis --- sinon c'est 1) --- et même $\frac1p +
\frac1q = \frac1r$, puisque $\|A\|_r$ décroît avec $r$. Soit $m$ un entier
avec $p, q \leq 2m$, et $L = (AB)^{*}(AB)$, $H = A^{*}A$, $K = B^{*}B$ ; alors
$r \leq m$ et $S_r(AB) = S_{r/2}(L) = S_{r/2m}(L^m)$, de sorte que, avec
$\lambda(x) = \pi/\sin\pi x$ et le corollaire du lemme 1,
\[
  S_r(AB) = \frac{r/2m}{\lambda(r/2m)}\int_0^\infty
  \frac{\log M_{L^m}(\rho)}{\rho^{1 + r/2m}}\,d\rho .
\]
Or $M_{L^m}(\rho) = \sum_n \alpha_n(L^m)\rho^n \leq \sum_n \alpha_n(H^m)\,
\alpha_n(K^m)\,\rho^n$ par le lemme 2, et avec $\frac{r}{p} + \frac{r}{q} = 1$,
\[
  \sum_n \alpha_n(H^m)\alpha_n(K^m)\rho^{n\frac rp + n\frac rq}
  \leq \Bigl(\sum_n \alpha_n(H^m)\rho^{n\frac rp}\Bigr)\Bigl(\sum_n \alpha_n(K^m)\rho^{n\frac rq}\Bigr)
  = M_{H^m}(\rho^{r/p})\,M_{K^m}(\rho^{r/q}),
\]
d'où
\[
  \frac{\lambda(r/2m)}{r/2m}\,S_r(AB) \leq
  \int_0^\infty \frac{\log M_{H^m}(\rho^{r/p})}{\rho^{1 + r/2m}}\,d\rho
  + \int_0^\infty \frac{\log M_{K^m}(\rho^{r/q})}{\rho^{1 + r/2m}}\,d\rho .
\]
Le changement de variable $\sigma = \rho^{r/p}$, $\frac{d\rho}{\rho} =
\frac{p}{r}\frac{d\sigma}{\sigma}$, transforme la première intégrale en
$\frac{p}{r}\int_0^\infty \log M_{H^m}(\sigma)\,\sigma^{-1-p/2m}\,d\sigma =
\frac{p}{r}\,\frac{\lambda(p/2m)}{p/2m}\,S_{p/2m}(H^m) = \frac{\lambda(p/2m)}{r/2m}
\,S_p(A)$, et de même pour la seconde ; donc
\begin{gather*}
  \lambda(r/2m)\,S_r(AB) \leq \lambda(p/2m)\,S_p(A) + \lambda(q/2m)\,S_q(B),
  \qquad\text{soit} \\
  S_r(AB) \leq \frac{\sin\pi\frac{r}{2m}}{\sin\pi\frac{p}{2m}}\,S_p(A)
  + \frac{\sin\pi\frac{r}{2m}}{\sin\pi\frac{q}{2m}}\,S_q(B),
\end{gather*}
ce qui donne déjà la finitude de $S_r(AB)$. Faisant $m \to \infty$ :
\[
  S_r(AB) \leq \frac{r}{p}\,S_p(A) + \frac{r}{q}\,S_q(B) .
\]
« C'est une inégalité remarquable, mais n'est pas celle que nous
cherchons. » On y remplace $A$ par $\lambda A$ et $B$ par $\mu B$, $\lambda,
\mu \geq 0$ :
\[
  (\lambda\mu)^r S_r(AB) \leq \frac{\lambda^p}{p}\,S_p(A) + \frac{\mu^q}{q}\,S_q(B),
\]
et avec $a = S_p(A)$, $b = S_q(B)$, $c = S_r(AB)$, $x = \lambda^p$, $y =
\mu^q$, $\alpha = r/p$, $\beta = r/q$, $\alpha + \beta = 1$ :
\[
  c \leq \inf_{x, y \geq 0} \frac{\alpha a x + \beta b y}{x^\alpha y^\beta}
  = \inf_{x^\alpha y^\beta = 1} (\alpha a x + \beta b y) .
\]
L'infimum est atteint, et au point critique $ax = by$, d'où $x = (b/a)^\beta$
et le minimum $ax = a^\alpha b^\beta$. Donc $S_r(AB) \leq S_p(A)^{r/p}
S_q(B)^{r/q}$, c'est-à-dire $\|AB\|_r \leq \|A\|_p\|B\|_q$, ce qui achève la
démonstration.\footnote{C'est l'inégalité de Young pondérée, $\alpha a x +
\beta b y \geq a^\alpha b^\beta x^\alpha y^\beta$, obtenue par un multiplicateur
de Lagrange. La note interlinéaire de la page 137 relève que l'inégalité
intermédiaire donne déjà $S_r(AB) < \infty$, c'est-à-dire l'appartenance
$AB \in \mathcal{S}_r$.}

\pagerange{141}{142}
\section*{Convergence dans $\mathcal{S}_p$ (pages 141 et 142)}

\emph{Proposition.} Pour qu'une suite $(H_i)$ d'hermitiens de
$\mathcal{S}_p(H)$, $0 < p < \infty$, converge vers $H_0 \in \mathcal{S}_p(H)$,
il faut et il suffit que
1°) $\|H_i\|_p \to \|H_0\|_p$, ou seulement $\limsup \|H_i\|_p \leq
\|H_0\|_p$ ;
2°) pour tout $\lambda \neq 0$ du spectre de $H_0$ et tout $\alpha > 0$ tel
que $[\lambda - \alpha, \lambda + \alpha]$ ne contienne pas d'autre valeur
spectrale de $H_0$ et que $\lambda \pm \alpha$ ne soient dans le spectre
d'aucun $H_i$, le projecteur spectral $h_i$ de $H_i$ sur cet intervalle
finit par avoir la dimension de celui de $H_0$, $h_0$, et $h_i \to h_0$ ---
la topologie n'a pas à être précisée, toutes coïncidant sur les projecteurs
de rang fini borné.\footnote{L'énoncé est correct : nécessité par la
convergence en norme et la continuité des projecteurs de Riesz ;
suffisance parce que, pour $\delta > 0$, les valeurs propres de $H_i$ dans
les intervalles autour des $\lambda$ avec $|\lambda| > \delta$ convergent
vers celles de $H_0$ avec leurs vecteurs propres, tandis que 1°) force la
somme des $|\mu|^p$ sur les autres valeurs propres de $H_i$ à tendre vers
$\sum_{|\lambda| \leq \delta} |\lambda|^p$, petite ; la sous-additivité de
$S_p$ conclut. La page écrit « tout $\lambda$ du spectre » ; pour $\lambda =
0$ l'intervalle n'est pas isolé, et c'est 1°) qui en tient lieu.}

\emph{Corollaire 1.} $H_i \to H_0$ dans $\mathcal{S}_p$ si et seulement si
$\|H_i\|_p \to \|H_0\|_p$ et $H_i \to H_0$ en norme d'opérateur.
\emph{Corollaire 2.} $A \mapsto |A|^r$ est continue de $\mathcal{S}_p$ dans
$\mathcal{S}_{p/r}$, comme composée de $A \mapsto A^{*}A$, continue de
$\mathcal{S}_p$ dans $\mathcal{S}_{p/2}$ par Hölder, et de $H \mapsto H^{r/2}$
sur le cône positif de $\mathcal{S}_{p/2}$, qui vérifie les deux conditions
de la proposition : la première trivialement, la seconde parce que les
projecteurs spectraux de $H^{r/2}$ sont ceux de $H$.

\emph{Remarque 1.} Si $1 \leq p \leq 2$, une suite d'hermitiens de
$\mathcal{S}_p$ converge dès qu'elle converge faiblement avec convergence
des normes. Est-ce vrai pour tout $p < \infty$ ? Uniforme convexité de
$\mathcal{S}_p$ ?\footnote{Oui, pour $1 < p < \infty$ : $\mathcal{S}_p$ est
uniformément convexe (McCarthy 1967, inégalités de Clarkson--McCarthy),
donc a la propriété de Kadec--Klee que demande la remarque ; pour $p = 1$,
$\mathcal{S}_1$ a encore la propriété de Kadec--Klee (Arazy 1981) sans être
uniformément convexe.}
\emph{Remarque 2.} Pour $H$, $K$ hermitiens positifs, $S_p(|H^r - K^r|^{1/r})$
est-elle fonction croissante de $r$ ? Il suffirait de le savoir pour $r <
1$, et cela donnerait
\[
  \bigl\|\,|H^r - K^r|^{1/r}\,\bigr\| \leq \|H - K\| ,
\]
et une précision notable du corollaire 1.\footnote{La réponse est oui pour
$0 < r \leq 1$ : $\|\,|H^r - K^r|\,\|_{p/r} \leq \|\,|H - K|^r\,\|_{p/r}$, soit
$S_{p/r}(H^r - K^r) \leq S_p(H - K)$, est l'inégalité de Birman, Koplienko et
Solomyak (1975), valable pour toute norme unitairement invariante ; la
monotonie en $r$ sur $]0, 1]$ s'en déduit en l'appliquant à $H^{r'}$,
$K^{r'}$. L'inégalité en norme d'opérateur, $\|H^r - K^r\| \leq \|H -
K\|^r$, est celle de Powers et Størmer (1970) pour $r = 1/2$ et de
Kittaneh en général.}

\pagerange{143}{147}
\section*{Les inégalités de Ky Fan, et les petites valeurs singulières (pages 143 à 147)}

\emph{Les deux inégalités.} Pour $u : E \to F$, $v : F \to G$ compacts entre
espaces de Hilbert,
\[
  s_{m+n-1}(vu) \leq s_m(v)\,s_n(u), \qquad s_{m+n-1}(u+v) \leq s_m(u) + s_n(v) .
\]
Par le minimax, $s_k(w) = \inf_{\dim E_{k-1} = k-1}\sup_{x \perp E_{k-1}, \|x\|
\leq 1}\|wx\|$. Soit $E_{n-1} \subset E$ réalisant $s_n(u)$ et $F_{m-1}
\subset F$ réalisant $s_m(v)$, et $H = u^{*}(F_{m-1})$. Pour $x \perp H$ et
$x \perp E_{n-1}$ on a $ux \perp F_{m-1}$, donc $\|vux\| \leq s_m(v)\|ux\|
\leq s_m(v)s_n(u)\|x\|$ ; et $H + E_{n-1}$ est de dimension $\leq m + n - 2$,
d'où la première inégalité. La seconde : pour $x \perp E_{n-1} + F_{m-1}$
--- ici $v : E \to F$ aussi, et $F_{m-1} \subset E$ réalise $s_m(v)$ ---,
$\|(u+v)x\| \leq (s_n(u) + s_m(v))\|x\|$.\footnote{La page 143 écrit
« $\operatorname{codim} H \wedge E_{n-1} \geq \operatorname{cod} H +
\operatorname{cod} E_{n-1}$ » là où l'argument demande $\leq$, et un $\bar{u}$
pour $u^{*}$ ; la transcription note les deux. Ce sont les inégalités de
Ky Fan (1951), que la page 38 cite.}

Sous un trait, une question au crayon : peut-on avoir $\limsup
|\lambda_i(u)|/s_i(u) = +\infty$ ? Une majoration $|\lambda_i(u)| \leq
s_{i/2}(u)$ « ou analogue » lui paraît invraisemblable.\footnote{Elle l'est.
Sur $\mathbf{C}^n$, le décalage cyclique pondéré $e_i \mapsto w_i e_{i+1}$,
$e_n \mapsto w_n e_1$, a pour valeurs singulières les $w_i$ et pour valeurs
propres les $n$ racines $n$-ièmes de $w_1 \cdots w_n$, toutes de module
$(w_1 \cdots w_n)^{1/n}$ ; avec $w = (1, \varepsilon, \varepsilon)$, $|\lambda_3|
= \varepsilon^{2/3} > s_2 = \varepsilon$, ce qui exclut $|\lambda_i| \leq
s_{\lceil i/2\rceil}$ ; et une somme directe de tels blocs, de tailles
croissantes et d'échelles décroissantes, donne $\limsup |\lambda_i|/s_i =
+\infty$ avec $u$ compact. Ce que Weyl laisse, et qui est optimal, ce sont
les majorations des sommes et des produits partiels.}

Les pages 144 et 145, au crayon pâle, reprennent la proposition pour un
espace de Banach $E$ --- si $u_i \to u$ dans $L_0(E)$, les spectres et les
projecteurs spectraux convergent --- et le corollaire 2 sous la forme
$u \mapsto |u|^r$ continue de $\mathcal{S}_p(E,F)$ dans $\mathcal{S}_{p/r}(E)$,
avec une question : $\sigma \mapsto \sigma^{r/2}$ est-elle continue de
$L(E)^{+}$ dans $L(E)$, et avec quelles inégalités précises ?\footnote{Oui :
c'est la continuité de $\sigma \mapsto \sigma^{s}$ pour $0 < s < 1$ sur le cône
positif, en norme d'opérateur, conséquence de l'inégalité de Powers--Størmer
citée à la remarque 2 ; en 1953 le calcul fonctionnel continu la donne
sur les bornés.} Un \emph{corollaire 2} affirme que pour $0 < p \leq 1$
une suite d'hermitiens $u_i \to u$ converge dans $\mathcal{S}_p(E)$ dès
qu'elle converge faiblement et que les normes convergent ; la page 145
demande dans quels cas $\mathcal{S}_p(E,F)$ est uniformément convexe pour
$0 < p < \infty$ --- « pas d'importance » ---, ramène la continuité de
$H \mapsto H^r$ au cas $0 < r < 1$, et reformule la remarque 2 : prouver que
$S_{p/r}(H^r - K^r) = S_p(|H^r - K^r|^{1/r})$ croît avec $r$, pour $p$
quelconque ou déjà pour $p = 1$.\footnote{Ce corollaire 2 pour $p \leq 1$ est
douteux tel quel : la convergence faible et celle des quasi-normes ne
donnent pas la convergence en norme d'opérateur, et l'énoncé de la page 141
demande cette dernière. Pour $p < 1$, $\mathcal{S}_p$ n'est pas localement
convexe et la question de la convexité uniforme n'a pas de sens ; la page
le dit à sa façon.}

\emph{Théorème.} Si $A_i \to A$ dans $\mathcal{S}_p(H)$, alors la suite des
valeurs singulières $\hat{A}_i = (s_\nu(A_i))_\nu$ converge vers $\hat{A}$
dans $\ell^p$.\footnote{Ni $\hat{A}$ ni l'espace $\Sigma_p$ de la page ne
sont définis ; on les lit ainsi. Pour $p \geq 1$ c'est l'inégalité de
Mirsky (1960), $\|(s_\nu(A) - s_\nu(B))\|_N \leq \|A - B\|_N$ pour toute
norme unitairement invariante ; pour $p < 1$ l'inégalité $\sum |s_\nu(A) -
s_\nu(B)|^p \leq S_p(A - B)$ est vraie aussi, et le lemme ci-dessous est ce
qui en tient lieu ici.} Cela résulte du

\emph{Lemme.} Pour tout $\varepsilon > 0$ il existe $\rho > 0$ et $i_0$ tels
que $\sum_{s_\nu(A_i) \leq \rho} s_\nu(A_i)^p \leq \varepsilon$ pour $i \geq
i_0$.

Soit $\omega$ la fonction convexe égale à $e^{pt}$ pour $t \leq \log\rho$ et
à sa tangente au-delà, $\omega(t) = \rho^p\bigl(1 + p(t - \log\rho)\bigr)$ ;
la majorisation de Weyl donne $\sum_\nu \omega(\log s_\nu(A_i)) \leq \sum_\nu
\omega(\log s_\nu(|A_i|))$, ce qui permet de supposer $A_i \geq 0$, et
$\sum_\nu \omega(\log s_\nu)$ se lit
\[
  \sum_{s_\nu \leq \rho} s_\nu^p + \rho^p\sum_{s_\nu > \rho} \log\Bigl(e\,\frac{s_\nu^p}{\rho^p}\Bigr) .
\]
\footnote{« $\log e(x)^p$ » est sa façon d'écrire $\log(e\,x^p)$, ce que la
première expression de $\omega$ confirme. La réduction à $A_i \geq 0$ est
la ligne la moins lisible de la page ; l'inégalité entre les sommes pour
$A_i$ et $|A_i|$ est en fait une égalité, $s_\nu(A_i) = s_\nu(|A_i|)$, et le
lemme se démontre tel quel pour $A_i$ quelconques.} Il suffit alors de
montrer :
a) $\rho^p\sum_{s_\nu > \rho} \log(e\,s_\nu^p/\rho^p) \to 0$ quand $\rho \to 0$,
pour la suite fixe $(s_\nu) = \hat{A}$ ; b) $\sum_{s_\nu(A_i) \leq \rho}
s_\nu(A_i)^p \to \sum_{s_\nu \leq \rho} s_\nu^p$ quand $i \to \infty$, à $\rho$
fixé.

\emph{Preuve de a).} Soit $n(r)$ le nombre des $\nu$ avec $s_\nu > r$ ; par
l'intégration de Stieltjes des pages 115 et 116, avec $g(r) = \log(e\,r^p/
\rho^p)$, $g(\rho) = 1$, $g'(r) = p/r$ :
\[
  \rho^p\sum_{s_\nu > \rho} \log\Bigl(e\,\frac{s_\nu^p}{\rho^p}\Bigr)
  = \rho^p\Bigl[-n(r)\,g(r)\Bigr]_{\rho+0}^{\infty} + p\rho^p\int_\rho^\infty \frac{n(r)}{r}\,dr
  = \rho^p\,n(\rho) + p\rho^p\int_\rho^{s_1} \frac{n(r)}{r}\,dr .
\]
Comme $\sum s_\nu^p < \infty$, on sait que $\rho^p n(\rho) \to 0$ et que
$\int_0^\infty n(r)\,r^{-1-p}\,dr < \infty$. Pour le second terme, on coupe
en $\delta$ : sur $[\rho, \delta]$, $\rho^p/r \leq (\rho\delta)^p\,r^{-1-p}$,
donc $\rho^p\int_\rho^\delta n(r)\,dr/r \leq (\rho\delta)^p\int_0^\infty n(r)\,
r^{-1-p}\,dr \to 0$ ; sur $[\delta, s_1]$, $\rho^p\int_\delta^{s_1} n(r)\,dr/r
\leq \rho^p\,n(\delta)\log(s_1/\delta) \to 0$.\footnote{La page 147 écrit
« $p\rho^p\int_\rho^\infty n(r)/r\,dr < p\rho^p\int_0^\infty n(r)/r\,dr \to 0$ »,
ce qui supposerait $\int_0 n(r)\,dr/r$ fini ; or $n(r) \to \infty$ quand $r \to
0$ dès que $A$ est de rang infini, et l'intégrale diverge. La coupure en
$\delta$ répare l'argument sans en changer l'idée. Un bloc de quatre lignes,
barré, tentait une autre voie par $\sum \|A_{i+1} - A_i\|_p < \infty$.}

\emph{Preuve de b).} On a
\[
  \sum_{s_\nu(A_i) \leq \rho} s_\nu(A_i)^p = S_p(A_i) - \sum_{s_\nu(A_i) > \rho} s_\nu(A_i)^p
  \longrightarrow S_p(A) - \sum_{s_\nu > \rho} s_\nu^p = \sum_{s_\nu \leq \rho} s_\nu^p,
\]
pour $\rho$ hors de l'ensemble des $s_\nu$, puisque $S_p(A_i) \to S_p(A)$ et
que les $s_\nu(A_i)$ en nombre fini au-dessus de $\rho$ convergent vers les
$s_\nu$ par convergence en norme ; il suffit alors de prendre $\rho$ avec
$\sum_{s_\nu \leq \rho} s_\nu^p < \varepsilon$, puis $i$ assez grand. Le
théorème en résulte : les $s_\nu(A_i)$ convergent terme à terme vers les
$s_\nu$, et les queues sont uniformément petites dans $\ell^p$.

\pagerange{148}{149}
\section*{L'opérateur $R(u) = \det(1+u)\,(1+u)^{-1}$ (pages 148 et 149)}

Deux pages d'une écriture serrée, à l'encre, sur un opérateur $u$ à trace.
Posons
\[
  P_n(u) = \sum_{k=0}^{n} (-1)^k\,\alpha_{n-k}(u)\,u^k, \qquad
  R(u) = \sum_{k \geq 0} P_k(u) = \frac{\det(1+u)}{1+u},
\]
la dernière égalité lorsque $-1$ n'est pas valeur propre ; et en général
$(1+u)^{-1} = 1 - u\,R(u)/\det(1+u)$. La série des $P_k$ converge en norme
par le corollaire ci-dessous.\footnote{La page 148 écrit $P_n(u) = \sum
\alpha_{n-k}(u)\,u^k$ sans le signe $(-1)^k$, et son lemme calcule $\sum_k
u^{n-k}\alpha_k(u)(-1)^k$ avec le signe. C'est le signe qui est juste :
$\det(1+zu)(1+zu)^{-1} = \sum_j \alpha_j z^j\sum_k (-1)^k z^k u^k$ a pour
coefficient de $z^n$ la somme alternée, et pour $u$ scalaire $\lambda$ en
dimension $1$, $P_1 = \alpha_1 - u = 0$ et $R = 1 = \det(1+u)/(1+u)$, tandis
que sans le signe on aurait $R = 1 + 2\lambda + 2\lambda^2 + \cdots$. En
dimension finie $R(u)$ est la matrice adjointe de $1 + u$.}

\emph{Lemme.} Si $u = \sum_i \lambda_i\,e_i \otimes e_i'$ est diagonalisable,
avec $(e_i)$ une base et $(e_i')$ la base duale, alors $P_n(u) = \sum_i
p_i(u)\,e_i \otimes e_i'$ où
\[
  p_i(u) = \sum_{I \not\ni i,\ |I| = n} \prod_{j \in I} \lambda_j
\]
est la $n$-ième fonction symétrique élémentaire des $\lambda_j$, $j \neq i$.
On le voit en calculant successivement $u^n - \alpha_1 u^{n-1}$, $u^n -
\alpha_1 u^{n-1} + \alpha_2 u^{n-2}$, et par récurrence.

\emph{Théorème.} 1) $|P_n(u)| \leq P_n(|u|)$ et $|R(u)| \leq R(|u|)$ pour $u$
normal, au sens des valeurs propres sur une base commune ; en particulier
$P_n(u) \geq 0$ et $R(u) \geq 0$ si $u \geq 0$. 2) Si $u \geq 0$, $0 \leq
P_n(u) \leq \alpha_n(u)\cdot 1$ et $0 \leq R(u) \leq \det(1+u)$, puisque
$p_i(u)$ est une somme partielle de $\alpha_n(u) = \sum_{|I| = n}\prod_I
\lambda_j$.\footnote{La page 148 justifie 2) en disant que les termes
$\alpha_{n-k}(u)\,u^k$ forment « une suite décroissante d'hermitiens », ce
qui est faux --- pour $\lambda = (10, 1)$ et $n = 2$, $\alpha_2 = 10 <
\lambda_1\alpha_1 = 110$ --- ; c'est la somme alternée qui est comprise
entre $0$ et son premier terme, par le lemme. Le point 1) est énoncé pour
$u$ à trace quelconque, et « devient évident » par le lemme, qui ne
couvre que le cas diagonalisable ; on l'a restreint aux $u$ normaux, où
$|u|$ se diagonalise dans la même base.}

\emph{Corollaire.} Pour $u$ à trace,
\[
  \|R(u)\| \leq \det(1 + |u|) \leq e^{\|u\|_1}, \qquad
  \|P_n(u)\| \leq \alpha_n(|u|) \leq \frac{\|u\|_1^n}{n!} .
\]
\footnote{La page les écrit avec un membre intermédiaire $\|R(|u|)\|$,
$\|P_n(|u|)\|$, qui suppose le point 1) pour $u$ quelconque, non démontré
ici. Les deux inégalités sont vraies sans lui : ce sont les estimations
classiques du premier mineur de Fredholm, $\|\det(1+zu)\,(1+zu)^{-1}\| \leq
\det(1 + |z|\,|u|)$ coefficient par coefficient, qu'on trouve, par les
puissances extérieures, dans Simon, \emph{Trace ideals}, chapitre 5 ; et
$\det(1+|u|) = \prod(1 + s_i) \leq e^{\sum s_i}$.}

\emph{Remarque.} Les valeurs propres de $R(u)$ sont les $p_i = \prod_{j
\neq i}(1 + \lambda_j) = \det(1+u)/(1 + \lambda_i)$, d'où $|p_i| \leq
\det(1 + |u|)$, ce qui redonne une partie du corollaire ; mais comme il y a
des $\lambda_i$ arbitrairement petits, il y a des $p_i$ arbitrairement
voisins de $\det(1+u)$, et
\[
  \|R(u)\| \geq |\det(1+u)| ,
\]
avec égalité pour $u \geq 0$ : l'estimation du corollaire est donc exacte
sur les positifs.\footnote{En dimension infinie, $0$ est valeur spectrale de
$u$ et $\det(1+u)$ est valeur spectrale de $R(u)$, ce qui donne
l'inégalité même si $u$ est de rang fini. C'est le dernier feuillet écrit du
carton.}

\end{document}
